% Authors: Dr. Denys Dutykh (Mathematics Department, Khalifa University of Science and Technology,
% Abu Dhabi, UAE) and Prof. Laurent Vuillon (Univ. Savoie Mont Blanc, CNRS, LAMA, Chambery, France)
\section{What the comparison isolates}\label{sec:outlook}

The optimization problem separates two sources of difficulty in
Markov uniqueness. The occurrence cost is explicit, but a fixed label
may have many primitive two-square representations. Gaussian
factorization describes that entire fiber without identifying its
minimum. The path comparison supplies the missing order for the thirteen
rows treated here. Its content is therefore stronger than a comparison
between selected Markov labels, while its conclusion remains restricted
to the family actually proved.

The parity-order theorem has a scope independent of the row cutoff. It
turns the normalized norm for any fixed defect count into a function
whose extrema on parity boxes occur at prescribed corners, even when
some gaps are infinite. The useful feature is the normalization itself:
bounding the numerator and Fibonacci denominator separately loses the
correlation that makes narrow source intervals accessible.

Extending the result to another row requires more than another finite
list of labels. One must first separate its source interval from every
relevant shorter-cost stratum, including the smaller costs in the final
assembly. The cover procedure has no proved termination theorem for
arbitrary source intervals. A box that remains inside such an interval
is an unresolved comparison, not evidence of a Markov collision.
This is the precise additional work required beyond row thirteen.

The broader question is whether the conclusion of
\cref{thm:thirteen-rows} holds on every row.

\begin{conjecture}[Euclidean minimality of Markov representations]
\label{conj:euclidean-minimality}
Every Markov occurrence weakly minimizes the Euclidean cost over the whole
primitive positive two-square fiber of its label.
\end{conjecture}

By \cref{cor:restricted-injectivity}, the conjecture implies the
Frobenius--Markov uniqueness conjecture. Beyond the thirteen rows the
evidence is computational. By complete factorization of the labels, every
one of the $4145$ occurrences with denominator at most $100$, or with
$P+Q\leqslant160$, is strictly minimal: its pair is the only representation of
minimum cost. \Cref{fig:margins} shows by how much, up to the denominator
$100$. The margin can be a single unit: the label
$m_{4/5}=985=12^2+29^2=16^2+27^2$ has the occurrence pair of cost $7$ and a
second pair of cost $8$. \c{C}anak\c{c}{\i} and Schiffler checked their
uniqueness conjectures by computer for all denominators below $70$
\cite[Remark~4.7]{CanakciSchiffler2020}.

\begin{figure}[tb]
\centering
% Authors: Dr. Denys Dutykh (Mathematics Department, Khalifa University of Science and
% Technology, Abu Dhabi, UAE) and Prof. Laurent Vuillon (Univ. Savoie Mont Blanc, CNRS, LAMA)
% Generated by supplement/code/make_figures.py from exact data; do not edit by hand.
\begin{tikzpicture}[font=\footnotesize, >={Stealth[length=4pt]}]
\fill[sand] (0.122,0.036) rectangle (12.322,0.965);
\draw[slate, line width=0.4pt] (0.244,0.071) circle (1.0pt);
\draw[slate, line width=0.4pt] (0.366,0.071) circle (1.0pt);
\draw[slate, line width=0.4pt] (0.366,0.143) circle (1.0pt);
\draw[slate, line width=0.4pt] (0.488,0.071) circle (1.0pt);
\draw[slate, line width=0.4pt] (0.488,0.214) circle (1.0pt);
\draw[slate, line width=0.4pt] (0.61,0.071) circle (1.0pt);
\draw[slate, line width=0.4pt] (0.61,0.143) circle (1.0pt);
\draw[slate, line width=0.4pt] (0.61,0.214) circle (1.0pt);
\draw[slate, line width=0.4pt] (0.732,0.071) circle (1.0pt);
\draw[slate, line width=0.4pt] (0.732,0.357) circle (1.0pt);
\draw[slate, line width=0.4pt] (0.854,0.214) circle (1.0pt);
\draw[slate, line width=0.4pt] (0.854,0.429) circle (1.0pt);
\draw[slate, line width=0.4pt] (0.976,0.071) circle (1.0pt);
\draw[slate, line width=0.4pt] (0.976,0.214) circle (1.0pt);
\draw[slate, line width=0.4pt] (1.098,0.286) circle (1.0pt);
\draw[slate, line width=0.4pt] (1.098,0.357) circle (1.0pt);
\draw[slate, line width=0.4pt] (1.342,0.071) circle (1.0pt);
\draw[slate, line width=0.4pt] (1.342,0.357) circle (1.0pt);
\draw[slate, line width=0.4pt] (1.342,0.5) circle (1.0pt);
\draw[slate, line width=0.4pt] (1.464,0.357) circle (1.0pt);
\draw[slate, line width=0.4pt] (1.586,0.143) circle (1.0pt);
\draw[slate, line width=0.4pt] (1.586,0.286) circle (1.0pt);
\draw[slate, line width=0.4pt] (1.586,0.5) circle (1.0pt);
\draw[slate, line width=0.4pt] (1.586,0.715) circle (1.0pt);
\draw[slate, line width=0.4pt] (1.708,0.071) circle (1.0pt);
\draw[slate, line width=0.4pt] (1.708,0.357) circle (1.0pt);
\draw[slate, line width=0.4pt] (1.708,0.786) circle (1.0pt);
\draw[slate, line width=0.4pt] (1.83,0.143) circle (1.0pt);
\draw[slate, line width=0.4pt] (1.83,0.572) circle (1.0pt);
\draw[slate, line width=0.4pt] (1.83,0.929) circle (1.0pt);
\draw[slate, line width=0.4pt] (1.83,1.001) circle (1.0pt);
\draw[slate, line width=0.4pt] (2.074,0.286) circle (1.0pt);
\draw[slate, line width=0.4pt] (2.074,0.929) circle (1.0pt);
\draw[slate, line width=0.4pt] (2.318,0.357) circle (1.0pt);
\draw[slate, line width=0.4pt] (2.318,0.5) circle (1.0pt);
\draw[slate, line width=0.4pt] (2.318,0.715) circle (1.0pt);
\draw[slate, line width=0.4pt] (2.318,1.073) circle (1.0pt);
\draw[slate, line width=0.4pt] (2.318,1.144) circle (1.0pt);
\draw[slate, line width=0.4pt] (2.44,0.214) circle (1.0pt);
\draw[slate, line width=0.4pt] (2.44,0.786) circle (1.0pt);
\draw[slate, line width=0.4pt] (2.44,0.929) circle (1.0pt);
\draw[slate, line width=0.4pt] (2.44,1.215) circle (1.0pt);
\draw[slate, line width=0.4pt] (2.562,0.071) circle (1.0pt);
\draw[slate, line width=0.4pt] (2.562,1.144) circle (1.0pt);
\draw[slate, line width=0.4pt] (2.562,1.43) circle (1.0pt);
\draw[slate, line width=0.4pt] (2.684,0.357) circle (1.0pt);
\draw[slate, line width=0.4pt] (2.806,0.071) circle (1.0pt);
\draw[slate, line width=0.4pt] (2.806,0.357) circle (1.0pt);
\draw[slate, line width=0.4pt] (2.806,0.429) circle (1.0pt);
\draw[slate, line width=0.4pt] (2.806,0.5) circle (1.0pt);
\draw[slate, line width=0.4pt] (2.928,0.357) circle (1.0pt);
\draw[slate, line width=0.4pt] (3.05,0.143) circle (1.0pt);
\draw[slate, line width=0.4pt] (3.05,1.287) circle (1.0pt);
\draw[slate, line width=0.4pt] (3.05,1.644) circle (1.0pt);
\draw[slate, line width=0.4pt] (3.294,0.572) circle (1.0pt);
\draw[slate, line width=0.4pt] (3.294,1.859) circle (1.0pt);
\draw[slate, line width=0.4pt] (3.538,0.143) circle (1.0pt);
\draw[slate, line width=0.4pt] (3.538,1.43) circle (1.0pt);
\draw[slate, line width=0.4pt] (3.538,1.859) circle (1.0pt);
\draw[slate, line width=0.4pt] (3.66,0.929) circle (1.0pt);
\draw[slate, line width=0.4pt] (3.66,1.644) circle (1.0pt);
\draw[slate, line width=0.4pt] (3.66,2.073) circle (1.0pt);
\draw[slate, line width=0.4pt] (3.782,0.429) circle (1.0pt);
\draw[slate, line width=0.4pt] (3.904,0.5) circle (1.0pt);
\draw[slate, line width=0.4pt] (4.026,0.143) circle (1.0pt);
\draw[slate, line width=0.4pt] (4.148,1.787) circle (1.0pt);
\draw[slate, line width=0.4pt] (4.148,2.216) circle (1.0pt);
\draw[slate, line width=0.4pt] (4.27,0.572) circle (1.0pt);
\draw[slate, line width=0.4pt] (4.514,0.715) circle (1.0pt);
\draw[slate, line width=0.4pt] (4.636,0.357) circle (1.0pt);
\draw[slate, line width=0.4pt] (4.636,0.5) circle (1.0pt);
\draw[slate, line width=0.4pt] (4.88,0.214) circle (1.0pt);
\draw[slate, line width=0.4pt] (5.002,0.071) circle (1.0pt);
\draw[slate, line width=0.4pt] (5.002,1.001) circle (1.0pt);
\draw[slate, line width=0.4pt] (5.002,2.073) circle (1.0pt);
\draw[slate, line width=0.4pt] (5.002,2.216) circle (1.0pt);
\draw[slate, line width=0.4pt] (5.124,0.357) circle (1.0pt);
\draw[slate, line width=0.4pt] (5.124,1.787) circle (1.0pt);
\draw[slate, line width=0.4pt] (5.246,0.572) circle (1.0pt);
\draw[slate, line width=0.4pt] (5.246,0.858) circle (1.0pt);
\draw[slate, line width=0.4pt] (5.246,1.573) circle (1.0pt);
\draw[slate, line width=0.4pt] (5.49,0.786) circle (1.0pt);
\draw[slate, line width=0.4pt] (5.49,2.216) circle (1.0pt);
\draw[slate, line width=0.4pt] (5.49,2.931) circle (1.0pt);
\draw[slate, line width=0.4pt] (5.49,3.146) circle (1.0pt);
\draw[slate, line width=0.4pt] (5.612,0.643) circle (1.0pt);
\draw[slate, line width=0.4pt] (5.612,2.216) circle (1.0pt);
\draw[slate, line width=0.4pt] (5.734,1.287) circle (1.0pt);
\draw[slate, line width=0.4pt] (5.734,2.145) circle (1.0pt);
\draw[slate, line width=0.4pt] (5.734,2.645) circle (1.0pt);
\draw[slate, line width=0.4pt] (5.978,3.432) circle (1.0pt);
\draw[slate, line width=0.4pt] (6.1,1.215) circle (1.0pt);
\draw[slate, line width=0.4pt] (6.1,3.074) circle (1.0pt);
\draw[slate, line width=0.4pt] (6.222,0.715) circle (1.0pt);
\draw[slate, line width=0.4pt] (6.222,1.43) circle (1.0pt);
\draw[slate, line width=0.4pt] (6.222,3.575) circle (1.0pt);
\draw[slate, line width=0.4pt] (6.466,0.357) circle (1.0pt);
\draw[slate, line width=0.4pt] (6.466,1.073) circle (1.0pt);
\draw[slate, line width=0.4pt] (6.466,1.358) circle (1.0pt);
\draw[slate, line width=0.4pt] (6.588,1.358) circle (1.0pt);
\draw[slate, line width=0.4pt] (6.588,2.502) circle (1.0pt);
\draw[slate, line width=0.4pt] (6.588,3.074) circle (1.0pt);
\draw[slate, line width=0.4pt] (6.71,2.288) circle (1.0pt);
\draw[slate, line width=0.4pt] (6.71,2.431) circle (1.0pt);
\draw[slate, line width=0.4pt] (6.832,0.643) circle (1.0pt);
\draw[slate, line width=0.4pt] (6.832,2.931) circle (1.0pt);
\draw[slate, line width=0.4pt] (6.954,1.001) circle (1.0pt);
\draw[slate, line width=0.4pt] (6.954,1.644) circle (1.0pt);
\draw[slate, line width=0.4pt] (6.954,2.073) circle (1.0pt);
\draw[slate, line width=0.4pt] (7.076,0.929) circle (1.0pt);
\draw[slate, line width=0.4pt] (7.076,1.358) circle (1.0pt);
\draw[slate, line width=0.4pt] (7.076,1.787) circle (1.0pt);
\draw[slate, line width=0.4pt] (7.076,2.931) circle (1.0pt);
\draw[slate, line width=0.4pt] (7.198,2.073) circle (1.0pt);
\draw[slate, line width=0.4pt] (7.32,2.216) circle (1.0pt);
\draw[slate, line width=0.4pt] (7.32,2.645) circle (1.0pt);
\draw[slate, line width=0.4pt] (7.442,0.286) circle (1.0pt);
\draw[slate, line width=0.4pt] (7.442,0.643) circle (1.0pt);
\draw[slate, line width=0.4pt] (7.442,1.287) circle (1.0pt);
\draw[slate, line width=0.4pt] (7.442,1.644) circle (1.0pt);
\draw[slate, line width=0.4pt] (7.442,2.145) circle (1.0pt);
\draw[slate, line width=0.4pt] (7.442,3.861) circle (1.0pt);
\draw[slate, line width=0.4pt] (7.564,0.643) circle (1.0pt);
\draw[slate, line width=0.4pt] (7.564,0.929) circle (1.0pt);
\draw[slate, line width=0.4pt] (7.686,1.644) circle (1.0pt);
\draw[slate, line width=0.4pt] (7.808,0.5) circle (1.0pt);
\draw[slate, line width=0.4pt] (7.808,3.074) circle (1.0pt);
\draw[slate, line width=0.4pt] (7.93,0.071) circle (1.0pt);
\draw[slate, line width=0.4pt] (8.052,1.644) circle (1.0pt);
\draw[slate, line width=0.4pt] (8.052,2.931) circle (1.0pt);
\draw[slate, line width=0.4pt] (8.052,3.789) circle (1.0pt);
\draw[slate, line width=0.4pt] (8.174,1.573) circle (1.0pt);
\draw[slate, line width=0.4pt] (8.174,2.931) circle (1.0pt);
\draw[slate, line width=0.4pt] (8.174,4.504) circle (1.0pt);
\draw[slate, line width=0.4pt] (8.296,0.071) circle (1.0pt);
\draw[slate, line width=0.4pt] (8.54,1.215) circle (1.0pt);
\draw[slate, line width=0.4pt] (8.662,0.214) circle (1.0pt);
\draw[slate, line width=0.4pt] (8.662,0.357) circle (1.0pt);
\draw[slate, line width=0.4pt] (8.662,0.786) circle (1.0pt);
\draw[slate, line width=0.4pt] (8.662,1.43) circle (1.0pt);
\draw[slate, line width=0.4pt] (8.784,3.932) circle (1.0pt);
\draw[slate, line width=0.4pt] (8.906,0.357) circle (1.0pt);
\draw[slate, line width=0.4pt] (8.906,1.144) circle (1.0pt);
\draw[slate, line width=0.4pt] (8.906,4.433) circle (1.0pt);
\draw[slate, line width=0.4pt] (9.028,0.357) circle (1.0pt);
\draw[slate, line width=0.4pt] (9.15,0.143) circle (1.0pt);
\draw[slate, line width=0.4pt] (9.15,1.859) circle (1.0pt);
\draw[slate, line width=0.4pt] (9.15,3.36) circle (1.0pt);
\draw[slate, line width=0.4pt] (9.394,1.073) circle (1.0pt);
\draw[slate, line width=0.4pt] (9.394,1.93) circle (1.0pt);
\draw[slate, line width=0.4pt] (9.394,2.288) circle (1.0pt);
\draw[slate, line width=0.4pt] (9.394,2.431) circle (1.0pt);
\draw[slate, line width=0.4pt] (9.516,1.787) circle (1.0pt);
\draw[slate, line width=0.4pt] (9.638,2.288) circle (1.0pt);
\draw[slate, line width=0.4pt] (9.638,2.502) circle (1.0pt);
\draw[slate, line width=0.4pt] (9.638,4.147) circle (1.0pt);
\draw[slate, line width=0.4pt] (9.638,4.361) circle (1.0pt);
\draw[slate, line width=0.4pt] (9.882,2.002) circle (1.0pt);
\draw[slate, line width=0.4pt] (9.882,3.074) circle (1.0pt);
\draw[slate, line width=0.4pt] (9.882,3.718) circle (1.0pt);
\draw[slate, line width=0.4pt] (10.004,0.5) circle (1.0pt);
\draw[slate, line width=0.4pt] (10.004,2.073) circle (1.0pt);
\draw[slate, line width=0.4pt] (10.126,2.216) circle (1.0pt);
\draw[slate, line width=0.4pt] (10.126,2.502) circle (1.0pt);
\draw[slate, line width=0.4pt] (10.248,1.358) circle (1.0pt);
\draw[slate, line width=0.4pt] (10.248,5.934) circle (1.0pt);
\draw[slate, line width=0.4pt] (10.37,4.719) circle (1.0pt);
\draw[slate, line width=0.4pt] (10.492,0.929) circle (1.0pt);
\draw[slate, line width=0.4pt] (10.492,3.789) circle (1.0pt);
\draw[slate, line width=0.4pt] (10.492,3.932) circle (1.0pt);
\draw[slate, line width=0.4pt] (10.614,5.72) circle (1.0pt);
\draw[slate, line width=0.4pt] (10.736,2.645) circle (1.0pt);
\draw[slate, line width=0.4pt] (10.858,5.076) circle (1.0pt);
\draw[slate, line width=0.4pt] (10.98,5.219) circle (1.0pt);
\draw[slate, line width=0.4pt] (11.102,5.219) circle (1.0pt);
\draw[slate, line width=0.4pt] (11.102,6.435) circle (1.0pt);
\draw[slate, line width=0.4pt] (11.224,0.643) circle (1.0pt);
\draw[slate, line width=0.4pt] (11.224,6.077) circle (1.0pt);
\draw[slate, line width=0.4pt] (11.346,3.789) circle (1.0pt);
\draw[slate, line width=0.4pt] (11.468,0.929) circle (1.0pt);
\draw[slate, line width=0.4pt] (11.468,1.073) circle (1.0pt);
\draw[slate, line width=0.4pt] (11.468,1.787) circle (1.0pt);
\draw[slate, line width=0.4pt] (11.468,1.93) circle (1.0pt);
\draw[slate, line width=0.4pt] (11.59,1.644) circle (1.0pt);
\draw[slate, line width=0.4pt] (11.59,2.431) circle (1.0pt);
\draw[slate, line width=0.4pt] (11.59,6.363) circle (1.0pt);
\draw[slate, line width=0.4pt] (11.712,6.792) circle (1.0pt);
\draw[slate, line width=0.4pt] (11.834,1.787) circle (1.0pt);
\draw[slate, line width=0.4pt] (11.834,6.292) circle (1.0pt);
\draw[slate, line width=0.4pt] (11.956,1.358) circle (1.0pt);
\draw[slate, line width=0.4pt] (11.956,1.644) circle (1.0pt);
\draw[slate, line width=0.4pt] (11.956,6.649) circle (1.0pt);
\draw[slate, line width=0.4pt] (12.078,1.787) circle (1.0pt);
\draw[slate, line width=0.4pt] (12.078,2.073) circle (1.0pt);
\draw[slate, line width=0.4pt] (12.078,5.005) circle (1.0pt);
\fill[navy] (1.22,0.643) circle (0.75pt);
\fill[navy] (1.342,0.214) circle (0.75pt);
\fill[navy] (2.562,1.215) circle (0.75pt);
\fill[navy] (2.806,0.286) circle (0.75pt);
\fill[navy] (2.928,0.929) circle (0.75pt);
\fill[navy] (3.05,0.572) circle (0.75pt);
\fill[navy] (3.172,1.501) circle (0.75pt);
\fill[navy] (3.172,1.644) circle (0.75pt);
\fill[navy] (3.66,1.215) circle (0.75pt);
\fill[navy] (3.782,0.643) circle (0.75pt);
\fill[navy] (3.782,0.786) circle (0.75pt);
\fill[navy] (3.904,0.214) circle (0.75pt);
\fill[navy] (4.026,1.144) circle (0.75pt);
\fill[navy] (4.026,1.43) circle (0.75pt);
\fill[navy] (4.026,2.216) circle (0.75pt);
\fill[navy] (4.148,0.357) circle (0.75pt);
\fill[navy] (4.148,2.073) circle (0.75pt);
\fill[navy] (4.148,2.359) circle (0.75pt);
\fill[navy] (4.27,0.643) circle (0.75pt);
\fill[navy] (4.392,1.358) circle (0.75pt);
\fill[navy] (4.392,2.216) circle (0.75pt);
\fill[navy] (4.514,0.357) circle (0.75pt);
\fill[navy] (4.514,1.001) circle (0.75pt);
\fill[navy] (4.514,1.644) circle (0.75pt);
\fill[navy] (4.514,1.859) circle (0.75pt);
\fill[navy] (4.514,2.502) circle (0.75pt);
\fill[navy] (4.514,2.574) circle (0.75pt);
\fill[navy] (4.636,1.644) circle (0.75pt);
\fill[navy] (4.636,2.502) circle (0.75pt);
\fill[navy] (4.758,1.644) circle (0.75pt);
\fill[navy] (4.758,1.787) circle (0.75pt);
\fill[navy] (4.758,2.002) circle (0.75pt);
\fill[navy] (4.758,2.216) circle (0.75pt);
\fill[navy] (4.88,2.645) circle (0.75pt);
\fill[navy] (5.002,0.214) circle (0.75pt);
\fill[navy] (5.002,0.429) circle (0.75pt);
\fill[navy] (5.002,1.859) circle (0.75pt);
\fill[navy] (5.002,1.93) circle (0.75pt);
\fill[navy] (5.002,2.288) circle (0.75pt);
\fill[navy] (5.124,0.786) circle (0.75pt);
\fill[navy] (5.124,2.216) circle (0.75pt);
\fill[navy] (5.124,2.931) circle (0.75pt);
\fill[navy] (5.246,0.357) circle (0.75pt);
\fill[navy] (5.246,0.715) circle (0.75pt);
\fill[navy] (5.246,1.215) circle (0.75pt);
\fill[navy] (5.246,1.287) circle (0.75pt);
\fill[navy] (5.246,2.288) circle (0.75pt);
\fill[navy] (5.246,2.717) circle (0.75pt);
\fill[navy] (5.368,0.357) circle (0.75pt);
\fill[navy] (5.368,2.502) circle (0.75pt);
\fill[navy] (5.368,2.788) circle (0.75pt);
\fill[navy] (5.49,1.144) circle (0.75pt);
\fill[navy] (5.49,2.717) circle (0.75pt);
\fill[navy] (5.612,2.502) circle (0.75pt);
\fill[navy] (5.612,3.074) circle (0.75pt);
\fill[navy] (5.734,0.214) circle (0.75pt);
\fill[navy] (5.734,0.5) circle (0.75pt);
\fill[navy] (5.734,0.786) circle (0.75pt);
\fill[navy] (5.734,2.717) circle (0.75pt);
\fill[navy] (5.734,2.788) circle (0.75pt);
\fill[navy] (5.734,2.86) circle (0.75pt);
\fill[navy] (5.734,3.217) circle (0.75pt);
\fill[navy] (5.856,1.787) circle (0.75pt);
\fill[navy] (5.978,2.788) circle (0.75pt);
\fill[navy] (5.978,2.931) circle (0.75pt);
\fill[navy] (5.978,3.146) circle (0.75pt);
\fill[navy] (5.978,3.217) circle (0.75pt);
\fill[navy] (5.978,3.289) circle (0.75pt);
\fill[navy] (5.978,3.36) circle (0.75pt);
\fill[navy] (6.1,0.214) circle (0.75pt);
\fill[navy] (6.1,2.073) circle (0.75pt);
\fill[navy] (6.222,1.787) circle (0.75pt);
\fill[navy] (6.222,2.073) circle (0.75pt);
\fill[navy] (6.222,2.931) circle (0.75pt);
\fill[navy] (6.344,0.357) circle (0.75pt);
\fill[navy] (6.344,1.93) circle (0.75pt);
\fill[navy] (6.344,2.645) circle (0.75pt);
\fill[navy] (6.344,3.646) circle (0.75pt);
\fill[navy] (6.466,0.643) circle (0.75pt);
\fill[navy] (6.466,0.858) circle (0.75pt);
\fill[navy] (6.466,0.929) circle (0.75pt);
\fill[navy] (6.466,1.144) circle (0.75pt);
\fill[navy] (6.466,1.787) circle (0.75pt);
\fill[navy] (6.466,1.93) circle (0.75pt);
\fill[navy] (6.466,2.431) circle (0.75pt);
\fill[navy] (6.466,3.003) circle (0.75pt);
\fill[navy] (6.466,3.503) circle (0.75pt);
\fill[navy] (6.588,0.071) circle (0.75pt);
\fill[navy] (6.71,0.572) circle (0.75pt);
\fill[navy] (6.71,0.929) circle (0.75pt);
\fill[navy] (6.71,1.501) circle (0.75pt);
\fill[navy] (6.71,2.216) circle (0.75pt);
\fill[navy] (6.71,3.003) circle (0.75pt);
\fill[navy] (6.71,3.074) circle (0.75pt);
\fill[navy] (6.71,3.503) circle (0.75pt);
\fill[navy] (6.71,3.646) circle (0.75pt);
\fill[navy] (6.832,0.071) circle (0.75pt);
\fill[navy] (6.832,1.787) circle (0.75pt);
\fill[navy] (6.832,3.074) circle (0.75pt);
\fill[navy] (6.832,3.789) circle (0.75pt);
\fill[navy] (6.954,0.786) circle (0.75pt);
\fill[navy] (6.954,2.645) circle (0.75pt);
\fill[navy] (6.954,3.146) circle (0.75pt);
\fill[navy] (6.954,3.718) circle (0.75pt);
\fill[navy] (7.076,0.643) circle (0.75pt);
\fill[navy] (7.076,3.789) circle (0.75pt);
\fill[navy] (7.076,3.932) circle (0.75pt);
\fill[navy] (7.198,0.929) circle (0.75pt);
\fill[navy] (7.198,1.144) circle (0.75pt);
\fill[navy] (7.198,1.358) circle (0.75pt);
\fill[navy] (7.198,1.43) circle (0.75pt);
\fill[navy] (7.198,1.787) circle (0.75pt);
\fill[navy] (7.198,2.288) circle (0.75pt);
\fill[navy] (7.198,2.717) circle (0.75pt);
\fill[navy] (7.198,2.788) circle (0.75pt);
\fill[navy] (7.198,3.146) circle (0.75pt);
\fill[navy] (7.198,3.217) circle (0.75pt);
\fill[navy] (7.198,3.36) circle (0.75pt);
\fill[navy] (7.198,3.432) circle (0.75pt);
\fill[navy] (7.198,3.646) circle (0.75pt);
\fill[navy] (7.198,3.932) circle (0.75pt);
\fill[navy] (7.198,4.004) circle (0.75pt);
\fill[navy] (7.32,3.36) circle (0.75pt);
\fill[navy] (7.32,3.789) circle (0.75pt);
\fill[navy] (7.442,0.357) circle (0.75pt);
\fill[navy] (7.442,1.144) circle (0.75pt);
\fill[navy] (7.442,1.787) circle (0.75pt);
\fill[navy] (7.442,1.93) circle (0.75pt);
\fill[navy] (7.442,2.073) circle (0.75pt);
\fill[navy] (7.442,2.216) circle (0.75pt);
\fill[navy] (7.442,2.359) circle (0.75pt);
\fill[navy] (7.442,2.574) circle (0.75pt);
\fill[navy] (7.442,2.645) circle (0.75pt);
\fill[navy] (7.442,2.86) circle (0.75pt);
\fill[navy] (7.442,3.003) circle (0.75pt);
\fill[navy] (7.442,3.36) circle (0.75pt);
\fill[navy] (7.442,3.432) circle (0.75pt);
\fill[navy] (7.442,3.503) circle (0.75pt);
\fill[navy] (7.442,3.646) circle (0.75pt);
\fill[navy] (7.442,3.789) circle (0.75pt);
\fill[navy] (7.564,0.786) circle (0.75pt);
\fill[navy] (7.564,1.644) circle (0.75pt);
\fill[navy] (7.564,2.359) circle (0.75pt);
\fill[navy] (7.564,2.931) circle (0.75pt);
\fill[navy] (7.564,3.217) circle (0.75pt);
\fill[navy] (7.564,3.503) circle (0.75pt);
\fill[navy] (7.564,3.646) circle (0.75pt);
\fill[navy] (7.686,0.071) circle (0.75pt);
\fill[navy] (7.686,0.715) circle (0.75pt);
\fill[navy] (7.686,1.144) circle (0.75pt);
\fill[navy] (7.686,1.787) circle (0.75pt);
\fill[navy] (7.686,1.859) circle (0.75pt);
\fill[navy] (7.686,2.431) circle (0.75pt);
\fill[navy] (7.686,3.074) circle (0.75pt);
\fill[navy] (7.686,3.718) circle (0.75pt);
\fill[navy] (7.686,4.218) circle (0.75pt);
\fill[navy] (7.686,4.361) circle (0.75pt);
\fill[navy] (7.808,0.929) circle (0.75pt);
\fill[navy] (7.808,1.073) circle (0.75pt);
\fill[navy] (7.808,1.358) circle (0.75pt);
\fill[navy] (7.808,1.787) circle (0.75pt);
\fill[navy] (7.808,2.073) circle (0.75pt);
\fill[navy] (7.808,2.216) circle (0.75pt);
\fill[navy] (7.808,2.645) circle (0.75pt);
\fill[navy] (7.808,2.931) circle (0.75pt);
\fill[navy] (7.808,3.217) circle (0.75pt);
\fill[navy] (7.808,3.36) circle (0.75pt);
\fill[navy] (7.808,3.503) circle (0.75pt);
\fill[navy] (7.808,3.646) circle (0.75pt);
\fill[navy] (7.808,4.504) circle (0.75pt);
\fill[navy] (7.93,0.286) circle (0.75pt);
\fill[navy] (7.93,0.858) circle (0.75pt);
\fill[navy] (7.93,1.144) circle (0.75pt);
\fill[navy] (7.93,1.215) circle (0.75pt);
\fill[navy] (7.93,1.358) circle (0.75pt);
\fill[navy] (7.93,1.573) circle (0.75pt);
\fill[navy] (7.93,1.644) circle (0.75pt);
\fill[navy] (7.93,2.359) circle (0.75pt);
\fill[navy] (7.93,2.574) circle (0.75pt);
\fill[navy] (7.93,2.645) circle (0.75pt);
\fill[navy] (7.93,3.289) circle (0.75pt);
\fill[navy] (7.93,3.36) circle (0.75pt);
\fill[navy] (7.93,3.432) circle (0.75pt);
\fill[navy] (7.93,4.004) circle (0.75pt);
\fill[navy] (7.93,4.075) circle (0.75pt);
\fill[navy] (7.93,4.218) circle (0.75pt);
\fill[navy] (7.93,4.361) circle (0.75pt);
\fill[navy] (7.93,4.433) circle (0.75pt);
\fill[navy] (7.93,4.504) circle (0.75pt);
\fill[navy] (8.052,0.357) circle (0.75pt);
\fill[navy] (8.052,3.503) circle (0.75pt);
\fill[navy] (8.052,4.218) circle (0.75pt);
\fill[navy] (8.052,4.647) circle (0.75pt);
\fill[navy] (8.174,0.429) circle (0.75pt);
\fill[navy] (8.174,1.001) circle (0.75pt);
\fill[navy] (8.174,1.215) circle (0.75pt);
\fill[navy] (8.174,1.287) circle (0.75pt);
\fill[navy] (8.174,2.288) circle (0.75pt);
\fill[navy] (8.174,2.359) circle (0.75pt);
\fill[navy] (8.174,2.788) circle (0.75pt);
\fill[navy] (8.174,3.074) circle (0.75pt);
\fill[navy] (8.174,3.646) circle (0.75pt);
\fill[navy] (8.174,3.932) circle (0.75pt);
\fill[navy] (8.174,4.361) circle (0.75pt);
\fill[navy] (8.174,4.647) circle (0.75pt);
\fill[navy] (8.296,0.5) circle (0.75pt);
\fill[navy] (8.296,1.073) circle (0.75pt);
\fill[navy] (8.296,1.644) circle (0.75pt);
\fill[navy] (8.296,1.787) circle (0.75pt);
\fill[navy] (8.296,1.93) circle (0.75pt);
\fill[navy] (8.296,2.645) circle (0.75pt);
\fill[navy] (8.296,2.788) circle (0.75pt);
\fill[navy] (8.296,3.074) circle (0.75pt);
\fill[navy] (8.296,4.218) circle (0.75pt);
\fill[navy] (8.296,4.647) circle (0.75pt);
\fill[navy] (8.418,0.572) circle (0.75pt);
\fill[navy] (8.418,1.358) circle (0.75pt);
\fill[navy] (8.418,1.43) circle (0.75pt);
\fill[navy] (8.418,1.573) circle (0.75pt);
\fill[navy] (8.418,1.787) circle (0.75pt);
\fill[navy] (8.418,2.073) circle (0.75pt);
\fill[navy] (8.418,2.717) circle (0.75pt);
\fill[navy] (8.418,2.931) circle (0.75pt);
\fill[navy] (8.418,3.503) circle (0.75pt);
\fill[navy] (8.418,4.147) circle (0.75pt);
\fill[navy] (8.418,4.361) circle (0.75pt);
\fill[navy] (8.418,4.647) circle (0.75pt);
\fill[navy] (8.418,4.862) circle (0.75pt);
\fill[navy] (8.54,1.644) circle (0.75pt);
\fill[navy] (8.54,1.93) circle (0.75pt);
\fill[navy] (8.54,2.073) circle (0.75pt);
\fill[navy] (8.54,2.645) circle (0.75pt);
\fill[navy] (8.54,2.931) circle (0.75pt);
\fill[navy] (8.54,3.36) circle (0.75pt);
\fill[navy] (8.54,4.075) circle (0.75pt);
\fill[navy] (8.54,4.218) circle (0.75pt);
\fill[navy] (8.54,4.361) circle (0.75pt);
\fill[navy] (8.54,4.933) circle (0.75pt);
\fill[navy] (8.662,0.143) circle (0.75pt);
\fill[navy] (8.662,0.286) circle (0.75pt);
\fill[navy] (8.662,0.929) circle (0.75pt);
\fill[navy] (8.662,1.501) circle (0.75pt);
\fill[navy] (8.662,2.431) circle (0.75pt);
\fill[navy] (8.662,2.931) circle (0.75pt);
\fill[navy] (8.662,3.718) circle (0.75pt);
\fill[navy] (8.662,3.789) circle (0.75pt);
\fill[navy] (8.662,3.861) circle (0.75pt);
\fill[navy] (8.662,3.932) circle (0.75pt);
\fill[navy] (8.662,4.004) circle (0.75pt);
\fill[navy] (8.662,4.433) circle (0.75pt);
\fill[navy] (8.662,4.79) circle (0.75pt);
\fill[navy] (8.662,4.862) circle (0.75pt);
\fill[navy] (8.662,4.933) circle (0.75pt);
\fill[navy] (8.784,0.929) circle (0.75pt);
\fill[navy] (8.784,1.215) circle (0.75pt);
\fill[navy] (8.784,1.358) circle (0.75pt);
\fill[navy] (8.784,1.644) circle (0.75pt);
\fill[navy] (8.784,2.073) circle (0.75pt);
\fill[navy] (8.784,2.931) circle (0.75pt);
\fill[navy] (8.784,3.503) circle (0.75pt);
\fill[navy] (8.784,4.218) circle (0.75pt);
\fill[navy] (8.784,4.361) circle (0.75pt);
\fill[navy] (8.784,4.647) circle (0.75pt);
\fill[navy] (8.784,4.79) circle (0.75pt);
\fill[navy] (8.906,0.643) circle (0.75pt);
\fill[navy] (8.906,1.215) circle (0.75pt);
\fill[navy] (8.906,1.501) circle (0.75pt);
\fill[navy] (8.906,1.644) circle (0.75pt);
\fill[navy] (8.906,1.716) circle (0.75pt);
\fill[navy] (8.906,1.787) circle (0.75pt);
\fill[navy] (8.906,1.859) circle (0.75pt);
\fill[navy] (8.906,2.931) circle (0.75pt);
\fill[navy] (8.906,3.074) circle (0.75pt);
\fill[navy] (8.906,3.217) circle (0.75pt);
\fill[navy] (8.906,3.289) circle (0.75pt);
\fill[navy] (8.906,3.36) circle (0.75pt);
\fill[navy] (8.906,3.575) circle (0.75pt);
\fill[navy] (8.906,3.646) circle (0.75pt);
\fill[navy] (8.906,3.718) circle (0.75pt);
\fill[navy] (8.906,3.789) circle (0.75pt);
\fill[navy] (8.906,3.861) circle (0.75pt);
\fill[navy] (8.906,4.004) circle (0.75pt);
\fill[navy] (8.906,4.147) circle (0.75pt);
\fill[navy] (8.906,4.29) circle (0.75pt);
\fill[navy] (8.906,4.361) circle (0.75pt);
\fill[navy] (8.906,4.576) circle (0.75pt);
\fill[navy] (8.906,4.647) circle (0.75pt);
\fill[navy] (8.906,4.79) circle (0.75pt);
\fill[navy] (8.906,4.933) circle (0.75pt);
\fill[navy] (8.906,5.148) circle (0.75pt);
\fill[navy] (9.028,0.643) circle (0.75pt);
\fill[navy] (9.028,1.358) circle (0.75pt);
\fill[navy] (9.028,2.359) circle (0.75pt);
\fill[navy] (9.028,2.502) circle (0.75pt);
\fill[navy] (9.028,2.931) circle (0.75pt);
\fill[navy] (9.028,3.503) circle (0.75pt);
\fill[navy] (9.028,4.361) circle (0.75pt);
\fill[navy] (9.028,4.933) circle (0.75pt);
\fill[navy] (9.15,0.572) circle (0.75pt);
\fill[navy] (9.15,2.073) circle (0.75pt);
\fill[navy] (9.15,2.216) circle (0.75pt);
\fill[navy] (9.15,2.288) circle (0.75pt);
\fill[navy] (9.15,2.431) circle (0.75pt);
\fill[navy] (9.15,3.074) circle (0.75pt);
\fill[navy] (9.15,3.289) circle (0.75pt);
\fill[navy] (9.15,3.718) circle (0.75pt);
\fill[navy] (9.15,4.218) circle (0.75pt);
\fill[navy] (9.15,4.433) circle (0.75pt);
\fill[navy] (9.15,4.79) circle (0.75pt);
\fill[navy] (9.15,5.291) circle (0.75pt);
\fill[navy] (9.272,0.929) circle (0.75pt);
\fill[navy] (9.272,1.215) circle (0.75pt);
\fill[navy] (9.272,1.644) circle (0.75pt);
\fill[navy] (9.272,1.787) circle (0.75pt);
\fill[navy] (9.272,2.359) circle (0.75pt);
\fill[navy] (9.272,2.502) circle (0.75pt);
\fill[navy] (9.272,2.645) circle (0.75pt);
\fill[navy] (9.272,3.503) circle (0.75pt);
\fill[navy] (9.272,3.932) circle (0.75pt);
\fill[navy] (9.272,4.218) circle (0.75pt);
\fill[navy] (9.272,4.361) circle (0.75pt);
\fill[navy] (9.272,4.647) circle (0.75pt);
\fill[navy] (9.272,4.79) circle (0.75pt);
\fill[navy] (9.272,5.076) circle (0.75pt);
\fill[navy] (9.272,5.219) circle (0.75pt);
\fill[navy] (9.272,5.362) circle (0.75pt);
\fill[navy] (9.394,0.429) circle (0.75pt);
\fill[navy] (9.394,1.43) circle (0.75pt);
\fill[navy] (9.394,2.073) circle (0.75pt);
\fill[navy] (9.394,2.216) circle (0.75pt);
\fill[navy] (9.394,2.717) circle (0.75pt);
\fill[navy] (9.394,2.788) circle (0.75pt);
\fill[navy] (9.394,3.432) circle (0.75pt);
\fill[navy] (9.394,3.718) circle (0.75pt);
\fill[navy] (9.394,3.789) circle (0.75pt);
\fill[navy] (9.394,4.433) circle (0.75pt);
\fill[navy] (9.394,4.647) circle (0.75pt);
\fill[navy] (9.394,4.862) circle (0.75pt);
\fill[navy] (9.394,5.148) circle (0.75pt);
\fill[navy] (9.394,5.219) circle (0.75pt);
\fill[navy] (9.394,5.362) circle (0.75pt);
\fill[navy] (9.516,0.357) circle (0.75pt);
\fill[navy] (9.516,0.786) circle (0.75pt);
\fill[navy] (9.516,1.644) circle (0.75pt);
\fill[navy] (9.516,3.074) circle (0.75pt);
\fill[navy] (9.516,3.789) circle (0.75pt);
\fill[navy] (9.516,3.932) circle (0.75pt);
\fill[navy] (9.516,4.361) circle (0.75pt);
\fill[navy] (9.516,4.79) circle (0.75pt);
\fill[navy] (9.516,5.076) circle (0.75pt);
\fill[navy] (9.638,0.643) circle (0.75pt);
\fill[navy] (9.638,1.001) circle (0.75pt);
\fill[navy] (9.638,1.287) circle (0.75pt);
\fill[navy] (9.638,1.358) circle (0.75pt);
\fill[navy] (9.638,1.573) circle (0.75pt);
\fill[navy] (9.638,1.859) circle (0.75pt);
\fill[navy] (9.638,2.073) circle (0.75pt);
\fill[navy] (9.638,2.216) circle (0.75pt);
\fill[navy] (9.638,2.431) circle (0.75pt);
\fill[navy] (9.638,2.717) circle (0.75pt);
\fill[navy] (9.638,2.931) circle (0.75pt);
\fill[navy] (9.638,3.074) circle (0.75pt);
\fill[navy] (9.638,3.432) circle (0.75pt);
\fill[navy] (9.638,3.575) circle (0.75pt);
\fill[navy] (9.638,3.646) circle (0.75pt);
\fill[navy] (9.638,3.861) circle (0.75pt);
\fill[navy] (9.638,4.218) circle (0.75pt);
\fill[navy] (9.638,4.29) circle (0.75pt);
\fill[navy] (9.638,4.433) circle (0.75pt);
\fill[navy] (9.638,4.504) circle (0.75pt);
\fill[navy] (9.638,4.576) circle (0.75pt);
\fill[navy] (9.638,4.719) circle (0.75pt);
\fill[navy] (9.638,4.862) circle (0.75pt);
\fill[navy] (9.638,5.005) circle (0.75pt);
\fill[navy] (9.638,5.076) circle (0.75pt);
\fill[navy] (9.638,5.219) circle (0.75pt);
\fill[navy] (9.638,5.434) circle (0.75pt);
\fill[navy] (9.638,5.505) circle (0.75pt);
\fill[navy] (9.638,5.577) circle (0.75pt);
\fill[navy] (9.76,0.929) circle (0.75pt);
\fill[navy] (9.76,1.358) circle (0.75pt);
\fill[navy] (9.76,1.501) circle (0.75pt);
\fill[navy] (9.76,1.644) circle (0.75pt);
\fill[navy] (9.76,2.359) circle (0.75pt);
\fill[navy] (9.76,2.645) circle (0.75pt);
\fill[navy] (9.76,3.36) circle (0.75pt);
\fill[navy] (9.76,3.646) circle (0.75pt);
\fill[navy] (9.76,3.789) circle (0.75pt);
\fill[navy] (9.76,4.218) circle (0.75pt);
\fill[navy] (9.76,4.504) circle (0.75pt);
\fill[navy] (9.76,4.79) circle (0.75pt);
\fill[navy] (9.76,5.076) circle (0.75pt);
\fill[navy] (9.76,5.219) circle (0.75pt);
\fill[navy] (9.76,5.505) circle (0.75pt);
\fill[navy] (9.882,0.071) circle (0.75pt);
\fill[navy] (9.882,0.143) circle (0.75pt);
\fill[navy] (9.882,0.929) circle (0.75pt);
\fill[navy] (9.882,1.144) circle (0.75pt);
\fill[navy] (9.882,1.215) circle (0.75pt);
\fill[navy] (9.882,1.358) circle (0.75pt);
\fill[navy] (9.882,1.43) circle (0.75pt);
\fill[navy] (9.882,1.644) circle (0.75pt);
\fill[navy] (9.882,1.859) circle (0.75pt);
\fill[navy] (9.882,2.216) circle (0.75pt);
\fill[navy] (9.882,2.288) circle (0.75pt);
\fill[navy] (9.882,2.502) circle (0.75pt);
\fill[navy] (9.882,2.645) circle (0.75pt);
\fill[navy] (9.882,2.717) circle (0.75pt);
\fill[navy] (9.882,2.86) circle (0.75pt);
\fill[navy] (9.882,3.146) circle (0.75pt);
\fill[navy] (9.882,3.289) circle (0.75pt);
\fill[navy] (9.882,4.004) circle (0.75pt);
\fill[navy] (9.882,4.433) circle (0.75pt);
\fill[navy] (9.882,4.576) circle (0.75pt);
\fill[navy] (9.882,4.647) circle (0.75pt);
\fill[navy] (9.882,5.434) circle (0.75pt);
\fill[navy] (9.882,5.505) circle (0.75pt);
\fill[navy] (9.882,5.648) circle (0.75pt);
\fill[navy] (10.004,2.359) circle (0.75pt);
\fill[navy] (10.004,3.074) circle (0.75pt);
\fill[navy] (10.004,3.36) circle (0.75pt);
\fill[navy] (10.004,3.932) circle (0.75pt);
\fill[navy] (10.004,4.361) circle (0.75pt);
\fill[navy] (10.004,4.504) circle (0.75pt);
\fill[navy] (10.004,4.647) circle (0.75pt);
\fill[navy] (10.004,4.79) circle (0.75pt);
\fill[navy] (10.004,5.076) circle (0.75pt);
\fill[navy] (10.004,5.505) circle (0.75pt);
\fill[navy] (10.004,5.648) circle (0.75pt);
\fill[navy] (10.004,5.791) circle (0.75pt);
\fill[navy] (10.126,0.643) circle (0.75pt);
\fill[navy] (10.126,1.144) circle (0.75pt);
\fill[navy] (10.126,1.287) circle (0.75pt);
\fill[navy] (10.126,1.644) circle (0.75pt);
\fill[navy] (10.126,1.716) circle (0.75pt);
\fill[navy] (10.126,2.073) circle (0.75pt);
\fill[navy] (10.126,2.288) circle (0.75pt);
\fill[navy] (10.126,2.359) circle (0.75pt);
\fill[navy] (10.126,3.074) circle (0.75pt);
\fill[navy] (10.126,3.217) circle (0.75pt);
\fill[navy] (10.126,3.432) circle (0.75pt);
\fill[navy] (10.126,3.503) circle (0.75pt);
\fill[navy] (10.126,3.575) circle (0.75pt);
\fill[navy] (10.126,3.718) circle (0.75pt);
\fill[navy] (10.126,4.004) circle (0.75pt);
\fill[navy] (10.126,4.147) circle (0.75pt);
\fill[navy] (10.126,4.29) circle (0.75pt);
\fill[navy] (10.126,4.433) circle (0.75pt);
\fill[navy] (10.126,4.504) circle (0.75pt);
\fill[navy] (10.126,4.647) circle (0.75pt);
\fill[navy] (10.126,4.719) circle (0.75pt);
\fill[navy] (10.126,4.79) circle (0.75pt);
\fill[navy] (10.126,4.862) circle (0.75pt);
\fill[navy] (10.126,4.933) circle (0.75pt);
\fill[navy] (10.126,5.005) circle (0.75pt);
\fill[navy] (10.126,5.076) circle (0.75pt);
\fill[navy] (10.126,5.219) circle (0.75pt);
\fill[navy] (10.126,5.434) circle (0.75pt);
\fill[navy] (10.126,5.577) circle (0.75pt);
\fill[navy] (10.126,5.648) circle (0.75pt);
\fill[navy] (10.126,5.72) circle (0.75pt);
\fill[navy] (10.248,0.357) circle (0.75pt);
\fill[navy] (10.248,2.216) circle (0.75pt);
\fill[navy] (10.248,2.931) circle (0.75pt);
\fill[navy] (10.248,3.932) circle (0.75pt);
\fill[navy] (10.248,4.361) circle (0.75pt);
\fill[navy] (10.248,4.647) circle (0.75pt);
\fill[navy] (10.248,4.79) circle (0.75pt);
\fill[navy] (10.248,5.076) circle (0.75pt);
\fill[navy] (10.248,5.219) circle (0.75pt);
\fill[navy] (10.248,5.648) circle (0.75pt);
\fill[navy] (10.37,0.143) circle (0.75pt);
\fill[navy] (10.37,0.286) circle (0.75pt);
\fill[navy] (10.37,0.572) circle (0.75pt);
\fill[navy] (10.37,0.643) circle (0.75pt);
\fill[navy] (10.37,0.858) circle (0.75pt);
\fill[navy] (10.37,1.287) circle (0.75pt);
\fill[navy] (10.37,1.358) circle (0.75pt);
\fill[navy] (10.37,1.501) circle (0.75pt);
\fill[navy] (10.37,1.573) circle (0.75pt);
\fill[navy] (10.37,1.644) circle (0.75pt);
\fill[navy] (10.37,1.716) circle (0.75pt);
\fill[navy] (10.37,2.073) circle (0.75pt);
\fill[navy] (10.37,2.288) circle (0.75pt);
\fill[navy] (10.37,2.359) circle (0.75pt);
\fill[navy] (10.37,2.574) circle (0.75pt);
\fill[navy] (10.37,3.003) circle (0.75pt);
\fill[navy] (10.37,3.289) circle (0.75pt);
\fill[navy] (10.37,3.718) circle (0.75pt);
\fill[navy] (10.37,3.789) circle (0.75pt);
\fill[navy] (10.37,3.861) circle (0.75pt);
\fill[navy] (10.37,4.075) circle (0.75pt);
\fill[navy] (10.37,4.147) circle (0.75pt);
\fill[navy] (10.37,4.218) circle (0.75pt);
\fill[navy] (10.37,4.433) circle (0.75pt);
\fill[navy] (10.37,4.576) circle (0.75pt);
\fill[navy] (10.37,4.79) circle (0.75pt);
\fill[navy] (10.37,5.076) circle (0.75pt);
\fill[navy] (10.37,5.148) circle (0.75pt);
\fill[navy] (10.37,5.291) circle (0.75pt);
\fill[navy] (10.37,5.434) circle (0.75pt);
\fill[navy] (10.37,5.505) circle (0.75pt);
\fill[navy] (10.37,5.577) circle (0.75pt);
\fill[navy] (10.37,5.648) circle (0.75pt);
\fill[navy] (10.37,5.791) circle (0.75pt);
\fill[navy] (10.37,5.863) circle (0.75pt);
\fill[navy] (10.37,5.934) circle (0.75pt);
\fill[navy] (10.37,6.006) circle (0.75pt);
\fill[navy] (10.492,0.214) circle (0.75pt);
\fill[navy] (10.492,1.501) circle (0.75pt);
\fill[navy] (10.492,1.644) circle (0.75pt);
\fill[navy] (10.492,2.073) circle (0.75pt);
\fill[navy] (10.492,2.359) circle (0.75pt);
\fill[navy] (10.492,2.502) circle (0.75pt);
\fill[navy] (10.492,2.788) circle (0.75pt);
\fill[navy] (10.492,2.931) circle (0.75pt);
\fill[navy] (10.492,4.075) circle (0.75pt);
\fill[navy] (10.492,4.218) circle (0.75pt);
\fill[navy] (10.492,4.361) circle (0.75pt);
\fill[navy] (10.492,4.504) circle (0.75pt);
\fill[navy] (10.492,4.647) circle (0.75pt);
\fill[navy] (10.492,5.076) circle (0.75pt);
\fill[navy] (10.492,5.219) circle (0.75pt);
\fill[navy] (10.492,5.362) circle (0.75pt);
\fill[navy] (10.492,5.648) circle (0.75pt);
\fill[navy] (10.492,5.791) circle (0.75pt);
\fill[navy] (10.492,5.934) circle (0.75pt);
\fill[navy] (10.614,0.357) circle (0.75pt);
\fill[navy] (10.614,0.5) circle (0.75pt);
\fill[navy] (10.614,0.929) circle (0.75pt);
\fill[navy] (10.614,1.001) circle (0.75pt);
\fill[navy] (10.614,1.358) circle (0.75pt);
\fill[navy] (10.614,1.43) circle (0.75pt);
\fill[navy] (10.614,1.787) circle (0.75pt);
\fill[navy] (10.614,2.002) circle (0.75pt);
\fill[navy] (10.614,2.216) circle (0.75pt);
\fill[navy] (10.614,2.288) circle (0.75pt);
\fill[navy] (10.614,2.431) circle (0.75pt);
\fill[navy] (10.614,2.502) circle (0.75pt);
\fill[navy] (10.614,2.645) circle (0.75pt);
\fill[navy] (10.614,2.717) circle (0.75pt);
\fill[navy] (10.614,2.931) circle (0.75pt);
\fill[navy] (10.614,3.146) circle (0.75pt);
\fill[navy] (10.614,3.36) circle (0.75pt);
\fill[navy] (10.614,3.575) circle (0.75pt);
\fill[navy] (10.614,4.361) circle (0.75pt);
\fill[navy] (10.614,4.433) circle (0.75pt);
\fill[navy] (10.614,4.576) circle (0.75pt);
\fill[navy] (10.614,4.647) circle (0.75pt);
\fill[navy] (10.614,4.862) circle (0.75pt);
\fill[navy] (10.614,5.005) circle (0.75pt);
\fill[navy] (10.614,5.076) circle (0.75pt);
\fill[navy] (10.614,5.219) circle (0.75pt);
\fill[navy] (10.614,5.434) circle (0.75pt);
\fill[navy] (10.614,5.648) circle (0.75pt);
\fill[navy] (10.614,5.863) circle (0.75pt);
\fill[navy] (10.614,5.934) circle (0.75pt);
\fill[navy] (10.614,6.149) circle (0.75pt);
\fill[navy] (10.736,1.215) circle (0.75pt);
\fill[navy] (10.736,1.644) circle (0.75pt);
\fill[navy] (10.736,1.787) circle (0.75pt);
\fill[navy] (10.736,2.502) circle (0.75pt);
\fill[navy] (10.736,3.36) circle (0.75pt);
\fill[navy] (10.736,3.503) circle (0.75pt);
\fill[navy] (10.736,4.218) circle (0.75pt);
\fill[navy] (10.736,4.647) circle (0.75pt);
\fill[navy] (10.736,4.79) circle (0.75pt);
\fill[navy] (10.736,4.933) circle (0.75pt);
\fill[navy] (10.736,5.076) circle (0.75pt);
\fill[navy] (10.736,5.219) circle (0.75pt);
\fill[navy] (10.736,5.362) circle (0.75pt);
\fill[navy] (10.736,5.648) circle (0.75pt);
\fill[navy] (10.736,5.791) circle (0.75pt);
\fill[navy] (10.736,6.077) circle (0.75pt);
\fill[navy] (10.858,0.071) circle (0.75pt);
\fill[navy] (10.858,0.143) circle (0.75pt);
\fill[navy] (10.858,0.214) circle (0.75pt);
\fill[navy] (10.858,0.286) circle (0.75pt);
\fill[navy] (10.858,0.572) circle (0.75pt);
\fill[navy] (10.858,0.643) circle (0.75pt);
\fill[navy] (10.858,1.073) circle (0.75pt);
\fill[navy] (10.858,1.215) circle (0.75pt);
\fill[navy] (10.858,1.358) circle (0.75pt);
\fill[navy] (10.858,1.43) circle (0.75pt);
\fill[navy] (10.858,1.501) circle (0.75pt);
\fill[navy] (10.858,1.644) circle (0.75pt);
\fill[navy] (10.858,1.93) circle (0.75pt);
\fill[navy] (10.858,2.002) circle (0.75pt);
\fill[navy] (10.858,2.073) circle (0.75pt);
\fill[navy] (10.858,2.216) circle (0.75pt);
\fill[navy] (10.858,2.288) circle (0.75pt);
\fill[navy] (10.858,2.359) circle (0.75pt);
\fill[navy] (10.858,2.431) circle (0.75pt);
\fill[navy] (10.858,2.788) circle (0.75pt);
\fill[navy] (10.858,2.86) circle (0.75pt);
\fill[navy] (10.858,2.931) circle (0.75pt);
\fill[navy] (10.858,3.074) circle (0.75pt);
\fill[navy] (10.858,3.146) circle (0.75pt);
\fill[navy] (10.858,3.432) circle (0.75pt);
\fill[navy] (10.858,3.575) circle (0.75pt);
\fill[navy] (10.858,3.646) circle (0.75pt);
\fill[navy] (10.858,3.718) circle (0.75pt);
\fill[navy] (10.858,3.789) circle (0.75pt);
\fill[navy] (10.858,3.932) circle (0.75pt);
\fill[navy] (10.858,4.004) circle (0.75pt);
\fill[navy] (10.858,4.075) circle (0.75pt);
\fill[navy] (10.858,4.147) circle (0.75pt);
\fill[navy] (10.858,4.218) circle (0.75pt);
\fill[navy] (10.858,4.361) circle (0.75pt);
\fill[navy] (10.858,4.647) circle (0.75pt);
\fill[navy] (10.858,4.719) circle (0.75pt);
\fill[navy] (10.858,4.862) circle (0.75pt);
\fill[navy] (10.858,5.005) circle (0.75pt);
\fill[navy] (10.858,5.148) circle (0.75pt);
\fill[navy] (10.858,5.219) circle (0.75pt);
\fill[navy] (10.858,5.291) circle (0.75pt);
\fill[navy] (10.858,5.434) circle (0.75pt);
\fill[navy] (10.858,5.505) circle (0.75pt);
\fill[navy] (10.858,5.577) circle (0.75pt);
\fill[navy] (10.858,5.648) circle (0.75pt);
\fill[navy] (10.858,5.72) circle (0.75pt);
\fill[navy] (10.858,5.863) circle (0.75pt);
\fill[navy] (10.858,5.934) circle (0.75pt);
\fill[navy] (10.858,6.006) circle (0.75pt);
\fill[navy] (10.858,6.077) circle (0.75pt);
\fill[navy] (10.858,6.149) circle (0.75pt);
\fill[navy] (10.858,6.22) circle (0.75pt);
\fill[navy] (10.98,0.071) circle (0.75pt);
\fill[navy] (10.98,0.786) circle (0.75pt);
\fill[navy] (10.98,1.644) circle (0.75pt);
\fill[navy] (10.98,2.645) circle (0.75pt);
\fill[navy] (10.98,2.931) circle (0.75pt);
\fill[navy] (10.98,3.36) circle (0.75pt);
\fill[navy] (10.98,4.218) circle (0.75pt);
\fill[navy] (10.98,4.79) circle (0.75pt);
\fill[navy] (10.98,5.076) circle (0.75pt);
\fill[navy] (10.98,5.648) circle (0.75pt);
\fill[navy] (10.98,5.934) circle (0.75pt);
\fill[navy] (10.98,6.363) circle (0.75pt);
\fill[navy] (11.102,0.286) circle (0.75pt);
\fill[navy] (11.102,0.643) circle (0.75pt);
\fill[navy] (11.102,0.715) circle (0.75pt);
\fill[navy] (11.102,0.786) circle (0.75pt);
\fill[navy] (11.102,0.858) circle (0.75pt);
\fill[navy] (11.102,1.073) circle (0.75pt);
\fill[navy] (11.102,1.215) circle (0.75pt);
\fill[navy] (11.102,1.573) circle (0.75pt);
\fill[navy] (11.102,1.644) circle (0.75pt);
\fill[navy] (11.102,1.716) circle (0.75pt);
\fill[navy] (11.102,1.787) circle (0.75pt);
\fill[navy] (11.102,2.073) circle (0.75pt);
\fill[navy] (11.102,2.145) circle (0.75pt);
\fill[navy] (11.102,2.288) circle (0.75pt);
\fill[navy] (11.102,2.431) circle (0.75pt);
\fill[navy] (11.102,3.217) circle (0.75pt);
\fill[navy] (11.102,3.289) circle (0.75pt);
\fill[navy] (11.102,3.36) circle (0.75pt);
\fill[navy] (11.102,3.432) circle (0.75pt);
\fill[navy] (11.102,3.646) circle (0.75pt);
\fill[navy] (11.102,3.789) circle (0.75pt);
\fill[navy] (11.102,3.861) circle (0.75pt);
\fill[navy] (11.102,3.932) circle (0.75pt);
\fill[navy] (11.102,4.075) circle (0.75pt);
\fill[navy] (11.102,4.433) circle (0.75pt);
\fill[navy] (11.102,4.576) circle (0.75pt);
\fill[navy] (11.102,4.719) circle (0.75pt);
\fill[navy] (11.102,4.79) circle (0.75pt);
\fill[navy] (11.102,4.862) circle (0.75pt);
\fill[navy] (11.102,4.933) circle (0.75pt);
\fill[navy] (11.102,5.148) circle (0.75pt);
\fill[navy] (11.102,5.362) circle (0.75pt);
\fill[navy] (11.102,5.72) circle (0.75pt);
\fill[navy] (11.102,5.791) circle (0.75pt);
\fill[navy] (11.102,5.863) circle (0.75pt);
\fill[navy] (11.102,5.934) circle (0.75pt);
\fill[navy] (11.102,6.077) circle (0.75pt);
\fill[navy] (11.102,6.149) circle (0.75pt);
\fill[navy] (11.102,6.22) circle (0.75pt);
\fill[navy] (11.102,6.363) circle (0.75pt);
\fill[navy] (11.224,0.214) circle (0.75pt);
\fill[navy] (11.224,0.786) circle (0.75pt);
\fill[navy] (11.224,1.073) circle (0.75pt);
\fill[navy] (11.224,1.501) circle (0.75pt);
\fill[navy] (11.224,2.073) circle (0.75pt);
\fill[navy] (11.224,2.359) circle (0.75pt);
\fill[navy] (11.224,2.788) circle (0.75pt);
\fill[navy] (11.224,3.36) circle (0.75pt);
\fill[navy] (11.224,3.503) circle (0.75pt);
\fill[navy] (11.224,3.789) circle (0.75pt);
\fill[navy] (11.224,4.218) circle (0.75pt);
\fill[navy] (11.224,4.361) circle (0.75pt);
\fill[navy] (11.224,4.504) circle (0.75pt);
\fill[navy] (11.224,4.79) circle (0.75pt);
\fill[navy] (11.224,5.076) circle (0.75pt);
\fill[navy] (11.224,5.219) circle (0.75pt);
\fill[navy] (11.224,5.362) circle (0.75pt);
\fill[navy] (11.224,5.505) circle (0.75pt);
\fill[navy] (11.224,5.648) circle (0.75pt);
\fill[navy] (11.224,5.791) circle (0.75pt);
\fill[navy] (11.224,5.934) circle (0.75pt);
\fill[navy] (11.224,6.363) circle (0.75pt);
\fill[navy] (11.346,0.357) circle (0.75pt);
\fill[navy] (11.346,0.715) circle (0.75pt);
\fill[navy] (11.346,0.786) circle (0.75pt);
\fill[navy] (11.346,1.001) circle (0.75pt);
\fill[navy] (11.346,1.144) circle (0.75pt);
\fill[navy] (11.346,1.43) circle (0.75pt);
\fill[navy] (11.346,1.573) circle (0.75pt);
\fill[navy] (11.346,1.787) circle (0.75pt);
\fill[navy] (11.346,1.859) circle (0.75pt);
\fill[navy] (11.346,2.073) circle (0.75pt);
\fill[navy] (11.346,2.288) circle (0.75pt);
\fill[navy] (11.346,2.431) circle (0.75pt);
\fill[navy] (11.346,2.645) circle (0.75pt);
\fill[navy] (11.346,2.717) circle (0.75pt);
\fill[navy] (11.346,2.86) circle (0.75pt);
\fill[navy] (11.346,3.074) circle (0.75pt);
\fill[navy] (11.346,3.289) circle (0.75pt);
\fill[navy] (11.346,3.36) circle (0.75pt);
\fill[navy] (11.346,3.503) circle (0.75pt);
\fill[navy] (11.346,3.575) circle (0.75pt);
\fill[navy] (11.346,3.718) circle (0.75pt);
\fill[navy] (11.346,4.147) circle (0.75pt);
\fill[navy] (11.346,4.218) circle (0.75pt);
\fill[navy] (11.346,4.647) circle (0.75pt);
\fill[navy] (11.346,5.005) circle (0.75pt);
\fill[navy] (11.346,5.291) circle (0.75pt);
\fill[navy] (11.346,5.434) circle (0.75pt);
\fill[navy] (11.346,5.505) circle (0.75pt);
\fill[navy] (11.346,5.863) circle (0.75pt);
\fill[navy] (11.346,5.934) circle (0.75pt);
\fill[navy] (11.346,6.149) circle (0.75pt);
\fill[navy] (11.346,6.292) circle (0.75pt);
\fill[navy] (11.346,6.363) circle (0.75pt);
\fill[navy] (11.468,0.786) circle (0.75pt);
\fill[navy] (11.468,1.215) circle (0.75pt);
\fill[navy] (11.468,1.358) circle (0.75pt);
\fill[navy] (11.468,2.359) circle (0.75pt);
\fill[navy] (11.468,2.502) circle (0.75pt);
\fill[navy] (11.468,2.645) circle (0.75pt);
\fill[navy] (11.468,2.788) circle (0.75pt);
\fill[navy] (11.468,3.074) circle (0.75pt);
\fill[navy] (11.468,3.217) circle (0.75pt);
\fill[navy] (11.468,3.646) circle (0.75pt);
\fill[navy] (11.468,3.932) circle (0.75pt);
\fill[navy] (11.468,4.075) circle (0.75pt);
\fill[navy] (11.468,4.218) circle (0.75pt);
\fill[navy] (11.468,4.361) circle (0.75pt);
\fill[navy] (11.468,4.933) circle (0.75pt);
\fill[navy] (11.468,5.076) circle (0.75pt);
\fill[navy] (11.468,5.219) circle (0.75pt);
\fill[navy] (11.468,5.362) circle (0.75pt);
\fill[navy] (11.468,5.648) circle (0.75pt);
\fill[navy] (11.468,5.934) circle (0.75pt);
\fill[navy] (11.468,6.077) circle (0.75pt);
\fill[navy] (11.468,6.22) circle (0.75pt);
\fill[navy] (11.468,6.363) circle (0.75pt);
\fill[navy] (11.468,6.506) circle (0.75pt);
\fill[navy] (11.59,0.214) circle (0.75pt);
\fill[navy] (11.59,0.429) circle (0.75pt);
\fill[navy] (11.59,0.5) circle (0.75pt);
\fill[navy] (11.59,0.929) circle (0.75pt);
\fill[navy] (11.59,1.001) circle (0.75pt);
\fill[navy] (11.59,1.501) circle (0.75pt);
\fill[navy] (11.59,1.573) circle (0.75pt);
\fill[navy] (11.59,1.859) circle (0.75pt);
\fill[navy] (11.59,1.93) circle (0.75pt);
\fill[navy] (11.59,2.216) circle (0.75pt);
\fill[navy] (11.59,2.288) circle (0.75pt);
\fill[navy] (11.59,2.359) circle (0.75pt);
\fill[navy] (11.59,2.788) circle (0.75pt);
\fill[navy] (11.59,2.931) circle (0.75pt);
\fill[navy] (11.59,3.074) circle (0.75pt);
\fill[navy] (11.59,3.432) circle (0.75pt);
\fill[navy] (11.59,3.503) circle (0.75pt);
\fill[navy] (11.59,3.646) circle (0.75pt);
\fill[navy] (11.59,4.004) circle (0.75pt);
\fill[navy] (11.59,4.504) circle (0.75pt);
\fill[navy] (11.59,4.576) circle (0.75pt);
\fill[navy] (11.59,4.79) circle (0.75pt);
\fill[navy] (11.59,5.148) circle (0.75pt);
\fill[navy] (11.59,5.219) circle (0.75pt);
\fill[navy] (11.59,5.291) circle (0.75pt);
\fill[navy] (11.59,5.505) circle (0.75pt);
\fill[navy] (11.59,5.577) circle (0.75pt);
\fill[navy] (11.59,5.648) circle (0.75pt);
\fill[navy] (11.59,5.863) circle (0.75pt);
\fill[navy] (11.59,6.006) circle (0.75pt);
\fill[navy] (11.59,6.149) circle (0.75pt);
\fill[navy] (11.59,6.22) circle (0.75pt);
\fill[navy] (11.59,6.292) circle (0.75pt);
\fill[navy] (11.59,6.506) circle (0.75pt);
\fill[navy] (11.59,6.578) circle (0.75pt);
\fill[navy] (11.59,6.649) circle (0.75pt);
\fill[navy] (11.712,0.786) circle (0.75pt);
\fill[navy] (11.712,1.215) circle (0.75pt);
\fill[navy] (11.712,2.073) circle (0.75pt);
\fill[navy] (11.712,3.503) circle (0.75pt);
\fill[navy] (11.712,3.932) circle (0.75pt);
\fill[navy] (11.712,4.218) circle (0.75pt);
\fill[navy] (11.712,4.361) circle (0.75pt);
\fill[navy] (11.712,4.647) circle (0.75pt);
\fill[navy] (11.712,5.076) circle (0.75pt);
\fill[navy] (11.712,5.219) circle (0.75pt);
\fill[navy] (11.712,5.505) circle (0.75pt);
\fill[navy] (11.712,5.648) circle (0.75pt);
\fill[navy] (11.712,5.934) circle (0.75pt);
\fill[navy] (11.712,6.077) circle (0.75pt);
\fill[navy] (11.712,6.363) circle (0.75pt);
\fill[navy] (11.712,6.506) circle (0.75pt);
\fill[navy] (11.834,0.214) circle (0.75pt);
\fill[navy] (11.834,0.643) circle (0.75pt);
\fill[navy] (11.834,0.858) circle (0.75pt);
\fill[navy] (11.834,1.93) circle (0.75pt);
\fill[navy] (11.834,2.145) circle (0.75pt);
\fill[navy] (11.834,2.288) circle (0.75pt);
\fill[navy] (11.834,2.502) circle (0.75pt);
\fill[navy] (11.834,2.717) circle (0.75pt);
\fill[navy] (11.834,2.788) circle (0.75pt);
\fill[navy] (11.834,2.86) circle (0.75pt);
\fill[navy] (11.834,3.003) circle (0.75pt);
\fill[navy] (11.834,3.074) circle (0.75pt);
\fill[navy] (11.834,3.146) circle (0.75pt);
\fill[navy] (11.834,3.289) circle (0.75pt);
\fill[navy] (11.834,3.503) circle (0.75pt);
\fill[navy] (11.834,3.646) circle (0.75pt);
\fill[navy] (11.834,3.718) circle (0.75pt);
\fill[navy] (11.834,3.789) circle (0.75pt);
\fill[navy] (11.834,3.932) circle (0.75pt);
\fill[navy] (11.834,4.004) circle (0.75pt);
\fill[navy] (11.834,4.075) circle (0.75pt);
\fill[navy] (11.834,4.147) circle (0.75pt);
\fill[navy] (11.834,4.218) circle (0.75pt);
\fill[navy] (11.834,4.29) circle (0.75pt);
\fill[navy] (11.834,4.361) circle (0.75pt);
\fill[navy] (11.834,4.433) circle (0.75pt);
\fill[navy] (11.834,4.504) circle (0.75pt);
\fill[navy] (11.834,4.576) circle (0.75pt);
\fill[navy] (11.834,4.647) circle (0.75pt);
\fill[navy] (11.834,4.79) circle (0.75pt);
\fill[navy] (11.834,4.862) circle (0.75pt);
\fill[navy] (11.834,4.933) circle (0.75pt);
\fill[navy] (11.834,5.005) circle (0.75pt);
\fill[navy] (11.834,5.076) circle (0.75pt);
\fill[navy] (11.834,5.148) circle (0.75pt);
\fill[navy] (11.834,5.291) circle (0.75pt);
\fill[navy] (11.834,5.577) circle (0.75pt);
\fill[navy] (11.834,5.648) circle (0.75pt);
\fill[navy] (11.834,5.72) circle (0.75pt);
\fill[navy] (11.834,5.791) circle (0.75pt);
\fill[navy] (11.834,5.863) circle (0.75pt);
\fill[navy] (11.834,6.006) circle (0.75pt);
\fill[navy] (11.834,6.077) circle (0.75pt);
\fill[navy] (11.834,6.149) circle (0.75pt);
\fill[navy] (11.834,6.22) circle (0.75pt);
\fill[navy] (11.834,6.363) circle (0.75pt);
\fill[navy] (11.834,6.435) circle (0.75pt);
\fill[navy] (11.834,6.578) circle (0.75pt);
\fill[navy] (11.834,6.721) circle (0.75pt);
\fill[navy] (11.834,6.864) circle (0.75pt);
\fill[navy] (11.956,0.357) circle (0.75pt);
\fill[navy] (11.956,0.643) circle (0.75pt);
\fill[navy] (11.956,0.786) circle (0.75pt);
\fill[navy] (11.956,0.929) circle (0.75pt);
\fill[navy] (11.956,1.787) circle (0.75pt);
\fill[navy] (11.956,1.93) circle (0.75pt);
\fill[navy] (11.956,2.359) circle (0.75pt);
\fill[navy] (11.956,2.645) circle (0.75pt);
\fill[navy] (11.956,2.931) circle (0.75pt);
\fill[navy] (11.956,3.217) circle (0.75pt);
\fill[navy] (11.956,3.646) circle (0.75pt);
\fill[navy] (11.956,3.789) circle (0.75pt);
\fill[navy] (11.956,4.075) circle (0.75pt);
\fill[navy] (11.956,4.79) circle (0.75pt);
\fill[navy] (11.956,4.933) circle (0.75pt);
\fill[navy] (11.956,5.076) circle (0.75pt);
\fill[navy] (11.956,5.219) circle (0.75pt);
\fill[navy] (11.956,5.362) circle (0.75pt);
\fill[navy] (11.956,5.648) circle (0.75pt);
\fill[navy] (11.956,5.791) circle (0.75pt);
\fill[navy] (11.956,5.934) circle (0.75pt);
\fill[navy] (11.956,6.077) circle (0.75pt);
\fill[navy] (11.956,6.22) circle (0.75pt);
\fill[navy] (11.956,6.363) circle (0.75pt);
\fill[navy] (11.956,6.792) circle (0.75pt);
\fill[navy] (12.078,0.071) circle (0.75pt);
\fill[navy] (12.078,0.143) circle (0.75pt);
\fill[navy] (12.078,0.286) circle (0.75pt);
\fill[navy] (12.078,0.357) circle (0.75pt);
\fill[navy] (12.078,0.5) circle (0.75pt);
\fill[navy] (12.078,0.572) circle (0.75pt);
\fill[navy] (12.078,0.715) circle (0.75pt);
\fill[navy] (12.078,0.929) circle (0.75pt);
\fill[navy] (12.078,1.001) circle (0.75pt);
\fill[navy] (12.078,1.144) circle (0.75pt);
\fill[navy] (12.078,1.43) circle (0.75pt);
\fill[navy] (12.078,1.644) circle (0.75pt);
\fill[navy] (12.078,1.859) circle (0.75pt);
\fill[navy] (12.078,2.002) circle (0.75pt);
\fill[navy] (12.078,2.645) circle (0.75pt);
\fill[navy] (12.078,2.86) circle (0.75pt);
\fill[navy] (12.078,3.074) circle (0.75pt);
\fill[navy] (12.078,3.575) circle (0.75pt);
\fill[navy] (12.078,3.718) circle (0.75pt);
\fill[navy] (12.078,4.147) circle (0.75pt);
\fill[navy] (12.078,4.218) circle (0.75pt);
\fill[navy] (12.078,4.361) circle (0.75pt);
\fill[navy] (12.078,4.433) circle (0.75pt);
\fill[navy] (12.078,4.576) circle (0.75pt);
\fill[navy] (12.078,4.647) circle (0.75pt);
\fill[navy] (12.078,4.79) circle (0.75pt);
\fill[navy] (12.078,4.862) circle (0.75pt);
\fill[navy] (12.078,5.076) circle (0.75pt);
\fill[navy] (12.078,5.219) circle (0.75pt);
\fill[navy] (12.078,5.434) circle (0.75pt);
\fill[navy] (12.078,5.648) circle (0.75pt);
\fill[navy] (12.078,5.72) circle (0.75pt);
\fill[navy] (12.078,5.863) circle (0.75pt);
\fill[navy] (12.078,5.934) circle (0.75pt);
\fill[navy] (12.078,6.077) circle (0.75pt);
\fill[navy] (12.078,6.149) circle (0.75pt);
\fill[navy] (12.078,6.363) circle (0.75pt);
\fill[navy] (12.078,6.506) circle (0.75pt);
\fill[navy] (12.078,6.578) circle (0.75pt);
\fill[navy] (12.078,6.721) circle (0.75pt);
\fill[navy] (12.078,6.935) circle (0.75pt);
\fill[navy] (12.078,7.007) circle (0.75pt);
\fill[navy] (12.2,0.643) circle (0.75pt);
\fill[navy] (12.2,0.786) circle (0.75pt);
\fill[navy] (12.2,1.215) circle (0.75pt);
\fill[navy] (12.2,1.644) circle (0.75pt);
\fill[navy] (12.2,1.93) circle (0.75pt);
\fill[navy] (12.2,2.216) circle (0.75pt);
\fill[navy] (12.2,2.359) circle (0.75pt);
\fill[navy] (12.2,2.645) circle (0.75pt);
\fill[navy] (12.2,2.931) circle (0.75pt);
\fill[navy] (12.2,3.074) circle (0.75pt);
\fill[navy] (12.2,3.36) circle (0.75pt);
\fill[navy] (12.2,3.646) circle (0.75pt);
\fill[navy] (12.2,3.789) circle (0.75pt);
\fill[navy] (12.2,4.075) circle (0.75pt);
\fill[navy] (12.2,4.218) circle (0.75pt);
\fill[navy] (12.2,4.504) circle (0.75pt);
\fill[navy] (12.2,4.79) circle (0.75pt);
\fill[navy] (12.2,4.933) circle (0.75pt);
\fill[navy] (12.2,5.219) circle (0.75pt);
\fill[navy] (12.2,5.505) circle (0.75pt);
\fill[navy] (12.2,5.648) circle (0.75pt);
\fill[navy] (12.2,5.791) circle (0.75pt);
\fill[navy] (12.2,5.934) circle (0.75pt);
\fill[navy] (12.2,6.22) circle (0.75pt);
\fill[navy] (12.2,6.363) circle (0.75pt);
\fill[navy] (12.2,6.506) circle (0.75pt);
\fill[navy] (12.2,6.649) circle (0.75pt);
\fill[navy] (12.2,6.935) circle (0.75pt);
\fill[navy] (12.2,7.078) circle (0.75pt);
\fill[teal] (0.976,0.357) circle (1.05pt);
\fill[teal] (1.22,0.5) circle (1.05pt);
\fill[teal] (1.342,0.286) circle (1.05pt);
\fill[teal] (1.342,0.643) circle (1.05pt);
\fill[teal] (1.586,0.429) circle (1.05pt);
\fill[teal] (1.586,0.643) circle (1.05pt);
\fill[teal] (1.83,0.5) circle (1.05pt);
\fill[teal] (1.952,0.643) circle (1.05pt);
\fill[teal] (1.952,0.786) circle (1.05pt);
\fill[teal] (1.952,1.073) circle (1.05pt);
\fill[teal] (2.074,0.357) circle (1.05pt);
\fill[teal] (2.074,0.5) circle (1.05pt);
\fill[teal] (2.074,0.572) circle (1.05pt);
\fill[teal] (2.074,0.643) circle (1.05pt);
\fill[teal] (2.074,0.786) circle (1.05pt);
\fill[teal] (2.074,0.858) circle (1.05pt);
\fill[teal] (2.074,1.001) circle (1.05pt);
\fill[teal] (2.196,0.5) circle (1.05pt);
\fill[teal] (2.196,0.786) circle (1.05pt);
\fill[teal] (2.196,0.929) circle (1.05pt);
\fill[teal] (2.318,0.286) circle (1.05pt);
\fill[teal] (2.318,0.643) circle (1.05pt);
\fill[teal] (2.318,0.786) circle (1.05pt);
\fill[teal] (2.318,0.858) circle (1.05pt);
\fill[teal] (2.318,0.929) circle (1.05pt);
\fill[teal] (2.318,1.001) circle (1.05pt);
\fill[teal] (2.318,1.215) circle (1.05pt);
\fill[teal] (2.318,1.287) circle (1.05pt);
\fill[teal] (2.44,0.071) circle (1.05pt);
\fill[teal] (2.44,0.643) circle (1.05pt);
\fill[teal] (2.44,1.358) circle (1.05pt);
\fill[teal] (2.562,0.143) circle (1.05pt);
\fill[teal] (2.562,0.286) circle (1.05pt);
\fill[teal] (2.562,1.358) circle (1.05pt);
\fill[teal] (2.684,0.929) circle (1.05pt);
\fill[teal] (2.684,1.215) circle (1.05pt);
\fill[teal] (2.684,1.358) circle (1.05pt);
\fill[teal] (2.684,1.501) circle (1.05pt);
\fill[teal] (2.806,0.143) circle (1.05pt);
\fill[teal] (2.806,0.572) circle (1.05pt);
\fill[teal] (2.806,0.643) circle (1.05pt);
\fill[teal] (2.806,0.715) circle (1.05pt);
\fill[teal] (2.806,0.786) circle (1.05pt);
\fill[teal] (2.806,0.858) circle (1.05pt);
\fill[teal] (2.806,0.929) circle (1.05pt);
\fill[teal] (2.806,1.001) circle (1.05pt);
\fill[teal] (2.806,1.073) circle (1.05pt);
\fill[teal] (2.806,1.144) circle (1.05pt);
\fill[teal] (2.806,1.287) circle (1.05pt);
\fill[teal] (2.806,1.358) circle (1.05pt);
\fill[teal] (2.806,1.43) circle (1.05pt);
\fill[teal] (2.806,1.501) circle (1.05pt);
\fill[teal] (2.928,0.5) circle (1.05pt);
\fill[teal] (2.928,0.786) circle (1.05pt);
\fill[teal] (2.928,1.215) circle (1.05pt);
\fill[teal] (2.928,1.358) circle (1.05pt);
\fill[teal] (2.928,1.644) circle (1.05pt);
\fill[teal] (3.05,0.429) circle (1.05pt);
\fill[teal] (3.05,0.643) circle (1.05pt);
\fill[teal] (3.05,0.786) circle (1.05pt);
\fill[teal] (3.05,0.858) circle (1.05pt);
\fill[teal] (3.05,0.929) circle (1.05pt);
\fill[teal] (3.05,1.001) circle (1.05pt);
\fill[teal] (3.05,1.144) circle (1.05pt);
\fill[teal] (3.05,1.215) circle (1.05pt);
\fill[teal] (3.05,1.358) circle (1.05pt);
\fill[teal] (3.05,1.573) circle (1.05pt);
\fill[teal] (3.05,1.716) circle (1.05pt);
\fill[teal] (3.172,0.071) circle (1.05pt);
\fill[teal] (3.172,0.357) circle (1.05pt);
\fill[teal] (3.172,0.5) circle (1.05pt);
\fill[teal] (3.172,0.643) circle (1.05pt);
\fill[teal] (3.172,0.786) circle (1.05pt);
\fill[teal] (3.172,1.073) circle (1.05pt);
\fill[teal] (3.172,1.215) circle (1.05pt);
\fill[teal] (3.172,1.358) circle (1.05pt);
\fill[teal] (3.172,1.787) circle (1.05pt);
\fill[teal] (3.294,0.143) circle (1.05pt);
\fill[teal] (3.294,0.286) circle (1.05pt);
\fill[teal] (3.294,0.357) circle (1.05pt);
\fill[teal] (3.294,0.5) circle (1.05pt);
\fill[teal] (3.294,0.715) circle (1.05pt);
\fill[teal] (3.294,0.786) circle (1.05pt);
\fill[teal] (3.294,1.001) circle (1.05pt);
\fill[teal] (3.294,1.144) circle (1.05pt);
\fill[teal] (3.294,1.215) circle (1.05pt);
\fill[teal] (3.294,1.43) circle (1.05pt);
\fill[teal] (3.294,1.573) circle (1.05pt);
\fill[teal] (3.294,1.644) circle (1.05pt);
\fill[teal] (3.416,0.214) circle (1.05pt);
\fill[teal] (3.416,0.357) circle (1.05pt);
\fill[teal] (3.416,0.643) circle (1.05pt);
\fill[teal] (3.416,0.786) circle (1.05pt);
\fill[teal] (3.416,0.929) circle (1.05pt);
\fill[teal] (3.416,1.073) circle (1.05pt);
\fill[teal] (3.416,1.358) circle (1.05pt);
\fill[teal] (3.416,1.644) circle (1.05pt);
\fill[teal] (3.416,1.787) circle (1.05pt);
\fill[teal] (3.416,1.93) circle (1.05pt);
\fill[teal] (3.538,0.071) circle (1.05pt);
\fill[teal] (3.538,0.214) circle (1.05pt);
\fill[teal] (3.538,0.286) circle (1.05pt);
\fill[teal] (3.538,0.429) circle (1.05pt);
\fill[teal] (3.538,0.5) circle (1.05pt);
\fill[teal] (3.538,0.572) circle (1.05pt);
\fill[teal] (3.538,0.643) circle (1.05pt);
\fill[teal] (3.538,0.715) circle (1.05pt);
\fill[teal] (3.538,0.786) circle (1.05pt);
\fill[teal] (3.538,0.858) circle (1.05pt);
\fill[teal] (3.538,0.929) circle (1.05pt);
\fill[teal] (3.538,1.001) circle (1.05pt);
\fill[teal] (3.538,1.073) circle (1.05pt);
\fill[teal] (3.538,1.144) circle (1.05pt);
\fill[teal] (3.538,1.215) circle (1.05pt);
\fill[teal] (3.538,1.287) circle (1.05pt);
\fill[teal] (3.538,1.358) circle (1.05pt);
\fill[teal] (3.538,1.501) circle (1.05pt);
\fill[teal] (3.538,1.573) circle (1.05pt);
\fill[teal] (3.538,1.644) circle (1.05pt);
\fill[teal] (3.538,1.716) circle (1.05pt);
\fill[teal] (3.538,1.787) circle (1.05pt);
\fill[teal] (3.538,1.93) circle (1.05pt);
\fill[teal] (3.538,2.002) circle (1.05pt);
\fill[teal] (3.66,0.071) circle (1.05pt);
\fill[teal] (3.66,0.5) circle (1.05pt);
\fill[teal] (3.66,0.786) circle (1.05pt);
\fill[teal] (3.66,1.358) circle (1.05pt);
\fill[teal] (3.782,0.143) circle (1.05pt);
\fill[teal] (3.782,0.214) circle (1.05pt);
\fill[teal] (3.782,0.286) circle (1.05pt);
\fill[teal] (3.782,0.357) circle (1.05pt);
\fill[teal] (3.782,0.5) circle (1.05pt);
\fill[teal] (3.782,0.572) circle (1.05pt);
\fill[teal] (3.782,0.715) circle (1.05pt);
\fill[teal] (3.782,0.858) circle (1.05pt);
\fill[teal] (3.782,0.929) circle (1.05pt);
\fill[teal] (3.782,1.001) circle (1.05pt);
\fill[teal] (3.782,1.073) circle (1.05pt);
\fill[teal] (3.782,1.144) circle (1.05pt);
\fill[teal] (3.782,1.215) circle (1.05pt);
\fill[teal] (3.782,1.287) circle (1.05pt);
\fill[teal] (3.782,1.358) circle (1.05pt);
\fill[teal] (3.782,1.43) circle (1.05pt);
\fill[teal] (3.782,1.501) circle (1.05pt);
\fill[teal] (3.782,1.573) circle (1.05pt);
\fill[teal] (3.782,1.644) circle (1.05pt);
\fill[teal] (3.782,1.716) circle (1.05pt);
\fill[teal] (3.782,1.787) circle (1.05pt);
\fill[teal] (3.782,1.859) circle (1.05pt);
\fill[teal] (3.782,1.93) circle (1.05pt);
\fill[teal] (3.782,2.002) circle (1.05pt);
\fill[teal] (3.782,2.145) circle (1.05pt);
\fill[teal] (3.904,0.071) circle (1.05pt);
\fill[teal] (3.904,0.643) circle (1.05pt);
\fill[teal] (3.904,0.786) circle (1.05pt);
\fill[teal] (3.904,0.929) circle (1.05pt);
\fill[teal] (3.904,1.073) circle (1.05pt);
\fill[teal] (3.904,1.501) circle (1.05pt);
\fill[teal] (3.904,1.787) circle (1.05pt);
\fill[teal] (3.904,1.93) circle (1.05pt);
\fill[teal] (3.904,2.073) circle (1.05pt);
\fill[teal] (3.904,2.216) circle (1.05pt);
\fill[teal] (4.026,0.071) circle (1.05pt);
\fill[teal] (4.026,0.286) circle (1.05pt);
\fill[teal] (4.026,0.357) circle (1.05pt);
\fill[teal] (4.026,0.5) circle (1.05pt);
\fill[teal] (4.026,0.572) circle (1.05pt);
\fill[teal] (4.026,0.715) circle (1.05pt);
\fill[teal] (4.026,0.929) circle (1.05pt);
\fill[teal] (4.026,1.001) circle (1.05pt);
\fill[teal] (4.026,1.215) circle (1.05pt);
\fill[teal] (4.026,1.358) circle (1.05pt);
\fill[teal] (4.026,1.644) circle (1.05pt);
\fill[teal] (4.026,1.787) circle (1.05pt);
\fill[teal] (4.026,1.859) circle (1.05pt);
\fill[teal] (4.026,2.002) circle (1.05pt);
\fill[teal] (4.026,2.073) circle (1.05pt);
\fill[teal] (4.026,2.288) circle (1.05pt);
\fill[teal] (4.148,0.071) circle (1.05pt);
\fill[teal] (4.148,0.214) circle (1.05pt);
\fill[teal] (4.148,0.643) circle (1.05pt);
\fill[teal] (4.148,0.786) circle (1.05pt);
\fill[teal] (4.148,0.929) circle (1.05pt);
\fill[teal] (4.148,1.073) circle (1.05pt);
\fill[teal] (4.148,1.358) circle (1.05pt);
\fill[teal] (4.148,1.501) circle (1.05pt);
\fill[teal] (4.148,1.644) circle (1.05pt);
\fill[teal] (4.148,1.93) circle (1.05pt);
\fill[teal] (4.27,0.071) circle (1.05pt);
\fill[teal] (4.27,0.143) circle (1.05pt);
\fill[teal] (4.27,0.214) circle (1.05pt);
\fill[teal] (4.27,0.286) circle (1.05pt);
\fill[teal] (4.27,0.429) circle (1.05pt);
\fill[teal] (4.27,0.786) circle (1.05pt);
\fill[teal] (4.27,0.858) circle (1.05pt);
\fill[teal] (4.27,0.929) circle (1.05pt);
\fill[teal] (4.27,1.144) circle (1.05pt);
\fill[teal] (4.27,1.215) circle (1.05pt);
\fill[teal] (4.27,1.287) circle (1.05pt);
\fill[teal] (4.27,1.358) circle (1.05pt);
\fill[teal] (4.27,1.573) circle (1.05pt);
\fill[teal] (4.27,1.644) circle (1.05pt);
\fill[teal] (4.27,1.716) circle (1.05pt);
\fill[teal] (4.27,1.859) circle (1.05pt);
\fill[teal] (4.27,1.93) circle (1.05pt);
\fill[teal] (4.27,2.073) circle (1.05pt);
\fill[teal] (4.27,2.216) circle (1.05pt);
\fill[teal] (4.27,2.288) circle (1.05pt);
\fill[teal] (4.27,2.359) circle (1.05pt);
\fill[teal] (4.27,2.431) circle (1.05pt);
\fill[teal] (4.392,0.071) circle (1.05pt);
\fill[teal] (4.392,0.357) circle (1.05pt);
\fill[teal] (4.392,0.5) circle (1.05pt);
\fill[teal] (4.392,0.786) circle (1.05pt);
\fill[teal] (4.392,1.215) circle (1.05pt);
\fill[teal] (4.392,1.644) circle (1.05pt);
\fill[teal] (4.392,1.787) circle (1.05pt);
\fill[teal] (4.392,2.073) circle (1.05pt);
\fill[teal] (4.392,2.502) circle (1.05pt);
\fill[teal] (4.514,0.143) circle (1.05pt);
\fill[teal] (4.514,0.214) circle (1.05pt);
\fill[teal] (4.514,0.286) circle (1.05pt);
\fill[teal] (4.514,0.572) circle (1.05pt);
\fill[teal] (4.514,0.643) circle (1.05pt);
\fill[teal] (4.514,0.786) circle (1.05pt);
\fill[teal] (4.514,0.858) circle (1.05pt);
\fill[teal] (4.514,0.929) circle (1.05pt);
\fill[teal] (4.514,1.073) circle (1.05pt);
\fill[teal] (4.514,1.144) circle (1.05pt);
\fill[teal] (4.514,1.215) circle (1.05pt);
\fill[teal] (4.514,1.287) circle (1.05pt);
\fill[teal] (4.514,1.358) circle (1.05pt);
\fill[teal] (4.514,1.43) circle (1.05pt);
\fill[teal] (4.514,1.501) circle (1.05pt);
\fill[teal] (4.514,1.573) circle (1.05pt);
\fill[teal] (4.514,1.716) circle (1.05pt);
\fill[teal] (4.514,1.787) circle (1.05pt);
\fill[teal] (4.514,1.93) circle (1.05pt);
\fill[teal] (4.514,2.002) circle (1.05pt);
\fill[teal] (4.514,2.073) circle (1.05pt);
\fill[teal] (4.514,2.145) circle (1.05pt);
\fill[teal] (4.514,2.216) circle (1.05pt);
\fill[teal] (4.514,2.288) circle (1.05pt);
\fill[teal] (4.514,2.359) circle (1.05pt);
\fill[teal] (4.514,2.431) circle (1.05pt);
\fill[teal] (4.636,0.214) circle (1.05pt);
\fill[teal] (4.636,0.643) circle (1.05pt);
\fill[teal] (4.636,0.786) circle (1.05pt);
\fill[teal] (4.636,0.929) circle (1.05pt);
\fill[teal] (4.636,1.073) circle (1.05pt);
\fill[teal] (4.636,1.215) circle (1.05pt);
\fill[teal] (4.636,1.501) circle (1.05pt);
\fill[teal] (4.636,1.787) circle (1.05pt);
\fill[teal] (4.636,2.073) circle (1.05pt);
\fill[teal] (4.636,2.216) circle (1.05pt);
\fill[teal] (4.636,2.359) circle (1.05pt);
\fill[teal] (4.636,2.645) circle (1.05pt);
\fill[teal] (4.758,0.071) circle (1.05pt);
\fill[teal] (4.758,0.143) circle (1.05pt);
\fill[teal] (4.758,0.286) circle (1.05pt);
\fill[teal] (4.758,0.357) circle (1.05pt);
\fill[teal] (4.758,0.5) circle (1.05pt);
\fill[teal] (4.758,0.572) circle (1.05pt);
\fill[teal] (4.758,0.715) circle (1.05pt);
\fill[teal] (4.758,0.786) circle (1.05pt);
\fill[teal] (4.758,1.001) circle (1.05pt);
\fill[teal] (4.758,1.144) circle (1.05pt);
\fill[teal] (4.758,1.215) circle (1.05pt);
\fill[teal] (4.758,1.358) circle (1.05pt);
\fill[teal] (4.758,1.43) circle (1.05pt);
\fill[teal] (4.758,1.573) circle (1.05pt);
\fill[teal] (4.758,2.073) circle (1.05pt);
\fill[teal] (4.758,2.288) circle (1.05pt);
\fill[teal] (4.758,2.431) circle (1.05pt);
\fill[teal] (4.758,2.502) circle (1.05pt);
\fill[teal] (4.758,2.645) circle (1.05pt);
\fill[teal] (4.758,2.717) circle (1.05pt);
\fill[teal] (4.88,0.071) circle (1.05pt);
\fill[teal] (4.88,0.5) circle (1.05pt);
\fill[teal] (4.88,0.643) circle (1.05pt);
\fill[teal] (4.88,0.786) circle (1.05pt);
\fill[teal] (4.88,0.929) circle (1.05pt);
\fill[teal] (4.88,1.215) circle (1.05pt);
\fill[teal] (4.88,1.358) circle (1.05pt);
\fill[teal] (4.88,1.501) circle (1.05pt);
\fill[teal] (4.88,1.644) circle (1.05pt);
\fill[teal] (4.88,1.93) circle (1.05pt);
\fill[teal] (4.88,2.073) circle (1.05pt);
\fill[teal] (4.88,2.216) circle (1.05pt);
\fill[teal] (4.88,2.359) circle (1.05pt);
\fill[teal] (4.88,2.788) circle (1.05pt);
\fill[teal] (5.002,0.143) circle (1.05pt);
\fill[teal] (5.002,0.286) circle (1.05pt);
\fill[teal] (5.002,0.357) circle (1.05pt);
\fill[teal] (5.002,0.5) circle (1.05pt);
\fill[teal] (5.002,0.572) circle (1.05pt);
\fill[teal] (5.002,0.643) circle (1.05pt);
\fill[teal] (5.002,0.715) circle (1.05pt);
\fill[teal] (5.002,0.786) circle (1.05pt);
\fill[teal] (5.002,0.858) circle (1.05pt);
\fill[teal] (5.002,0.929) circle (1.05pt);
\fill[teal] (5.002,1.073) circle (1.05pt);
\fill[teal] (5.002,1.144) circle (1.05pt);
\fill[teal] (5.002,1.215) circle (1.05pt);
\fill[teal] (5.002,1.287) circle (1.05pt);
\fill[teal] (5.002,1.358) circle (1.05pt);
\fill[teal] (5.002,1.43) circle (1.05pt);
\fill[teal] (5.002,1.501) circle (1.05pt);
\fill[teal] (5.002,1.573) circle (1.05pt);
\fill[teal] (5.002,1.644) circle (1.05pt);
\fill[teal] (5.002,1.716) circle (1.05pt);
\fill[teal] (5.002,1.787) circle (1.05pt);
\fill[teal] (5.002,2.002) circle (1.05pt);
\fill[teal] (5.002,2.145) circle (1.05pt);
\fill[teal] (5.002,2.359) circle (1.05pt);
\fill[teal] (5.002,2.431) circle (1.05pt);
\fill[teal] (5.002,2.502) circle (1.05pt);
\fill[teal] (5.002,2.574) circle (1.05pt);
\fill[teal] (5.002,2.645) circle (1.05pt);
\fill[teal] (5.002,2.717) circle (1.05pt);
\fill[teal] (5.002,2.788) circle (1.05pt);
\fill[teal] (5.002,2.86) circle (1.05pt);
\fill[teal] (5.124,0.071) circle (1.05pt);
\fill[teal] (5.124,0.929) circle (1.05pt);
\fill[teal] (5.124,1.215) circle (1.05pt);
\fill[teal] (5.124,1.358) circle (1.05pt);
\fill[teal] (5.124,1.644) circle (1.05pt);
\fill[teal] (5.124,2.073) circle (1.05pt);
\fill[teal] (5.124,2.645) circle (1.05pt);
\fill[teal] (5.246,0.071) circle (1.05pt);
\fill[teal] (5.246,0.143) circle (1.05pt);
\fill[teal] (5.246,0.214) circle (1.05pt);
\fill[teal] (5.246,0.286) circle (1.05pt);
\fill[teal] (5.246,0.429) circle (1.05pt);
\fill[teal] (5.246,0.643) circle (1.05pt);
\fill[teal] (5.246,0.786) circle (1.05pt);
\fill[teal] (5.246,0.929) circle (1.05pt);
\fill[teal] (5.246,1.001) circle (1.05pt);
\fill[teal] (5.246,1.073) circle (1.05pt);
\fill[teal] (5.246,1.144) circle (1.05pt);
\fill[teal] (5.246,1.358) circle (1.05pt);
\fill[teal] (5.246,1.43) circle (1.05pt);
\fill[teal] (5.246,1.501) circle (1.05pt);
\fill[teal] (5.246,1.644) circle (1.05pt);
\fill[teal] (5.246,1.716) circle (1.05pt);
\fill[teal] (5.246,1.787) circle (1.05pt);
\fill[teal] (5.246,1.93) circle (1.05pt);
\fill[teal] (5.246,2.002) circle (1.05pt);
\fill[teal] (5.246,2.073) circle (1.05pt);
\fill[teal] (5.246,2.145) circle (1.05pt);
\fill[teal] (5.246,2.216) circle (1.05pt);
\fill[teal] (5.246,2.359) circle (1.05pt);
\fill[teal] (5.246,2.431) circle (1.05pt);
\fill[teal] (5.246,2.502) circle (1.05pt);
\fill[teal] (5.246,2.574) circle (1.05pt);
\fill[teal] (5.246,2.645) circle (1.05pt);
\fill[teal] (5.246,2.788) circle (1.05pt);
\fill[teal] (5.246,2.86) circle (1.05pt);
\fill[teal] (5.246,2.931) circle (1.05pt);
\fill[teal] (5.246,3.003) circle (1.05pt);
\fill[teal] (5.368,0.071) circle (1.05pt);
\fill[teal] (5.368,0.214) circle (1.05pt);
\fill[teal] (5.368,0.5) circle (1.05pt);
\fill[teal] (5.368,0.643) circle (1.05pt);
\fill[teal] (5.368,0.929) circle (1.05pt);
\fill[teal] (5.368,1.073) circle (1.05pt);
\fill[teal] (5.368,1.215) circle (1.05pt);
\fill[teal] (5.368,1.358) circle (1.05pt);
\fill[teal] (5.368,1.501) circle (1.05pt);
\fill[teal] (5.368,1.644) circle (1.05pt);
\fill[teal] (5.368,1.787) circle (1.05pt);
\fill[teal] (5.368,1.93) circle (1.05pt);
\fill[teal] (5.368,2.073) circle (1.05pt);
\fill[teal] (5.368,2.216) circle (1.05pt);
\fill[teal] (5.368,2.645) circle (1.05pt);
\fill[teal] (5.368,2.931) circle (1.05pt);
\fill[teal] (5.368,3.074) circle (1.05pt);
\fill[teal] (5.49,0.071) circle (1.05pt);
\fill[teal] (5.49,0.143) circle (1.05pt);
\fill[teal] (5.49,0.286) circle (1.05pt);
\fill[teal] (5.49,0.5) circle (1.05pt);
\fill[teal] (5.49,0.572) circle (1.05pt);
\fill[teal] (5.49,0.929) circle (1.05pt);
\fill[teal] (5.49,1.001) circle (1.05pt);
\fill[teal] (5.49,1.215) circle (1.05pt);
\fill[teal] (5.49,1.358) circle (1.05pt);
\fill[teal] (5.49,1.573) circle (1.05pt);
\fill[teal] (5.49,1.644) circle (1.05pt);
\fill[teal] (5.49,1.859) circle (1.05pt);
\fill[teal] (5.49,2.002) circle (1.05pt);
\fill[teal] (5.49,2.073) circle (1.05pt);
\fill[teal] (5.49,2.288) circle (1.05pt);
\fill[teal] (5.49,2.431) circle (1.05pt);
\fill[teal] (5.49,2.645) circle (1.05pt);
\fill[teal] (5.49,3.074) circle (1.05pt);
\fill[teal] (5.612,0.071) circle (1.05pt);
\fill[teal] (5.612,0.214) circle (1.05pt);
\fill[teal] (5.612,0.357) circle (1.05pt);
\fill[teal] (5.612,0.5) circle (1.05pt);
\fill[teal] (5.612,0.786) circle (1.05pt);
\fill[teal] (5.612,0.929) circle (1.05pt);
\fill[teal] (5.612,1.073) circle (1.05pt);
\fill[teal] (5.612,1.215) circle (1.05pt);
\fill[teal] (5.612,1.358) circle (1.05pt);
\fill[teal] (5.612,1.501) circle (1.05pt);
\fill[teal] (5.612,1.787) circle (1.05pt);
\fill[teal] (5.612,1.93) circle (1.05pt);
\fill[teal] (5.612,2.073) circle (1.05pt);
\fill[teal] (5.612,2.359) circle (1.05pt);
\fill[teal] (5.612,2.645) circle (1.05pt);
\fill[teal] (5.612,2.788) circle (1.05pt);
\fill[teal] (5.612,2.931) circle (1.05pt);
\fill[teal] (5.612,3.217) circle (1.05pt);
\fill[teal] (5.734,0.071) circle (1.05pt);
\fill[teal] (5.734,0.143) circle (1.05pt);
\fill[teal] (5.734,0.286) circle (1.05pt);
\fill[teal] (5.734,0.357) circle (1.05pt);
\fill[teal] (5.734,0.429) circle (1.05pt);
\fill[teal] (5.734,0.643) circle (1.05pt);
\fill[teal] (5.734,0.715) circle (1.05pt);
\fill[teal] (5.734,0.858) circle (1.05pt);
\fill[teal] (5.734,0.929) circle (1.05pt);
\fill[teal] (5.734,1.001) circle (1.05pt);
\fill[teal] (5.734,1.073) circle (1.05pt);
\fill[teal] (5.734,1.144) circle (1.05pt);
\fill[teal] (5.734,1.215) circle (1.05pt);
\fill[teal] (5.734,1.358) circle (1.05pt);
\fill[teal] (5.734,1.43) circle (1.05pt);
\fill[teal] (5.734,1.501) circle (1.05pt);
\fill[teal] (5.734,1.573) circle (1.05pt);
\fill[teal] (5.734,1.644) circle (1.05pt);
\fill[teal] (5.734,1.716) circle (1.05pt);
\fill[teal] (5.734,1.787) circle (1.05pt);
\fill[teal] (5.734,1.859) circle (1.05pt);
\fill[teal] (5.734,1.93) circle (1.05pt);
\fill[teal] (5.734,2.073) circle (1.05pt);
\fill[teal] (5.734,2.216) circle (1.05pt);
\fill[teal] (5.734,2.288) circle (1.05pt);
\fill[teal] (5.734,2.359) circle (1.05pt);
\fill[teal] (5.734,2.431) circle (1.05pt);
\fill[teal] (5.734,2.502) circle (1.05pt);
\fill[teal] (5.734,2.574) circle (1.05pt);
\fill[teal] (5.734,2.931) circle (1.05pt);
\fill[teal] (5.734,3.003) circle (1.05pt);
\fill[teal] (5.734,3.074) circle (1.05pt);
\fill[teal] (5.734,3.146) circle (1.05pt);
\fill[teal] (5.734,3.289) circle (1.05pt);
\fill[teal] (5.856,0.071) circle (1.05pt);
\fill[teal] (5.856,0.357) circle (1.05pt);
\fill[teal] (5.856,0.5) circle (1.05pt);
\fill[teal] (5.856,0.786) circle (1.05pt);
\fill[teal] (5.856,0.929) circle (1.05pt);
\fill[teal] (5.856,1.215) circle (1.05pt);
\fill[teal] (5.856,1.358) circle (1.05pt);
\fill[teal] (5.856,1.644) circle (1.05pt);
\fill[teal] (5.856,2.073) circle (1.05pt);
\fill[teal] (5.856,2.216) circle (1.05pt);
\fill[teal] (5.856,2.502) circle (1.05pt);
\fill[teal] (5.856,2.645) circle (1.05pt);
\fill[teal] (5.856,2.931) circle (1.05pt);
\fill[teal] (5.856,3.074) circle (1.05pt);
\fill[teal] (5.856,3.36) circle (1.05pt);
\fill[teal] (5.978,0.071) circle (1.05pt);
\fill[teal] (5.978,0.143) circle (1.05pt);
\fill[teal] (5.978,0.214) circle (1.05pt);
\fill[teal] (5.978,0.286) circle (1.05pt);
\fill[teal] (5.978,0.357) circle (1.05pt);
\fill[teal] (5.978,0.429) circle (1.05pt);
\fill[teal] (5.978,0.572) circle (1.05pt);
\fill[teal] (5.978,0.643) circle (1.05pt);
\fill[teal] (5.978,0.715) circle (1.05pt);
\fill[teal] (5.978,0.786) circle (1.05pt);
\fill[teal] (5.978,0.858) circle (1.05pt);
\fill[teal] (5.978,0.929) circle (1.05pt);
\fill[teal] (5.978,1.073) circle (1.05pt);
\fill[teal] (5.978,1.144) circle (1.05pt);
\fill[teal] (5.978,1.215) circle (1.05pt);
\fill[teal] (5.978,1.287) circle (1.05pt);
\fill[teal] (5.978,1.358) circle (1.05pt);
\fill[teal] (5.978,1.43) circle (1.05pt);
\fill[teal] (5.978,1.573) circle (1.05pt);
\fill[teal] (5.978,1.644) circle (1.05pt);
\fill[teal] (5.978,1.716) circle (1.05pt);
\fill[teal] (5.978,1.787) circle (1.05pt);
\fill[teal] (5.978,1.859) circle (1.05pt);
\fill[teal] (5.978,1.93) circle (1.05pt);
\fill[teal] (5.978,2.073) circle (1.05pt);
\fill[teal] (5.978,2.145) circle (1.05pt);
\fill[teal] (5.978,2.216) circle (1.05pt);
\fill[teal] (5.978,2.288) circle (1.05pt);
\fill[teal] (5.978,2.359) circle (1.05pt);
\fill[teal] (5.978,2.431) circle (1.05pt);
\fill[teal] (5.978,2.574) circle (1.05pt);
\fill[teal] (5.978,2.645) circle (1.05pt);
\fill[teal] (5.978,2.717) circle (1.05pt);
\fill[teal] (5.978,2.86) circle (1.05pt);
\fill[teal] (5.978,3.074) circle (1.05pt);
\fill[teal] (6.1,0.071) circle (1.05pt);
\fill[teal] (6.1,0.5) circle (1.05pt);
\fill[teal] (6.1,0.643) circle (1.05pt);
\fill[teal] (6.1,0.786) circle (1.05pt);
\fill[teal] (6.1,0.929) circle (1.05pt);
\fill[teal] (6.1,1.358) circle (1.05pt);
\fill[teal] (6.1,1.501) circle (1.05pt);
\fill[teal] (6.1,1.644) circle (1.05pt);
\fill[teal] (6.1,1.93) circle (1.05pt);
\fill[teal] (6.1,2.216) circle (1.05pt);
\fill[teal] (6.1,2.359) circle (1.05pt);
\fill[teal] (6.1,2.645) circle (1.05pt);
\fill[teal] (6.1,2.788) circle (1.05pt);
\fill[teal] (6.1,2.931) circle (1.05pt);
\fill[teal] (6.1,3.36) circle (1.05pt);
\fill[teal] (6.1,3.503) circle (1.05pt);
\fill[teal] (6.222,0.071) circle (1.05pt);
\fill[teal] (6.222,0.143) circle (1.05pt);
\fill[teal] (6.222,0.286) circle (1.05pt);
\fill[teal] (6.222,0.357) circle (1.05pt);
\fill[teal] (6.222,0.5) circle (1.05pt);
\fill[teal] (6.222,0.572) circle (1.05pt);
\fill[teal] (6.222,0.786) circle (1.05pt);
\fill[teal] (6.222,0.929) circle (1.05pt);
\fill[teal] (6.222,1.001) circle (1.05pt);
\fill[teal] (6.222,1.144) circle (1.05pt);
\fill[teal] (6.222,1.358) circle (1.05pt);
\fill[teal] (6.222,1.573) circle (1.05pt);
\fill[teal] (6.222,1.644) circle (1.05pt);
\fill[teal] (6.222,1.859) circle (1.05pt);
\fill[teal] (6.222,2.002) circle (1.05pt);
\fill[teal] (6.222,2.216) circle (1.05pt);
\fill[teal] (6.222,2.288) circle (1.05pt);
\fill[teal] (6.222,2.502) circle (1.05pt);
\fill[teal] (6.222,2.645) circle (1.05pt);
\fill[teal] (6.222,2.717) circle (1.05pt);
\fill[teal] (6.222,2.86) circle (1.05pt);
\fill[teal] (6.222,3.074) circle (1.05pt);
\fill[teal] (6.222,3.146) circle (1.05pt);
\fill[teal] (6.222,3.289) circle (1.05pt);
\fill[teal] (6.222,3.36) circle (1.05pt);
\fill[teal] (6.222,3.503) circle (1.05pt);
\fill[teal] (6.344,0.071) circle (1.05pt);
\fill[teal] (6.344,0.214) circle (1.05pt);
\fill[teal] (6.344,0.5) circle (1.05pt);
\fill[teal] (6.344,0.643) circle (1.05pt);
\fill[teal] (6.344,0.786) circle (1.05pt);
\fill[teal] (6.344,1.073) circle (1.05pt);
\fill[teal] (6.344,1.215) circle (1.05pt);
\fill[teal] (6.344,1.358) circle (1.05pt);
\fill[teal] (6.344,1.501) circle (1.05pt);
\fill[teal] (6.344,1.644) circle (1.05pt);
\fill[teal] (6.344,1.787) circle (1.05pt);
\fill[teal] (6.344,2.073) circle (1.05pt);
\fill[teal] (6.344,2.216) circle (1.05pt);
\fill[teal] (6.344,2.359) circle (1.05pt);
\fill[teal] (6.344,2.502) circle (1.05pt);
\fill[teal] (6.344,2.931) circle (1.05pt);
\fill[teal] (6.344,3.074) circle (1.05pt);
\fill[teal] (6.344,3.217) circle (1.05pt);
\fill[teal] (6.344,3.36) circle (1.05pt);
\fill[teal] (6.344,3.503) circle (1.05pt);
\fill[teal] (6.466,0.071) circle (1.05pt);
\fill[teal] (6.466,0.143) circle (1.05pt);
\fill[teal] (6.466,0.214) circle (1.05pt);
\fill[teal] (6.466,0.286) circle (1.05pt);
\fill[teal] (6.466,0.429) circle (1.05pt);
\fill[teal] (6.466,0.5) circle (1.05pt);
\fill[teal] (6.466,0.572) circle (1.05pt);
\fill[teal] (6.466,0.715) circle (1.05pt);
\fill[teal] (6.466,0.786) circle (1.05pt);
\fill[teal] (6.466,1.001) circle (1.05pt);
\fill[teal] (6.466,1.215) circle (1.05pt);
\fill[teal] (6.466,1.287) circle (1.05pt);
\fill[teal] (6.466,1.43) circle (1.05pt);
\fill[teal] (6.466,1.501) circle (1.05pt);
\fill[teal] (6.466,1.573) circle (1.05pt);
\fill[teal] (6.466,1.644) circle (1.05pt);
\fill[teal] (6.466,1.716) circle (1.05pt);
\fill[teal] (6.466,1.859) circle (1.05pt);
\fill[teal] (6.466,2.002) circle (1.05pt);
\fill[teal] (6.466,2.073) circle (1.05pt);
\fill[teal] (6.466,2.145) circle (1.05pt);
\fill[teal] (6.466,2.216) circle (1.05pt);
\fill[teal] (6.466,2.288) circle (1.05pt);
\fill[teal] (6.466,2.359) circle (1.05pt);
\fill[teal] (6.466,2.502) circle (1.05pt);
\fill[teal] (6.466,2.574) circle (1.05pt);
\fill[teal] (6.466,2.645) circle (1.05pt);
\fill[teal] (6.466,2.717) circle (1.05pt);
\fill[teal] (6.466,2.788) circle (1.05pt);
\fill[teal] (6.466,2.86) circle (1.05pt);
\fill[teal] (6.466,2.931) circle (1.05pt);
\fill[teal] (6.466,3.074) circle (1.05pt);
\fill[teal] (6.466,3.146) circle (1.05pt);
\fill[teal] (6.466,3.217) circle (1.05pt);
\fill[teal] (6.466,3.289) circle (1.05pt);
\fill[teal] (6.466,3.36) circle (1.05pt);
\fill[teal] (6.466,3.432) circle (1.05pt);
\fill[teal] (6.466,3.575) circle (1.05pt);
\fill[teal] (6.466,3.646) circle (1.05pt);
\fill[teal] (6.466,3.718) circle (1.05pt);
\fill[teal] (6.588,0.357) circle (1.05pt);
\fill[teal] (6.588,0.5) circle (1.05pt);
\fill[teal] (6.588,0.786) circle (1.05pt);
\fill[teal] (6.588,0.929) circle (1.05pt);
\fill[teal] (6.588,1.215) circle (1.05pt);
\fill[teal] (6.588,1.644) circle (1.05pt);
\fill[teal] (6.588,1.787) circle (1.05pt);
\fill[teal] (6.588,2.073) circle (1.05pt);
\fill[teal] (6.588,2.216) circle (1.05pt);
\fill[teal] (6.588,2.645) circle (1.05pt);
\fill[teal] (6.588,2.931) circle (1.05pt);
\fill[teal] (6.588,3.36) circle (1.05pt);
\fill[teal] (6.588,3.503) circle (1.05pt);
\fill[teal] (6.588,3.789) circle (1.05pt);
\fill[teal] (6.71,0.143) circle (1.05pt);
\fill[teal] (6.71,0.214) circle (1.05pt);
\fill[teal] (6.71,0.286) circle (1.05pt);
\fill[teal] (6.71,0.429) circle (1.05pt);
\fill[teal] (6.71,0.5) circle (1.05pt);
\fill[teal] (6.71,0.643) circle (1.05pt);
\fill[teal] (6.71,0.858) circle (1.05pt);
\fill[teal] (6.71,1.001) circle (1.05pt);
\fill[teal] (6.71,1.144) circle (1.05pt);
\fill[teal] (6.71,1.215) circle (1.05pt);
\fill[teal] (6.71,1.287) circle (1.05pt);
\fill[teal] (6.71,1.358) circle (1.05pt);
\fill[teal] (6.71,1.644) circle (1.05pt);
\fill[teal] (6.71,1.716) circle (1.05pt);
\fill[teal] (6.71,1.859) circle (1.05pt);
\fill[teal] (6.71,1.93) circle (1.05pt);
\fill[teal] (6.71,2.002) circle (1.05pt);
\fill[teal] (6.71,2.073) circle (1.05pt);
\fill[teal] (6.71,2.574) circle (1.05pt);
\fill[teal] (6.71,2.645) circle (1.05pt);
\fill[teal] (6.71,2.717) circle (1.05pt);
\fill[teal] (6.71,2.788) circle (1.05pt);
\fill[teal] (6.71,2.931) circle (1.05pt);
\fill[teal] (6.71,3.289) circle (1.05pt);
\fill[teal] (6.71,3.36) circle (1.05pt);
\fill[teal] (6.71,3.432) circle (1.05pt);
\fill[teal] (6.71,3.718) circle (1.05pt);
\fill[teal] (6.71,3.789) circle (1.05pt);
\fill[teal] (6.71,3.861) circle (1.05pt);
\fill[teal] (6.832,0.214) circle (1.05pt);
\fill[teal] (6.832,0.357) circle (1.05pt);
\fill[teal] (6.832,0.786) circle (1.05pt);
\fill[teal] (6.832,0.929) circle (1.05pt);
\fill[teal] (6.832,1.073) circle (1.05pt);
\fill[teal] (6.832,1.215) circle (1.05pt);
\fill[teal] (6.832,1.358) circle (1.05pt);
\fill[teal] (6.832,1.644) circle (1.05pt);
\fill[teal] (6.832,1.93) circle (1.05pt);
\fill[teal] (6.832,2.073) circle (1.05pt);
\fill[teal] (6.832,2.216) circle (1.05pt);
\fill[teal] (6.832,2.359) circle (1.05pt);
\fill[teal] (6.832,2.645) circle (1.05pt);
\fill[teal] (6.832,2.788) circle (1.05pt);
\fill[teal] (6.832,3.217) circle (1.05pt);
\fill[teal] (6.832,3.36) circle (1.05pt);
\fill[teal] (6.832,3.646) circle (1.05pt);
\fill[teal] (6.832,3.932) circle (1.05pt);
\fill[teal] (6.954,0.071) circle (1.05pt);
\fill[teal] (6.954,0.143) circle (1.05pt);
\fill[teal] (6.954,0.286) circle (1.05pt);
\fill[teal] (6.954,0.357) circle (1.05pt);
\fill[teal] (6.954,0.5) circle (1.05pt);
\fill[teal] (6.954,0.572) circle (1.05pt);
\fill[teal] (6.954,0.715) circle (1.05pt);
\fill[teal] (6.954,0.929) circle (1.05pt);
\fill[teal] (6.954,1.144) circle (1.05pt);
\fill[teal] (6.954,1.215) circle (1.05pt);
\fill[teal] (6.954,1.43) circle (1.05pt);
\fill[teal] (6.954,1.573) circle (1.05pt);
\fill[teal] (6.954,1.787) circle (1.05pt);
\fill[teal] (6.954,1.859) circle (1.05pt);
\fill[teal] (6.954,2.002) circle (1.05pt);
\fill[teal] (6.954,2.216) circle (1.05pt);
\fill[teal] (6.954,2.288) circle (1.05pt);
\fill[teal] (6.954,2.431) circle (1.05pt);
\fill[teal] (6.954,2.502) circle (1.05pt);
\fill[teal] (6.954,2.86) circle (1.05pt);
\fill[teal] (6.954,2.931) circle (1.05pt);
\fill[teal] (6.954,3.074) circle (1.05pt);
\fill[teal] (6.954,3.289) circle (1.05pt);
\fill[teal] (6.954,3.36) circle (1.05pt);
\fill[teal] (6.954,3.503) circle (1.05pt);
\fill[teal] (6.954,3.575) circle (1.05pt);
\fill[teal] (6.954,3.789) circle (1.05pt);
\fill[teal] (6.954,3.932) circle (1.05pt);
\fill[teal] (6.954,4.004) circle (1.05pt);
\fill[teal] (7.076,0.071) circle (1.05pt);
\fill[teal] (7.076,0.214) circle (1.05pt);
\fill[teal] (7.076,0.357) circle (1.05pt);
\fill[teal] (7.076,0.5) circle (1.05pt);
\fill[teal] (7.076,0.786) circle (1.05pt);
\fill[teal] (7.076,1.073) circle (1.05pt);
\fill[teal] (7.076,1.215) circle (1.05pt);
\fill[teal] (7.076,1.501) circle (1.05pt);
\fill[teal] (7.076,1.644) circle (1.05pt);
\fill[teal] (7.076,1.93) circle (1.05pt);
\fill[teal] (7.076,2.216) circle (1.05pt);
\fill[teal] (7.076,2.359) circle (1.05pt);
\fill[teal] (7.076,2.645) circle (1.05pt);
\fill[teal] (7.076,2.788) circle (1.05pt);
\fill[teal] (7.076,3.074) circle (1.05pt);
\fill[teal] (7.076,3.217) circle (1.05pt);
\fill[teal] (7.076,3.36) circle (1.05pt);
\fill[teal] (7.076,3.503) circle (1.05pt);
\fill[teal] (7.076,3.646) circle (1.05pt);
\fill[teal] (7.076,4.075) circle (1.05pt);
\fill[teal] (7.198,0.071) circle (1.05pt);
\fill[teal] (7.198,0.143) circle (1.05pt);
\fill[teal] (7.198,0.214) circle (1.05pt);
\fill[teal] (7.198,0.286) circle (1.05pt);
\fill[teal] (7.198,0.357) circle (1.05pt);
\fill[teal] (7.198,0.429) circle (1.05pt);
\fill[teal] (7.198,0.5) circle (1.05pt);
\fill[teal] (7.198,0.572) circle (1.05pt);
\fill[teal] (7.198,0.643) circle (1.05pt);
\fill[teal] (7.198,0.715) circle (1.05pt);
\fill[teal] (7.198,0.786) circle (1.05pt);
\fill[teal] (7.198,0.858) circle (1.05pt);
\fill[teal] (7.198,1.001) circle (1.05pt);
\fill[teal] (7.198,1.073) circle (1.05pt);
\fill[teal] (7.198,1.215) circle (1.05pt);
\fill[teal] (7.198,1.287) circle (1.05pt);
\fill[teal] (7.198,1.501) circle (1.05pt);
\fill[teal] (7.198,1.573) circle (1.05pt);
\fill[teal] (7.198,1.644) circle (1.05pt);
\fill[teal] (7.198,1.716) circle (1.05pt);
\fill[teal] (7.198,1.859) circle (1.05pt);
\fill[teal] (7.198,1.93) circle (1.05pt);
\fill[teal] (7.198,2.002) circle (1.05pt);
\fill[teal] (7.198,2.145) circle (1.05pt);
\fill[teal] (7.198,2.216) circle (1.05pt);
\fill[teal] (7.198,2.359) circle (1.05pt);
\fill[teal] (7.198,2.431) circle (1.05pt);
\fill[teal] (7.198,2.502) circle (1.05pt);
\fill[teal] (7.198,2.574) circle (1.05pt);
\fill[teal] (7.198,2.645) circle (1.05pt);
\fill[teal] (7.198,2.86) circle (1.05pt);
\fill[teal] (7.198,2.931) circle (1.05pt);
\fill[teal] (7.198,3.003) circle (1.05pt);
\fill[teal] (7.198,3.074) circle (1.05pt);
\fill[teal] (7.198,3.289) circle (1.05pt);
\fill[teal] (7.198,3.503) circle (1.05pt);
\fill[teal] (7.198,3.575) circle (1.05pt);
\fill[teal] (7.198,3.718) circle (1.05pt);
\fill[teal] (7.198,3.789) circle (1.05pt);
\fill[teal] (7.198,3.861) circle (1.05pt);
\fill[teal] (7.198,4.075) circle (1.05pt);
\fill[teal] (7.198,4.147) circle (1.05pt);
\fill[teal] (7.32,0.071) circle (1.05pt);
\fill[teal] (7.32,0.5) circle (1.05pt);
\fill[teal] (7.32,0.786) circle (1.05pt);
\fill[teal] (7.32,0.929) circle (1.05pt);
\fill[teal] (7.32,1.215) circle (1.05pt);
\fill[teal] (7.32,1.358) circle (1.05pt);
\fill[teal] (7.32,1.644) circle (1.05pt);
\fill[teal] (7.32,2.073) circle (1.05pt);
\fill[teal] (7.32,2.931) circle (1.05pt);
\fill[teal] (7.32,3.074) circle (1.05pt);
\fill[teal] (7.32,3.503) circle (1.05pt);
\fill[teal] (7.32,4.218) circle (1.05pt);
\fill[teal] (7.442,0.071) circle (1.05pt);
\fill[teal] (7.442,0.143) circle (1.05pt);
\fill[teal] (7.442,0.214) circle (1.05pt);
\fill[teal] (7.442,0.429) circle (1.05pt);
\fill[teal] (7.442,0.5) circle (1.05pt);
\fill[teal] (7.442,0.572) circle (1.05pt);
\fill[teal] (7.442,0.786) circle (1.05pt);
\fill[teal] (7.442,0.858) circle (1.05pt);
\fill[teal] (7.442,0.929) circle (1.05pt);
\fill[teal] (7.442,1.001) circle (1.05pt);
\fill[teal] (7.442,1.073) circle (1.05pt);
\fill[teal] (7.442,1.215) circle (1.05pt);
\fill[teal] (7.442,1.358) circle (1.05pt);
\fill[teal] (7.442,1.43) circle (1.05pt);
\fill[teal] (7.442,1.501) circle (1.05pt);
\fill[teal] (7.442,1.573) circle (1.05pt);
\fill[teal] (7.442,1.716) circle (1.05pt);
\fill[teal] (7.442,1.859) circle (1.05pt);
\fill[teal] (7.442,2.002) circle (1.05pt);
\fill[teal] (7.442,2.288) circle (1.05pt);
\fill[teal] (7.442,2.431) circle (1.05pt);
\fill[teal] (7.442,2.502) circle (1.05pt);
\fill[teal] (7.442,2.717) circle (1.05pt);
\fill[teal] (7.442,2.788) circle (1.05pt);
\fill[teal] (7.442,2.931) circle (1.05pt);
\fill[teal] (7.442,3.074) circle (1.05pt);
\fill[teal] (7.442,3.146) circle (1.05pt);
\fill[teal] (7.442,3.217) circle (1.05pt);
\fill[teal] (7.442,3.289) circle (1.05pt);
\fill[teal] (7.442,3.575) circle (1.05pt);
\fill[teal] (7.442,3.718) circle (1.05pt);
\fill[teal] (7.442,3.932) circle (1.05pt);
\fill[teal] (7.442,4.004) circle (1.05pt);
\fill[teal] (7.442,4.075) circle (1.05pt);
\fill[teal] (7.442,4.147) circle (1.05pt);
\fill[teal] (7.442,4.218) circle (1.05pt);
\fill[teal] (7.442,4.29) circle (1.05pt);
\fill[teal] (7.564,0.071) circle (1.05pt);
\fill[teal] (7.564,0.214) circle (1.05pt);
\fill[teal] (7.564,0.357) circle (1.05pt);
\fill[teal] (7.564,0.5) circle (1.05pt);
\fill[teal] (7.564,1.073) circle (1.05pt);
\fill[teal] (7.564,1.215) circle (1.05pt);
\fill[teal] (7.564,1.358) circle (1.05pt);
\fill[teal] (7.564,1.501) circle (1.05pt);
\fill[teal] (7.564,1.787) circle (1.05pt);
\fill[teal] (7.564,1.93) circle (1.05pt);
\fill[teal] (7.564,2.073) circle (1.05pt);
\fill[teal] (7.564,2.502) circle (1.05pt);
\fill[teal] (7.564,2.788) circle (1.05pt);
\fill[teal] (7.564,3.074) circle (1.05pt);
\fill[teal] (7.564,3.36) circle (1.05pt);
\fill[teal] (7.564,3.789) circle (1.05pt);
\fill[teal] (7.564,3.932) circle (1.05pt);
\fill[teal] (7.564,4.075) circle (1.05pt);
\fill[teal] (7.564,4.218) circle (1.05pt);
\fill[teal] (7.564,4.361) circle (1.05pt);
\fill[teal] (7.686,0.143) circle (1.05pt);
\fill[teal] (7.686,0.286) circle (1.05pt);
\fill[teal] (7.686,0.357) circle (1.05pt);
\fill[teal] (7.686,0.572) circle (1.05pt);
\fill[teal] (7.686,0.786) circle (1.05pt);
\fill[teal] (7.686,0.929) circle (1.05pt);
\fill[teal] (7.686,1.215) circle (1.05pt);
\fill[teal] (7.686,1.358) circle (1.05pt);
\fill[teal] (7.686,1.43) circle (1.05pt);
\fill[teal] (7.686,1.573) circle (1.05pt);
\fill[teal] (7.686,2.073) circle (1.05pt);
\fill[teal] (7.686,2.216) circle (1.05pt);
\fill[teal] (7.686,2.288) circle (1.05pt);
\fill[teal] (7.686,2.645) circle (1.05pt);
\fill[teal] (7.686,2.717) circle (1.05pt);
\fill[teal] (7.686,2.86) circle (1.05pt);
\fill[teal] (7.686,2.931) circle (1.05pt);
\fill[teal] (7.686,3.146) circle (1.05pt);
\fill[teal] (7.686,3.289) circle (1.05pt);
\fill[teal] (7.686,3.36) circle (1.05pt);
\fill[teal] (7.686,3.575) circle (1.05pt);
\fill[teal] (7.686,3.789) circle (1.05pt);
\fill[teal] (7.686,3.932) circle (1.05pt);
\fill[teal] (7.686,4.147) circle (1.05pt);
\fill[teal] (7.686,4.433) circle (1.05pt);
\fill[teal] (7.808,0.071) circle (1.05pt);
\fill[teal] (7.808,0.214) circle (1.05pt);
\fill[teal] (7.808,0.357) circle (1.05pt);
\fill[teal] (7.808,0.643) circle (1.05pt);
\fill[teal] (7.808,0.786) circle (1.05pt);
\fill[teal] (7.808,1.215) circle (1.05pt);
\fill[teal] (7.808,1.501) circle (1.05pt);
\fill[teal] (7.808,1.644) circle (1.05pt);
\fill[teal] (7.808,1.93) circle (1.05pt);
\fill[teal] (7.808,2.359) circle (1.05pt);
\fill[teal] (7.808,2.502) circle (1.05pt);
\fill[teal] (7.808,2.788) circle (1.05pt);
\fill[teal] (7.808,3.789) circle (1.05pt);
\fill[teal] (7.808,3.932) circle (1.05pt);
\fill[teal] (7.808,4.075) circle (1.05pt);
\fill[teal] (7.808,4.218) circle (1.05pt);
\fill[teal] (7.808,4.361) circle (1.05pt);
\fill[teal] (7.93,0.143) circle (1.05pt);
\fill[teal] (7.93,0.214) circle (1.05pt);
\fill[teal] (7.93,0.429) circle (1.05pt);
\fill[teal] (7.93,0.5) circle (1.05pt);
\fill[teal] (7.93,0.572) circle (1.05pt);
\fill[teal] (7.93,0.643) circle (1.05pt);
\fill[teal] (7.93,1.001) circle (1.05pt);
\fill[teal] (7.93,1.287) circle (1.05pt);
\fill[teal] (7.93,1.501) circle (1.05pt);
\fill[teal] (7.93,1.716) circle (1.05pt);
\fill[teal] (7.93,1.93) circle (1.05pt);
\fill[teal] (7.93,2.002) circle (1.05pt);
\fill[teal] (7.93,2.073) circle (1.05pt);
\fill[teal] (7.93,2.216) circle (1.05pt);
\fill[teal] (7.93,2.288) circle (1.05pt);
\fill[teal] (7.93,2.431) circle (1.05pt);
\fill[teal] (7.93,2.717) circle (1.05pt);
\fill[teal] (7.93,2.931) circle (1.05pt);
\fill[teal] (7.93,3.003) circle (1.05pt);
\fill[teal] (7.93,3.074) circle (1.05pt);
\fill[teal] (7.93,3.146) circle (1.05pt);
\fill[teal] (7.93,3.503) circle (1.05pt);
\fill[teal] (7.93,3.646) circle (1.05pt);
\fill[teal] (7.93,3.789) circle (1.05pt);
\fill[teal] (7.93,3.861) circle (1.05pt);
\fill[teal] (7.93,4.147) circle (1.05pt);
\fill[teal] (7.93,4.576) circle (1.05pt);
\fill[teal] (8.052,0.071) circle (1.05pt);
\fill[teal] (8.052,0.5) circle (1.05pt);
\fill[teal] (8.052,0.929) circle (1.05pt);
\fill[teal] (8.052,1.215) circle (1.05pt);
\fill[teal] (8.052,1.358) circle (1.05pt);
\fill[teal] (8.052,1.787) circle (1.05pt);
\fill[teal] (8.052,2.073) circle (1.05pt);
\fill[teal] (8.052,2.216) circle (1.05pt);
\fill[teal] (8.052,2.502) circle (1.05pt);
\fill[teal] (8.052,2.645) circle (1.05pt);
\fill[teal] (8.052,3.074) circle (1.05pt);
\fill[teal] (8.052,3.36) circle (1.05pt);
\fill[teal] (8.052,4.361) circle (1.05pt);
\fill[teal] (8.174,0.071) circle (1.05pt);
\fill[teal] (8.174,0.143) circle (1.05pt);
\fill[teal] (8.174,0.214) circle (1.05pt);
\fill[teal] (8.174,0.286) circle (1.05pt);
\fill[teal] (8.174,0.357) circle (1.05pt);
\fill[teal] (8.174,0.5) circle (1.05pt);
\fill[teal] (8.174,0.572) circle (1.05pt);
\fill[teal] (8.174,0.643) circle (1.05pt);
\fill[teal] (8.174,0.715) circle (1.05pt);
\fill[teal] (8.174,0.786) circle (1.05pt);
\fill[teal] (8.174,0.858) circle (1.05pt);
\fill[teal] (8.174,0.929) circle (1.05pt);
\fill[teal] (8.174,1.073) circle (1.05pt);
\fill[teal] (8.174,1.144) circle (1.05pt);
\fill[teal] (8.174,1.358) circle (1.05pt);
\fill[teal] (8.174,1.43) circle (1.05pt);
\fill[teal] (8.174,1.501) circle (1.05pt);
\fill[teal] (8.174,1.644) circle (1.05pt);
\fill[teal] (8.174,1.716) circle (1.05pt);
\fill[teal] (8.174,1.787) circle (1.05pt);
\fill[teal] (8.174,1.859) circle (1.05pt);
\fill[teal] (8.174,1.93) circle (1.05pt);
\fill[teal] (8.174,2.002) circle (1.05pt);
\fill[teal] (8.174,2.073) circle (1.05pt);
\fill[teal] (8.174,2.145) circle (1.05pt);
\fill[teal] (8.174,2.216) circle (1.05pt);
\fill[teal] (8.174,2.431) circle (1.05pt);
\fill[teal] (8.174,2.502) circle (1.05pt);
\fill[teal] (8.174,2.574) circle (1.05pt);
\fill[teal] (8.174,2.645) circle (1.05pt);
\fill[teal] (8.174,2.717) circle (1.05pt);
\fill[teal] (8.174,2.86) circle (1.05pt);
\fill[teal] (8.174,3.003) circle (1.05pt);
\fill[teal] (8.174,3.146) circle (1.05pt);
\fill[teal] (8.174,3.217) circle (1.05pt);
\fill[teal] (8.174,3.289) circle (1.05pt);
\fill[teal] (8.174,3.36) circle (1.05pt);
\fill[teal] (8.174,3.432) circle (1.05pt);
\fill[teal] (8.174,3.503) circle (1.05pt);
\fill[teal] (8.174,3.575) circle (1.05pt);
\fill[teal] (8.174,3.718) circle (1.05pt);
\fill[teal] (8.174,3.789) circle (1.05pt);
\fill[teal] (8.174,3.861) circle (1.05pt);
\fill[teal] (8.174,4.004) circle (1.05pt);
\fill[teal] (8.174,4.075) circle (1.05pt);
\fill[teal] (8.174,4.147) circle (1.05pt);
\fill[teal] (8.174,4.218) circle (1.05pt);
\fill[teal] (8.174,4.29) circle (1.05pt);
\fill[teal] (8.174,4.433) circle (1.05pt);
\fill[teal] (8.174,4.576) circle (1.05pt);
\fill[teal] (8.174,4.719) circle (1.05pt);
\fill[teal] (8.296,0.214) circle (1.05pt);
\fill[teal] (8.296,0.357) circle (1.05pt);
\fill[teal] (8.296,0.643) circle (1.05pt);
\fill[teal] (8.296,0.786) circle (1.05pt);
\fill[teal] (8.296,0.929) circle (1.05pt);
\fill[teal] (8.296,1.358) circle (1.05pt);
\fill[teal] (8.296,1.501) circle (1.05pt);
\fill[teal] (8.296,2.073) circle (1.05pt);
\fill[teal] (8.296,2.216) circle (1.05pt);
\fill[teal] (8.296,2.359) circle (1.05pt);
\fill[teal] (8.296,2.502) circle (1.05pt);
\fill[teal] (8.296,2.931) circle (1.05pt);
\fill[teal] (8.296,3.217) circle (1.05pt);
\fill[teal] (8.296,3.36) circle (1.05pt);
\fill[teal] (8.296,3.503) circle (1.05pt);
\fill[teal] (8.296,3.789) circle (1.05pt);
\fill[teal] (8.296,3.932) circle (1.05pt);
\fill[teal] (8.296,4.075) circle (1.05pt);
\fill[teal] (8.296,4.361) circle (1.05pt);
\fill[teal] (8.296,4.504) circle (1.05pt);
\fill[teal] (8.296,4.79) circle (1.05pt);
\fill[teal] (8.418,0.071) circle (1.05pt);
\fill[teal] (8.418,0.143) circle (1.05pt);
\fill[teal] (8.418,0.286) circle (1.05pt);
\fill[teal] (8.418,0.357) circle (1.05pt);
\fill[teal] (8.418,0.5) circle (1.05pt);
\fill[teal] (8.418,0.715) circle (1.05pt);
\fill[teal] (8.418,0.786) circle (1.05pt);
\fill[teal] (8.418,0.929) circle (1.05pt);
\fill[teal] (8.418,1.001) circle (1.05pt);
\fill[teal] (8.418,1.144) circle (1.05pt);
\fill[teal] (8.418,1.215) circle (1.05pt);
\fill[teal] (8.418,1.859) circle (1.05pt);
\fill[teal] (8.418,2.002) circle (1.05pt);
\fill[teal] (8.418,2.216) circle (1.05pt);
\fill[teal] (8.418,2.288) circle (1.05pt);
\fill[teal] (8.418,2.431) circle (1.05pt);
\fill[teal] (8.418,2.502) circle (1.05pt);
\fill[teal] (8.418,2.645) circle (1.05pt);
\fill[teal] (8.418,2.86) circle (1.05pt);
\fill[teal] (8.418,3.074) circle (1.05pt);
\fill[teal] (8.418,3.146) circle (1.05pt);
\fill[teal] (8.418,3.36) circle (1.05pt);
\fill[teal] (8.418,3.575) circle (1.05pt);
\fill[teal] (8.418,3.718) circle (1.05pt);
\fill[teal] (8.418,3.789) circle (1.05pt);
\fill[teal] (8.418,3.932) circle (1.05pt);
\fill[teal] (8.418,4.004) circle (1.05pt);
\fill[teal] (8.418,4.218) circle (1.05pt);
\fill[teal] (8.418,4.433) circle (1.05pt);
\fill[teal] (8.418,4.576) circle (1.05pt);
\fill[teal] (8.418,4.79) circle (1.05pt);
\fill[teal] (8.54,0.071) circle (1.05pt);
\fill[teal] (8.54,0.214) circle (1.05pt);
\fill[teal] (8.54,0.643) circle (1.05pt);
\fill[teal] (8.54,0.786) circle (1.05pt);
\fill[teal] (8.54,0.929) circle (1.05pt);
\fill[teal] (8.54,1.358) circle (1.05pt);
\fill[teal] (8.54,2.216) circle (1.05pt);
\fill[teal] (8.54,2.359) circle (1.05pt);
\fill[teal] (8.54,2.788) circle (1.05pt);
\fill[teal] (8.54,3.074) circle (1.05pt);
\fill[teal] (8.54,3.646) circle (1.05pt);
\fill[teal] (8.54,3.789) circle (1.05pt);
\fill[teal] (8.54,4.79) circle (1.05pt);
\fill[teal] (8.662,0.071) circle (1.05pt);
\fill[teal] (8.662,0.429) circle (1.05pt);
\fill[teal] (8.662,0.5) circle (1.05pt);
\fill[teal] (8.662,0.572) circle (1.05pt);
\fill[teal] (8.662,0.643) circle (1.05pt);
\fill[teal] (8.662,0.715) circle (1.05pt);
\fill[teal] (8.662,0.858) circle (1.05pt);
\fill[teal] (8.662,1.001) circle (1.05pt);
\fill[teal] (8.662,1.073) circle (1.05pt);
\fill[teal] (8.662,1.144) circle (1.05pt);
\fill[teal] (8.662,1.215) circle (1.05pt);
\fill[teal] (8.662,1.287) circle (1.05pt);
\fill[teal] (8.662,1.358) circle (1.05pt);
\fill[teal] (8.662,1.573) circle (1.05pt);
\fill[teal] (8.662,1.644) circle (1.05pt);
\fill[teal] (8.662,1.716) circle (1.05pt);
\fill[teal] (8.662,1.787) circle (1.05pt);
\fill[teal] (8.662,1.859) circle (1.05pt);
\fill[teal] (8.662,1.93) circle (1.05pt);
\fill[teal] (8.662,2.002) circle (1.05pt);
\fill[teal] (8.662,2.073) circle (1.05pt);
\fill[teal] (8.662,2.145) circle (1.05pt);
\fill[teal] (8.662,2.216) circle (1.05pt);
\fill[teal] (8.662,2.288) circle (1.05pt);
\fill[teal] (8.662,2.359) circle (1.05pt);
\fill[teal] (8.662,2.502) circle (1.05pt);
\fill[teal] (8.662,2.574) circle (1.05pt);
\fill[teal] (8.662,2.645) circle (1.05pt);
\fill[teal] (8.662,2.717) circle (1.05pt);
\fill[teal] (8.662,2.788) circle (1.05pt);
\fill[teal] (8.662,2.86) circle (1.05pt);
\fill[teal] (8.662,3.003) circle (1.05pt);
\fill[teal] (8.662,3.074) circle (1.05pt);
\fill[teal] (8.662,3.146) circle (1.05pt);
\fill[teal] (8.662,3.217) circle (1.05pt);
\fill[teal] (8.662,3.289) circle (1.05pt);
\fill[teal] (8.662,3.36) circle (1.05pt);
\fill[teal] (8.662,3.432) circle (1.05pt);
\fill[teal] (8.662,3.503) circle (1.05pt);
\fill[teal] (8.662,3.575) circle (1.05pt);
\fill[teal] (8.662,3.646) circle (1.05pt);
\fill[teal] (8.662,4.075) circle (1.05pt);
\fill[teal] (8.662,4.147) circle (1.05pt);
\fill[teal] (8.662,4.218) circle (1.05pt);
\fill[teal] (8.662,4.29) circle (1.05pt);
\fill[teal] (8.662,4.361) circle (1.05pt);
\fill[teal] (8.662,4.504) circle (1.05pt);
\fill[teal] (8.662,4.576) circle (1.05pt);
\fill[teal] (8.662,4.647) circle (1.05pt);
\fill[teal] (8.662,4.719) circle (1.05pt);
\fill[teal] (8.662,5.005) circle (1.05pt);
\fill[teal] (8.784,0.071) circle (1.05pt);
\fill[teal] (8.784,0.357) circle (1.05pt);
\fill[teal] (8.784,0.5) circle (1.05pt);
\fill[teal] (8.784,0.786) circle (1.05pt);
\fill[teal] (8.784,1.787) circle (1.05pt);
\fill[teal] (8.784,2.216) circle (1.05pt);
\fill[teal] (8.784,2.502) circle (1.05pt);
\fill[teal] (8.784,2.645) circle (1.05pt);
\fill[teal] (8.784,3.074) circle (1.05pt);
\fill[teal] (8.784,3.36) circle (1.05pt);
\fill[teal] (8.784,3.789) circle (1.05pt);
\fill[teal] (8.784,5.076) circle (1.05pt);
\fill[teal] (8.906,0.071) circle (1.05pt);
\fill[teal] (8.906,0.143) circle (1.05pt);
\fill[teal] (8.906,0.214) circle (1.05pt);
\fill[teal] (8.906,0.286) circle (1.05pt);
\fill[teal] (8.906,0.429) circle (1.05pt);
\fill[teal] (8.906,0.5) circle (1.05pt);
\fill[teal] (8.906,0.572) circle (1.05pt);
\fill[teal] (8.906,0.715) circle (1.05pt);
\fill[teal] (8.906,0.786) circle (1.05pt);
\fill[teal] (8.906,0.858) circle (1.05pt);
\fill[teal] (8.906,0.929) circle (1.05pt);
\fill[teal] (8.906,1.001) circle (1.05pt);
\fill[teal] (8.906,1.073) circle (1.05pt);
\fill[teal] (8.906,1.287) circle (1.05pt);
\fill[teal] (8.906,1.358) circle (1.05pt);
\fill[teal] (8.906,1.43) circle (1.05pt);
\fill[teal] (8.906,1.573) circle (1.05pt);
\fill[teal] (8.906,1.93) circle (1.05pt);
\fill[teal] (8.906,2.002) circle (1.05pt);
\fill[teal] (8.906,2.073) circle (1.05pt);
\fill[teal] (8.906,2.145) circle (1.05pt);
\fill[teal] (8.906,2.216) circle (1.05pt);
\fill[teal] (8.906,2.288) circle (1.05pt);
\fill[teal] (8.906,2.359) circle (1.05pt);
\fill[teal] (8.906,2.431) circle (1.05pt);
\fill[teal] (8.906,2.502) circle (1.05pt);
\fill[teal] (8.906,2.574) circle (1.05pt);
\fill[teal] (8.906,2.645) circle (1.05pt);
\fill[teal] (8.906,2.717) circle (1.05pt);
\fill[teal] (8.906,2.788) circle (1.05pt);
\fill[teal] (8.906,2.86) circle (1.05pt);
\fill[teal] (8.906,3.003) circle (1.05pt);
\fill[teal] (8.906,3.432) circle (1.05pt);
\fill[teal] (8.906,3.503) circle (1.05pt);
\fill[teal] (8.906,3.932) circle (1.05pt);
\fill[teal] (8.906,4.075) circle (1.05pt);
\fill[teal] (8.906,4.218) circle (1.05pt);
\fill[teal] (8.906,4.504) circle (1.05pt);
\fill[teal] (8.906,4.719) circle (1.05pt);
\fill[teal] (8.906,4.862) circle (1.05pt);
\fill[teal] (8.906,5.005) circle (1.05pt);
\fill[teal] (8.906,5.076) circle (1.05pt);
\fill[teal] (9.028,0.071) circle (1.05pt);
\fill[teal] (9.028,0.214) circle (1.05pt);
\fill[teal] (9.028,0.5) circle (1.05pt);
\fill[teal] (9.028,0.786) circle (1.05pt);
\fill[teal] (9.028,0.929) circle (1.05pt);
\fill[teal] (9.028,1.073) circle (1.05pt);
\fill[teal] (9.028,1.215) circle (1.05pt);
\fill[teal] (9.028,1.501) circle (1.05pt);
\fill[teal] (9.028,1.644) circle (1.05pt);
\fill[teal] (9.028,1.787) circle (1.05pt);
\fill[teal] (9.028,1.93) circle (1.05pt);
\fill[teal] (9.028,2.073) circle (1.05pt);
\fill[teal] (9.028,2.216) circle (1.05pt);
\fill[teal] (9.028,2.788) circle (1.05pt);
\fill[teal] (9.028,3.074) circle (1.05pt);
\fill[teal] (9.028,3.217) circle (1.05pt);
\fill[teal] (9.028,3.36) circle (1.05pt);
\fill[teal] (9.028,3.646) circle (1.05pt);
\fill[teal] (9.028,3.789) circle (1.05pt);
\fill[teal] (9.028,3.932) circle (1.05pt);
\fill[teal] (9.028,4.075) circle (1.05pt);
\fill[teal] (9.028,4.218) circle (1.05pt);
\fill[teal] (9.028,4.504) circle (1.05pt);
\fill[teal] (9.028,4.647) circle (1.05pt);
\fill[teal] (9.028,4.79) circle (1.05pt);
\fill[teal] (9.028,5.076) circle (1.05pt);
\fill[teal] (9.028,5.219) circle (1.05pt);
\fill[teal] (9.15,0.071) circle (1.05pt);
\fill[teal] (9.15,0.286) circle (1.05pt);
\fill[teal] (9.15,0.5) circle (1.05pt);
\fill[teal] (9.15,0.786) circle (1.05pt);
\fill[teal] (9.15,0.929) circle (1.05pt);
\fill[teal] (9.15,1.001) circle (1.05pt);
\fill[teal] (9.15,1.144) circle (1.05pt);
\fill[teal] (9.15,1.215) circle (1.05pt);
\fill[teal] (9.15,1.358) circle (1.05pt);
\fill[teal] (9.15,1.573) circle (1.05pt);
\fill[teal] (9.15,1.644) circle (1.05pt);
\fill[teal] (9.15,2.002) circle (1.05pt);
\fill[teal] (9.15,2.645) circle (1.05pt);
\fill[teal] (9.15,2.717) circle (1.05pt);
\fill[teal] (9.15,2.931) circle (1.05pt);
\fill[teal] (9.15,3.146) circle (1.05pt);
\fill[teal] (9.15,3.503) circle (1.05pt);
\fill[teal] (9.15,3.789) circle (1.05pt);
\fill[teal] (9.15,4.004) circle (1.05pt);
\fill[teal] (9.15,4.147) circle (1.05pt);
\fill[teal] (9.15,4.361) circle (1.05pt);
\fill[teal] (9.15,4.576) circle (1.05pt);
\fill[teal] (9.15,4.862) circle (1.05pt);
\fill[teal] (9.15,5.076) circle (1.05pt);
\fill[teal] (9.15,5.219) circle (1.05pt);
\fill[teal] (9.272,0.071) circle (1.05pt);
\fill[teal] (9.272,0.214) circle (1.05pt);
\fill[teal] (9.272,0.357) circle (1.05pt);
\fill[teal] (9.272,0.5) circle (1.05pt);
\fill[teal] (9.272,0.643) circle (1.05pt);
\fill[teal] (9.272,0.786) circle (1.05pt);
\fill[teal] (9.272,1.073) circle (1.05pt);
\fill[teal] (9.272,1.501) circle (1.05pt);
\fill[teal] (9.272,1.93) circle (1.05pt);
\fill[teal] (9.272,2.073) circle (1.05pt);
\fill[teal] (9.272,2.216) circle (1.05pt);
\fill[teal] (9.272,2.788) circle (1.05pt);
\fill[teal] (9.272,2.931) circle (1.05pt);
\fill[teal] (9.272,3.074) circle (1.05pt);
\fill[teal] (9.272,3.217) circle (1.05pt);
\fill[teal] (9.272,3.36) circle (1.05pt);
\fill[teal] (9.272,3.646) circle (1.05pt);
\fill[teal] (9.272,3.789) circle (1.05pt);
\fill[teal] (9.272,4.504) circle (1.05pt);
\fill[teal] (9.272,4.933) circle (1.05pt);
\fill[teal] (9.394,0.071) circle (1.05pt);
\fill[teal] (9.394,0.143) circle (1.05pt);
\fill[teal] (9.394,0.214) circle (1.05pt);
\fill[teal] (9.394,0.286) circle (1.05pt);
\fill[teal] (9.394,0.357) circle (1.05pt);
\fill[teal] (9.394,0.572) circle (1.05pt);
\fill[teal] (9.394,0.643) circle (1.05pt);
\fill[teal] (9.394,0.715) circle (1.05pt);
\fill[teal] (9.394,0.858) circle (1.05pt);
\fill[teal] (9.394,0.929) circle (1.05pt);
\fill[teal] (9.394,1.144) circle (1.05pt);
\fill[teal] (9.394,1.215) circle (1.05pt);
\fill[teal] (9.394,1.287) circle (1.05pt);
\fill[teal] (9.394,1.358) circle (1.05pt);
\fill[teal] (9.394,1.644) circle (1.05pt);
\fill[teal] (9.394,1.716) circle (1.05pt);
\fill[teal] (9.394,1.787) circle (1.05pt);
\fill[teal] (9.394,1.859) circle (1.05pt);
\fill[teal] (9.394,2.145) circle (1.05pt);
\fill[teal] (9.394,2.574) circle (1.05pt);
\fill[teal] (9.394,2.645) circle (1.05pt);
\fill[teal] (9.394,2.86) circle (1.05pt);
\fill[teal] (9.394,2.931) circle (1.05pt);
\fill[teal] (9.394,3.074) circle (1.05pt);
\fill[teal] (9.394,3.217) circle (1.05pt);
\fill[teal] (9.394,3.36) circle (1.05pt);
\fill[teal] (9.394,3.575) circle (1.05pt);
\fill[teal] (9.394,3.646) circle (1.05pt);
\fill[teal] (9.394,3.861) circle (1.05pt);
\fill[teal] (9.394,4.075) circle (1.05pt);
\fill[teal] (9.394,4.147) circle (1.05pt);
\fill[teal] (9.394,4.218) circle (1.05pt);
\fill[teal] (9.394,4.29) circle (1.05pt);
\fill[teal] (9.394,4.361) circle (1.05pt);
\fill[teal] (9.394,4.576) circle (1.05pt);
\fill[teal] (9.394,4.79) circle (1.05pt);
\fill[teal] (9.394,4.933) circle (1.05pt);
\fill[teal] (9.394,5.076) circle (1.05pt);
\fill[teal] (9.394,5.291) circle (1.05pt);
\fill[teal] (9.394,5.434) circle (1.05pt);
\fill[teal] (9.516,0.071) circle (1.05pt);
\fill[teal] (9.516,0.5) circle (1.05pt);
\fill[teal] (9.516,1.215) circle (1.05pt);
\fill[teal] (9.516,1.358) circle (1.05pt);
\fill[teal] (9.516,2.073) circle (1.05pt);
\fill[teal] (9.516,2.216) circle (1.05pt);
\fill[teal] (9.516,2.502) circle (1.05pt);
\fill[teal] (9.516,2.645) circle (1.05pt);
\fill[teal] (9.516,2.931) circle (1.05pt);
\fill[teal] (9.516,3.36) circle (1.05pt);
\fill[teal] (9.516,3.503) circle (1.05pt);
\fill[teal] (9.516,4.218) circle (1.05pt);
\fill[teal] (9.516,5.219) circle (1.05pt);
\fill[teal] (9.516,5.505) circle (1.05pt);
\fill[teal] (9.638,0.071) circle (1.05pt);
\fill[teal] (9.638,0.143) circle (1.05pt);
\fill[teal] (9.638,0.214) circle (1.05pt);
\fill[teal] (9.638,0.286) circle (1.05pt);
\fill[teal] (9.638,0.357) circle (1.05pt);
\fill[teal] (9.638,0.429) circle (1.05pt);
\fill[teal] (9.638,0.5) circle (1.05pt);
\fill[teal] (9.638,0.572) circle (1.05pt);
\fill[teal] (9.638,0.715) circle (1.05pt);
\fill[teal] (9.638,0.786) circle (1.05pt);
\fill[teal] (9.638,0.858) circle (1.05pt);
\fill[teal] (9.638,1.073) circle (1.05pt);
\fill[teal] (9.638,1.144) circle (1.05pt);
\fill[teal] (9.638,1.215) circle (1.05pt);
\fill[teal] (9.638,1.43) circle (1.05pt);
\fill[teal] (9.638,1.501) circle (1.05pt);
\fill[teal] (9.638,1.644) circle (1.05pt);
\fill[teal] (9.638,1.716) circle (1.05pt);
\fill[teal] (9.638,1.787) circle (1.05pt);
\fill[teal] (9.638,1.93) circle (1.05pt);
\fill[teal] (9.638,2.002) circle (1.05pt);
\fill[teal] (9.638,2.145) circle (1.05pt);
\fill[teal] (9.638,2.359) circle (1.05pt);
\fill[teal] (9.638,2.574) circle (1.05pt);
\fill[teal] (9.638,2.645) circle (1.05pt);
\fill[teal] (9.638,2.788) circle (1.05pt);
\fill[teal] (9.638,2.86) circle (1.05pt);
\fill[teal] (9.638,3.003) circle (1.05pt);
\fill[teal] (9.638,3.146) circle (1.05pt);
\fill[teal] (9.638,3.217) circle (1.05pt);
\fill[teal] (9.638,3.289) circle (1.05pt);
\fill[teal] (9.638,3.36) circle (1.05pt);
\fill[teal] (9.638,3.503) circle (1.05pt);
\fill[teal] (9.638,3.718) circle (1.05pt);
\fill[teal] (9.638,3.789) circle (1.05pt);
\fill[teal] (9.638,3.932) circle (1.05pt);
\fill[teal] (9.638,4.004) circle (1.05pt);
\fill[teal] (9.638,4.075) circle (1.05pt);
\fill[teal] (9.638,4.647) circle (1.05pt);
\fill[teal] (9.638,4.79) circle (1.05pt);
\fill[teal] (9.638,4.933) circle (1.05pt);
\fill[teal] (9.638,5.148) circle (1.05pt);
\fill[teal] (9.638,5.291) circle (1.05pt);
\fill[teal] (9.638,5.362) circle (1.05pt);
\fill[teal] (9.76,0.071) circle (1.05pt);
\fill[teal] (9.76,0.214) circle (1.05pt);
\fill[teal] (9.76,0.5) circle (1.05pt);
\fill[teal] (9.76,0.643) circle (1.05pt);
\fill[teal] (9.76,0.786) circle (1.05pt);
\fill[teal] (9.76,1.215) circle (1.05pt);
\fill[teal] (9.76,1.93) circle (1.05pt);
\fill[teal] (9.76,2.073) circle (1.05pt);
\fill[teal] (9.76,2.216) circle (1.05pt);
\fill[teal] (9.76,2.788) circle (1.05pt);
\fill[teal] (9.76,2.931) circle (1.05pt);
\fill[teal] (9.76,3.074) circle (1.05pt);
\fill[teal] (9.76,3.503) circle (1.05pt);
\fill[teal] (9.76,4.075) circle (1.05pt);
\fill[teal] (9.76,4.361) circle (1.05pt);
\fill[teal] (9.76,4.933) circle (1.05pt);
\fill[teal] (9.76,5.648) circle (1.05pt);
\fill[teal] (9.882,0.286) circle (1.05pt);
\fill[teal] (9.882,0.357) circle (1.05pt);
\fill[teal] (9.882,0.5) circle (1.05pt);
\fill[teal] (9.882,0.572) circle (1.05pt);
\fill[teal] (9.882,0.715) circle (1.05pt);
\fill[teal] (9.882,0.786) circle (1.05pt);
\fill[teal] (9.882,1.001) circle (1.05pt);
\fill[teal] (9.882,1.573) circle (1.05pt);
\fill[teal] (9.882,1.787) circle (1.05pt);
\fill[teal] (9.882,2.073) circle (1.05pt);
\fill[teal] (9.882,2.431) circle (1.05pt);
\fill[teal] (9.882,2.931) circle (1.05pt);
\fill[teal] (9.882,3.36) circle (1.05pt);
\fill[teal] (9.882,3.503) circle (1.05pt);
\fill[teal] (9.882,3.575) circle (1.05pt);
\fill[teal] (9.882,3.789) circle (1.05pt);
\fill[teal] (9.882,3.932) circle (1.05pt);
\fill[teal] (9.882,4.147) circle (1.05pt);
\fill[teal] (9.882,4.218) circle (1.05pt);
\fill[teal] (9.882,4.361) circle (1.05pt);
\fill[teal] (9.882,4.79) circle (1.05pt);
\fill[teal] (9.882,4.862) circle (1.05pt);
\fill[teal] (9.882,5.005) circle (1.05pt);
\fill[teal] (9.882,5.076) circle (1.05pt);
\fill[teal] (9.882,5.219) circle (1.05pt);
\fill[teal] (9.882,5.291) circle (1.05pt);
\fill[teal] (9.882,5.72) circle (1.05pt);
\fill[teal] (10.004,0.071) circle (1.05pt);
\fill[teal] (10.004,0.214) circle (1.05pt);
\fill[teal] (10.004,0.357) circle (1.05pt);
\fill[teal] (10.004,0.643) circle (1.05pt);
\fill[teal] (10.004,0.786) circle (1.05pt);
\fill[teal] (10.004,0.929) circle (1.05pt);
\fill[teal] (10.004,1.073) circle (1.05pt);
\fill[teal] (10.004,1.215) circle (1.05pt);
\fill[teal] (10.004,1.358) circle (1.05pt);
\fill[teal] (10.004,1.501) circle (1.05pt);
\fill[teal] (10.004,1.644) circle (1.05pt);
\fill[teal] (10.004,1.787) circle (1.05pt);
\fill[teal] (10.004,1.93) circle (1.05pt);
\fill[teal] (10.004,2.216) circle (1.05pt);
\fill[teal] (10.004,2.502) circle (1.05pt);
\fill[teal] (10.004,2.645) circle (1.05pt);
\fill[teal] (10.004,2.788) circle (1.05pt);
\fill[teal] (10.004,3.217) circle (1.05pt);
\fill[teal] (10.004,3.503) circle (1.05pt);
\fill[teal] (10.004,3.646) circle (1.05pt);
\fill[teal] (10.004,3.789) circle (1.05pt);
\fill[teal] (10.004,4.075) circle (1.05pt);
\fill[teal] (10.004,4.218) circle (1.05pt);
\fill[teal] (10.004,4.933) circle (1.05pt);
\fill[teal] (10.004,5.219) circle (1.05pt);
\fill[teal] (10.004,5.362) circle (1.05pt);
\fill[teal] (10.126,0.071) circle (1.05pt);
\fill[teal] (10.126,0.143) circle (1.05pt);
\fill[teal] (10.126,0.214) circle (1.05pt);
\fill[teal] (10.126,0.286) circle (1.05pt);
\fill[teal] (10.126,0.357) circle (1.05pt);
\fill[teal] (10.126,0.429) circle (1.05pt);
\fill[teal] (10.126,0.5) circle (1.05pt);
\fill[teal] (10.126,0.572) circle (1.05pt);
\fill[teal] (10.126,0.715) circle (1.05pt);
\fill[teal] (10.126,0.786) circle (1.05pt);
\fill[teal] (10.126,0.858) circle (1.05pt);
\fill[teal] (10.126,0.929) circle (1.05pt);
\fill[teal] (10.126,1.073) circle (1.05pt);
\fill[teal] (10.126,1.215) circle (1.05pt);
\fill[teal] (10.126,1.358) circle (1.05pt);
\fill[teal] (10.126,1.43) circle (1.05pt);
\fill[teal] (10.126,1.501) circle (1.05pt);
\fill[teal] (10.126,1.573) circle (1.05pt);
\fill[teal] (10.126,1.787) circle (1.05pt);
\fill[teal] (10.126,1.859) circle (1.05pt);
\fill[teal] (10.126,1.93) circle (1.05pt);
\fill[teal] (10.126,2.002) circle (1.05pt);
\fill[teal] (10.126,2.145) circle (1.05pt);
\fill[teal] (10.126,2.431) circle (1.05pt);
\fill[teal] (10.126,2.574) circle (1.05pt);
\fill[teal] (10.126,2.645) circle (1.05pt);
\fill[teal] (10.126,2.717) circle (1.05pt);
\fill[teal] (10.126,2.788) circle (1.05pt);
\fill[teal] (10.126,2.86) circle (1.05pt);
\fill[teal] (10.126,2.931) circle (1.05pt);
\fill[teal] (10.126,3.003) circle (1.05pt);
\fill[teal] (10.126,3.146) circle (1.05pt);
\fill[teal] (10.126,3.289) circle (1.05pt);
\fill[teal] (10.126,3.36) circle (1.05pt);
\fill[teal] (10.126,3.646) circle (1.05pt);
\fill[teal] (10.126,3.789) circle (1.05pt);
\fill[teal] (10.126,3.861) circle (1.05pt);
\fill[teal] (10.126,3.932) circle (1.05pt);
\fill[teal] (10.126,4.075) circle (1.05pt);
\fill[teal] (10.126,4.218) circle (1.05pt);
\fill[teal] (10.126,4.361) circle (1.05pt);
\fill[teal] (10.126,4.576) circle (1.05pt);
\fill[teal] (10.126,5.148) circle (1.05pt);
\fill[teal] (10.126,5.291) circle (1.05pt);
\fill[teal] (10.126,5.362) circle (1.05pt);
\fill[teal] (10.126,5.505) circle (1.05pt);
\fill[teal] (10.126,5.791) circle (1.05pt);
\fill[teal] (10.126,5.863) circle (1.05pt);
\fill[teal] (10.248,0.071) circle (1.05pt);
\fill[teal] (10.248,0.786) circle (1.05pt);
\fill[teal] (10.248,0.929) circle (1.05pt);
\fill[teal] (10.248,1.215) circle (1.05pt);
\fill[teal] (10.248,1.644) circle (1.05pt);
\fill[teal] (10.248,1.787) circle (1.05pt);
\fill[teal] (10.248,2.073) circle (1.05pt);
\fill[teal] (10.248,2.645) circle (1.05pt);
\fill[teal] (10.248,3.074) circle (1.05pt);
\fill[teal] (10.248,3.36) circle (1.05pt);
\fill[teal] (10.248,3.789) circle (1.05pt);
\fill[teal] (10.248,4.218) circle (1.05pt);
\fill[teal] (10.37,0.071) circle (1.05pt);
\fill[teal] (10.37,0.214) circle (1.05pt);
\fill[teal] (10.37,0.429) circle (1.05pt);
\fill[teal] (10.37,0.5) circle (1.05pt);
\fill[teal] (10.37,0.786) circle (1.05pt);
\fill[teal] (10.37,0.929) circle (1.05pt);
\fill[teal] (10.37,1.001) circle (1.05pt);
\fill[teal] (10.37,1.144) circle (1.05pt);
\fill[teal] (10.37,1.859) circle (1.05pt);
\fill[teal] (10.37,1.93) circle (1.05pt);
\fill[teal] (10.37,2.002) circle (1.05pt);
\fill[teal] (10.37,2.216) circle (1.05pt);
\fill[teal] (10.37,2.645) circle (1.05pt);
\fill[teal] (10.37,2.717) circle (1.05pt);
\fill[teal] (10.37,2.788) circle (1.05pt);
\fill[teal] (10.37,2.931) circle (1.05pt);
\fill[teal] (10.37,3.074) circle (1.05pt);
\fill[teal] (10.37,3.146) circle (1.05pt);
\fill[teal] (10.37,3.36) circle (1.05pt);
\fill[teal] (10.37,3.432) circle (1.05pt);
\fill[teal] (10.37,3.503) circle (1.05pt);
\fill[teal] (10.37,4.004) circle (1.05pt);
\fill[teal] (10.37,4.361) circle (1.05pt);
\fill[teal] (10.37,4.504) circle (1.05pt);
\fill[teal] (10.37,4.933) circle (1.05pt);
\fill[teal] (10.37,5.219) circle (1.05pt);
\fill[teal] (10.492,0.071) circle (1.05pt);
\fill[teal] (10.492,0.357) circle (1.05pt);
\fill[teal] (10.492,0.5) circle (1.05pt);
\fill[teal] (10.492,0.643) circle (1.05pt);
\fill[teal] (10.492,0.786) circle (1.05pt);
\fill[teal] (10.492,1.073) circle (1.05pt);
\fill[teal] (10.492,1.215) circle (1.05pt);
\fill[teal] (10.492,1.358) circle (1.05pt);
\fill[teal] (10.492,1.787) circle (1.05pt);
\fill[teal] (10.492,1.93) circle (1.05pt);
\fill[teal] (10.492,2.216) circle (1.05pt);
\fill[teal] (10.492,2.645) circle (1.05pt);
\fill[teal] (10.492,3.217) circle (1.05pt);
\fill[teal] (10.492,3.36) circle (1.05pt);
\fill[teal] (10.492,3.503) circle (1.05pt);
\fill[teal] (10.492,3.646) circle (1.05pt);
\fill[teal] (10.492,4.79) circle (1.05pt);
\fill[teal] (10.492,4.933) circle (1.05pt);
\fill[teal] (10.492,5.505) circle (1.05pt);
\fill[teal] (10.492,6.077) circle (1.05pt);
\fill[teal] (10.614,0.071) circle (1.05pt);
\fill[teal] (10.614,0.143) circle (1.05pt);
\fill[teal] (10.614,0.286) circle (1.05pt);
\fill[teal] (10.614,0.572) circle (1.05pt);
\fill[teal] (10.614,0.715) circle (1.05pt);
\fill[teal] (10.614,0.786) circle (1.05pt);
\fill[teal] (10.614,1.144) circle (1.05pt);
\fill[teal] (10.614,1.215) circle (1.05pt);
\fill[teal] (10.614,1.573) circle (1.05pt);
\fill[teal] (10.614,1.644) circle (1.05pt);
\fill[teal] (10.614,1.859) circle (1.05pt);
\fill[teal] (10.614,2.86) circle (1.05pt);
\fill[teal] (10.614,3.074) circle (1.05pt);
\fill[teal] (10.614,3.289) circle (1.05pt);
\fill[teal] (10.614,3.503) circle (1.05pt);
\fill[teal] (10.614,3.718) circle (1.05pt);
\fill[teal] (10.614,3.789) circle (1.05pt);
\fill[teal] (10.614,3.932) circle (1.05pt);
\fill[teal] (10.614,4.004) circle (1.05pt);
\fill[teal] (10.614,4.218) circle (1.05pt);
\fill[teal] (10.614,4.79) circle (1.05pt);
\fill[teal] (10.614,5.291) circle (1.05pt);
\fill[teal] (10.614,5.505) circle (1.05pt);
\fill[teal] (10.614,6.077) circle (1.05pt);
\fill[teal] (10.736,0.071) circle (1.05pt);
\fill[teal] (10.736,0.214) circle (1.05pt);
\fill[teal] (10.736,0.357) circle (1.05pt);
\fill[teal] (10.736,0.5) circle (1.05pt);
\fill[teal] (10.736,0.643) circle (1.05pt);
\fill[teal] (10.736,0.929) circle (1.05pt);
\fill[teal] (10.736,1.073) circle (1.05pt);
\fill[teal] (10.736,1.358) circle (1.05pt);
\fill[teal] (10.736,1.501) circle (1.05pt);
\fill[teal] (10.736,1.93) circle (1.05pt);
\fill[teal] (10.736,2.073) circle (1.05pt);
\fill[teal] (10.736,2.216) circle (1.05pt);
\fill[teal] (10.736,2.788) circle (1.05pt);
\fill[teal] (10.736,2.931) circle (1.05pt);
\fill[teal] (10.736,3.074) circle (1.05pt);
\fill[teal] (10.736,3.217) circle (1.05pt);
\fill[teal] (10.736,3.646) circle (1.05pt);
\fill[teal] (10.736,4.075) circle (1.05pt);
\fill[teal] (10.736,4.361) circle (1.05pt);
\fill[teal] (10.736,4.504) circle (1.05pt);
\fill[teal] (10.736,5.934) circle (1.05pt);
\fill[teal] (10.736,6.22) circle (1.05pt);
\fill[teal] (10.858,0.357) circle (1.05pt);
\fill[teal] (10.858,0.429) circle (1.05pt);
\fill[teal] (10.858,0.5) circle (1.05pt);
\fill[teal] (10.858,0.715) circle (1.05pt);
\fill[teal] (10.858,0.786) circle (1.05pt);
\fill[teal] (10.858,0.858) circle (1.05pt);
\fill[teal] (10.858,0.929) circle (1.05pt);
\fill[teal] (10.858,1.001) circle (1.05pt);
\fill[teal] (10.858,1.144) circle (1.05pt);
\fill[teal] (10.858,1.287) circle (1.05pt);
\fill[teal] (10.858,1.573) circle (1.05pt);
\fill[teal] (10.858,1.716) circle (1.05pt);
\fill[teal] (10.858,1.787) circle (1.05pt);
\fill[teal] (10.858,1.859) circle (1.05pt);
\fill[teal] (10.858,2.145) circle (1.05pt);
\fill[teal] (10.858,2.502) circle (1.05pt);
\fill[teal] (10.858,2.574) circle (1.05pt);
\fill[teal] (10.858,2.645) circle (1.05pt);
\fill[teal] (10.858,2.717) circle (1.05pt);
\fill[teal] (10.858,3.003) circle (1.05pt);
\fill[teal] (10.858,3.217) circle (1.05pt);
\fill[teal] (10.858,3.289) circle (1.05pt);
\fill[teal] (10.858,3.36) circle (1.05pt);
\fill[teal] (10.858,3.503) circle (1.05pt);
\fill[teal] (10.858,3.861) circle (1.05pt);
\fill[teal] (10.858,4.29) circle (1.05pt);
\fill[teal] (10.858,4.433) circle (1.05pt);
\fill[teal] (10.858,4.504) circle (1.05pt);
\fill[teal] (10.858,4.576) circle (1.05pt);
\fill[teal] (10.858,4.79) circle (1.05pt);
\fill[teal] (10.858,4.933) circle (1.05pt);
\fill[teal] (10.858,5.362) circle (1.05pt);
\fill[teal] (10.858,5.791) circle (1.05pt);
\fill[teal] (10.858,6.292) circle (1.05pt);
\fill[teal] (10.98,0.5) circle (1.05pt);
\fill[teal] (10.98,0.929) circle (1.05pt);
\fill[teal] (10.98,1.215) circle (1.05pt);
\fill[teal] (10.98,1.358) circle (1.05pt);
\fill[teal] (10.98,2.073) circle (1.05pt);
\fill[teal] (10.98,2.216) circle (1.05pt);
\fill[teal] (10.98,3.074) circle (1.05pt);
\fill[teal] (10.98,3.503) circle (1.05pt);
\fill[teal] (10.98,3.789) circle (1.05pt);
\fill[teal] (10.98,4.361) circle (1.05pt);
\fill[teal] (10.98,5.505) circle (1.05pt);
\fill[teal] (11.102,0.071) circle (1.05pt);
\fill[teal] (11.102,0.143) circle (1.05pt);
\fill[teal] (11.102,0.214) circle (1.05pt);
\fill[teal] (11.102,0.357) circle (1.05pt);
\fill[teal] (11.102,0.429) circle (1.05pt);
\fill[teal] (11.102,0.572) circle (1.05pt);
\fill[teal] (11.102,1.144) circle (1.05pt);
\fill[teal] (11.102,1.287) circle (1.05pt);
\fill[teal] (11.102,1.358) circle (1.05pt);
\fill[teal] (11.102,1.43) circle (1.05pt);
\fill[teal] (11.102,1.93) circle (1.05pt);
\fill[teal] (11.102,2.216) circle (1.05pt);
\fill[teal] (11.102,2.359) circle (1.05pt);
\fill[teal] (11.102,2.574) circle (1.05pt);
\fill[teal] (11.102,2.645) circle (1.05pt);
\fill[teal] (11.102,2.717) circle (1.05pt);
\fill[teal] (11.102,2.86) circle (1.05pt);
\fill[teal] (11.102,2.931) circle (1.05pt);
\fill[teal] (11.102,3.074) circle (1.05pt);
\fill[teal] (11.102,3.146) circle (1.05pt);
\fill[teal] (11.102,3.575) circle (1.05pt);
\fill[teal] (11.102,4.147) circle (1.05pt);
\fill[teal] (11.102,4.218) circle (1.05pt);
\fill[teal] (11.102,4.29) circle (1.05pt);
\fill[teal] (11.102,4.361) circle (1.05pt);
\fill[teal] (11.102,5.076) circle (1.05pt);
\fill[teal] (11.102,5.291) circle (1.05pt);
\fill[teal] (11.102,5.434) circle (1.05pt);
\fill[teal] (11.102,5.648) circle (1.05pt);
\fill[teal] (11.102,6.292) circle (1.05pt);
\fill[teal] (11.224,0.071) circle (1.05pt);
\fill[teal] (11.224,0.357) circle (1.05pt);
\fill[teal] (11.224,0.5) circle (1.05pt);
\fill[teal] (11.224,0.929) circle (1.05pt);
\fill[teal] (11.224,1.215) circle (1.05pt);
\fill[teal] (11.224,1.358) circle (1.05pt);
\fill[teal] (11.224,1.787) circle (1.05pt);
\fill[teal] (11.224,1.93) circle (1.05pt);
\fill[teal] (11.224,2.216) circle (1.05pt);
\fill[teal] (11.224,2.502) circle (1.05pt);
\fill[teal] (11.224,2.645) circle (1.05pt);
\fill[teal] (11.224,2.931) circle (1.05pt);
\fill[teal] (11.224,3.074) circle (1.05pt);
\fill[teal] (11.224,3.217) circle (1.05pt);
\fill[teal] (11.224,3.646) circle (1.05pt);
\fill[teal] (11.224,4.075) circle (1.05pt);
\fill[teal] (11.224,4.647) circle (1.05pt);
\fill[teal] (11.224,6.22) circle (1.05pt);
\fill[teal] (11.224,6.506) circle (1.05pt);
\fill[teal] (11.346,0.071) circle (1.05pt);
\fill[teal] (11.346,0.143) circle (1.05pt);
\fill[teal] (11.346,0.286) circle (1.05pt);
\fill[teal] (11.346,0.5) circle (1.05pt);
\fill[teal] (11.346,0.572) circle (1.05pt);
\fill[teal] (11.346,0.929) circle (1.05pt);
\fill[teal] (11.346,1.215) circle (1.05pt);
\fill[teal] (11.346,1.358) circle (1.05pt);
\fill[teal] (11.346,1.644) circle (1.05pt);
\fill[teal] (11.346,2.002) circle (1.05pt);
\fill[teal] (11.346,2.502) circle (1.05pt);
\fill[teal] (11.346,2.931) circle (1.05pt);
\fill[teal] (11.346,3.146) circle (1.05pt);
\fill[teal] (11.346,3.932) circle (1.05pt);
\fill[teal] (11.346,4.004) circle (1.05pt);
\fill[teal] (11.346,4.361) circle (1.05pt);
\fill[teal] (11.346,4.576) circle (1.05pt);
\fill[teal] (11.346,4.79) circle (1.05pt);
\fill[teal] (11.346,4.862) circle (1.05pt);
\fill[teal] (11.346,5.076) circle (1.05pt);
\fill[teal] (11.346,5.219) circle (1.05pt);
\fill[teal] (11.346,5.648) circle (1.05pt);
\fill[teal] (11.346,5.72) circle (1.05pt);
\fill[teal] (11.346,6.077) circle (1.05pt);
\fill[teal] (11.346,6.506) circle (1.05pt);
\fill[teal] (11.346,6.578) circle (1.05pt);
\fill[teal] (11.468,0.071) circle (1.05pt);
\fill[teal] (11.468,0.214) circle (1.05pt);
\fill[teal] (11.468,0.357) circle (1.05pt);
\fill[teal] (11.468,0.5) circle (1.05pt);
\fill[teal] (11.468,0.643) circle (1.05pt);
\fill[teal] (11.468,1.501) circle (1.05pt);
\fill[teal] (11.468,1.644) circle (1.05pt);
\fill[teal] (11.468,2.073) circle (1.05pt);
\fill[teal] (11.468,2.216) circle (1.05pt);
\fill[teal] (11.468,2.931) circle (1.05pt);
\fill[teal] (11.468,3.503) circle (1.05pt);
\fill[teal] (11.468,3.789) circle (1.05pt);
\fill[teal] (11.468,4.504) circle (1.05pt);
\fill[teal] (11.468,4.647) circle (1.05pt);
\fill[teal] (11.468,4.79) circle (1.05pt);
\fill[teal] (11.468,5.505) circle (1.05pt);
\fill[teal] (11.468,5.791) circle (1.05pt);
\fill[teal] (11.468,6.649) circle (1.05pt);
\fill[teal] (11.59,0.071) circle (1.05pt);
\fill[teal] (11.59,0.143) circle (1.05pt);
\fill[teal] (11.59,0.286) circle (1.05pt);
\fill[teal] (11.59,0.572) circle (1.05pt);
\fill[teal] (11.59,0.643) circle (1.05pt);
\fill[teal] (11.59,0.786) circle (1.05pt);
\fill[teal] (11.59,0.858) circle (1.05pt);
\fill[teal] (11.59,1.144) circle (1.05pt);
\fill[teal] (11.59,1.215) circle (1.05pt);
\fill[teal] (11.59,1.287) circle (1.05pt);
\fill[teal] (11.59,1.716) circle (1.05pt);
\fill[teal] (11.59,2.002) circle (1.05pt);
\fill[teal] (11.59,2.073) circle (1.05pt);
\fill[teal] (11.59,2.574) circle (1.05pt);
\fill[teal] (11.59,2.645) circle (1.05pt);
\fill[teal] (11.59,3.003) circle (1.05pt);
\fill[teal] (11.59,3.146) circle (1.05pt);
\fill[teal] (11.59,3.289) circle (1.05pt);
\fill[teal] (11.59,3.36) circle (1.05pt);
\fill[teal] (11.59,3.718) circle (1.05pt);
\fill[teal] (11.59,3.789) circle (1.05pt);
\fill[teal] (11.59,3.861) circle (1.05pt);
\fill[teal] (11.59,4.147) circle (1.05pt);
\fill[teal] (11.59,4.218) circle (1.05pt);
\fill[teal] (11.59,4.361) circle (1.05pt);
\fill[teal] (11.59,4.433) circle (1.05pt);
\fill[teal] (11.59,4.719) circle (1.05pt);
\fill[teal] (11.59,4.862) circle (1.05pt);
\fill[teal] (11.59,4.933) circle (1.05pt);
\fill[teal] (11.59,5.076) circle (1.05pt);
\fill[teal] (11.59,5.791) circle (1.05pt);
\fill[teal] (11.59,5.934) circle (1.05pt);
\fill[teal] (11.59,6.721) circle (1.05pt);
\fill[teal] (11.712,0.071) circle (1.05pt);
\fill[teal] (11.712,0.357) circle (1.05pt);
\fill[teal] (11.712,0.5) circle (1.05pt);
\fill[teal] (11.712,0.929) circle (1.05pt);
\fill[teal] (11.712,1.358) circle (1.05pt);
\fill[teal] (11.712,1.644) circle (1.05pt);
\fill[teal] (11.712,1.787) circle (1.05pt);
\fill[teal] (11.712,2.216) circle (1.05pt);
\fill[teal] (11.712,2.502) circle (1.05pt);
\fill[teal] (11.712,2.645) circle (1.05pt);
\fill[teal] (11.712,2.931) circle (1.05pt);
\fill[teal] (11.712,3.074) circle (1.05pt);
\fill[teal] (11.712,3.36) circle (1.05pt);
\fill[teal] (11.712,3.789) circle (1.05pt);
\fill[teal] (11.712,4.79) circle (1.05pt);
\fill[teal] (11.834,0.071) circle (1.05pt);
\fill[teal] (11.834,0.143) circle (1.05pt);
\fill[teal] (11.834,0.286) circle (1.05pt);
\fill[teal] (11.834,0.357) circle (1.05pt);
\fill[teal] (11.834,0.429) circle (1.05pt);
\fill[teal] (11.834,0.5) circle (1.05pt);
\fill[teal] (11.834,0.572) circle (1.05pt);
\fill[teal] (11.834,0.715) circle (1.05pt);
\fill[teal] (11.834,0.786) circle (1.05pt);
\fill[teal] (11.834,0.929) circle (1.05pt);
\fill[teal] (11.834,1.001) circle (1.05pt);
\fill[teal] (11.834,1.073) circle (1.05pt);
\fill[teal] (11.834,1.144) circle (1.05pt);
\fill[teal] (11.834,1.215) circle (1.05pt);
\fill[teal] (11.834,1.287) circle (1.05pt);
\fill[teal] (11.834,1.358) circle (1.05pt);
\fill[teal] (11.834,1.43) circle (1.05pt);
\fill[teal] (11.834,1.501) circle (1.05pt);
\fill[teal] (11.834,1.573) circle (1.05pt);
\fill[teal] (11.834,1.644) circle (1.05pt);
\fill[teal] (11.834,1.716) circle (1.05pt);
\fill[teal] (11.834,1.859) circle (1.05pt);
\fill[teal] (11.834,2.002) circle (1.05pt);
\fill[teal] (11.834,2.073) circle (1.05pt);
\fill[teal] (11.834,2.216) circle (1.05pt);
\fill[teal] (11.834,2.359) circle (1.05pt);
\fill[teal] (11.834,2.431) circle (1.05pt);
\fill[teal] (11.834,2.574) circle (1.05pt);
\fill[teal] (11.834,2.645) circle (1.05pt);
\fill[teal] (11.834,2.931) circle (1.05pt);
\fill[teal] (11.834,3.217) circle (1.05pt);
\fill[teal] (11.834,3.36) circle (1.05pt);
\fill[teal] (11.834,3.432) circle (1.05pt);
\fill[teal] (11.834,3.575) circle (1.05pt);
\fill[teal] (11.834,3.861) circle (1.05pt);
\fill[teal] (11.834,4.719) circle (1.05pt);
\fill[teal] (11.834,5.219) circle (1.05pt);
\fill[teal] (11.834,5.362) circle (1.05pt);
\fill[teal] (11.834,5.434) circle (1.05pt);
\fill[teal] (11.834,5.505) circle (1.05pt);
\fill[teal] (11.834,5.934) circle (1.05pt);
\fill[teal] (11.834,6.506) circle (1.05pt);
\fill[teal] (11.834,6.649) circle (1.05pt);
\fill[teal] (11.834,6.792) circle (1.05pt);
\fill[teal] (11.956,0.071) circle (1.05pt);
\fill[teal] (11.956,0.214) circle (1.05pt);
\fill[teal] (11.956,1.073) circle (1.05pt);
\fill[teal] (11.956,1.215) circle (1.05pt);
\fill[teal] (11.956,2.073) circle (1.05pt);
\fill[teal] (11.956,2.216) circle (1.05pt);
\fill[teal] (11.956,2.788) circle (1.05pt);
\fill[teal] (11.956,3.074) circle (1.05pt);
\fill[teal] (11.956,3.36) circle (1.05pt);
\fill[teal] (11.956,3.932) circle (1.05pt);
\fill[teal] (11.956,4.218) circle (1.05pt);
\fill[teal] (11.956,4.361) circle (1.05pt);
\fill[teal] (11.956,4.647) circle (1.05pt);
\fill[teal] (11.956,6.935) circle (1.05pt);
\fill[teal] (12.078,1.215) circle (1.05pt);
\fill[teal] (12.078,1.358) circle (1.05pt);
\fill[teal] (12.078,2.216) circle (1.05pt);
\fill[teal] (12.078,2.288) circle (1.05pt);
\fill[teal] (12.078,2.431) circle (1.05pt);
\fill[teal] (12.078,2.502) circle (1.05pt);
\fill[teal] (12.078,2.717) circle (1.05pt);
\fill[teal] (12.078,2.931) circle (1.05pt);
\fill[teal] (12.078,3.289) circle (1.05pt);
\fill[teal] (12.078,3.36) circle (1.05pt);
\fill[teal] (12.078,3.503) circle (1.05pt);
\fill[teal] (12.078,3.789) circle (1.05pt);
\fill[teal] (12.078,4.004) circle (1.05pt);
\fill[teal] (12.078,5.291) circle (1.05pt);
\fill[teal] (12.078,6.792) circle (1.05pt);
\fill[teal] (12.2,0.071) circle (1.05pt);
\fill[teal] (12.2,0.214) circle (1.05pt);
\fill[teal] (12.2,0.5) circle (1.05pt);
\fill[teal] (12.2,0.929) circle (1.05pt);
\fill[teal] (12.2,1.358) circle (1.05pt);
\fill[teal] (12.2,1.501) circle (1.05pt);
\fill[teal] (12.2,2.073) circle (1.05pt);
\fill[teal] (12.2,2.788) circle (1.05pt);
\fill[teal] (12.2,3.503) circle (1.05pt);
\fill[teal] (12.2,4.361) circle (1.05pt);
\fill[teal] (12.2,5.076) circle (1.05pt);
\fill[ochre] (0.854,0.143) ++(-1.15pt,-1.15pt) rectangle ++(2.3pt,2.3pt);
\fill[ochre] (1.098,0.143) ++(-1.15pt,-1.15pt) rectangle ++(2.3pt,2.3pt);
\fill[ochre] (1.098,0.5) ++(-1.15pt,-1.15pt) rectangle ++(2.3pt,2.3pt);
\fill[ochre] (1.342,0.429) ++(-1.15pt,-1.15pt) rectangle ++(2.3pt,2.3pt);
\fill[ochre] (1.342,0.572) ++(-1.15pt,-1.15pt) rectangle ++(2.3pt,2.3pt);
\fill[ochre] (1.342,0.715) ++(-1.15pt,-1.15pt) rectangle ++(2.3pt,2.3pt);
\fill[ochre] (1.464,0.786) ++(-1.15pt,-1.15pt) rectangle ++(2.3pt,2.3pt);
\fill[ochre] (1.586,0.071) ++(-1.15pt,-1.15pt) rectangle ++(2.3pt,2.3pt);
\fill[ochre] (1.586,0.214) ++(-1.15pt,-1.15pt) rectangle ++(2.3pt,2.3pt);
\fill[ochre] (1.586,0.357) ++(-1.15pt,-1.15pt) rectangle ++(2.3pt,2.3pt);
\fill[ochre] (1.586,0.786) ++(-1.15pt,-1.15pt) rectangle ++(2.3pt,2.3pt);
\fill[ochre] (1.586,0.858) ++(-1.15pt,-1.15pt) rectangle ++(2.3pt,2.3pt);
\fill[ochre] (1.708,0.214) ++(-1.15pt,-1.15pt) rectangle ++(2.3pt,2.3pt);
\fill[ochre] (1.708,0.643) ++(-1.15pt,-1.15pt) rectangle ++(2.3pt,2.3pt);
\fill[ochre] (1.708,0.929) ++(-1.15pt,-1.15pt) rectangle ++(2.3pt,2.3pt);
\fill[ochre] (1.83,0.071) ++(-1.15pt,-1.15pt) rectangle ++(2.3pt,2.3pt);
\fill[ochre] (1.83,0.286) ++(-1.15pt,-1.15pt) rectangle ++(2.3pt,2.3pt);
\fill[ochre] (1.83,0.786) ++(-1.15pt,-1.15pt) rectangle ++(2.3pt,2.3pt);
\fill[ochre] (1.952,0.071) ++(-1.15pt,-1.15pt) rectangle ++(2.3pt,2.3pt);
\fill[ochre] (1.952,0.214) ++(-1.15pt,-1.15pt) rectangle ++(2.3pt,2.3pt);
\fill[ochre] (1.952,0.357) ++(-1.15pt,-1.15pt) rectangle ++(2.3pt,2.3pt);
\fill[ochre] (1.952,0.5) ++(-1.15pt,-1.15pt) rectangle ++(2.3pt,2.3pt);
\fill[ochre] (1.952,0.929) ++(-1.15pt,-1.15pt) rectangle ++(2.3pt,2.3pt);
\fill[ochre] (2.074,0.071) ++(-1.15pt,-1.15pt) rectangle ++(2.3pt,2.3pt);
\fill[ochre] (2.074,0.214) ++(-1.15pt,-1.15pt) rectangle ++(2.3pt,2.3pt);
\fill[ochre] (2.074,0.429) ++(-1.15pt,-1.15pt) rectangle ++(2.3pt,2.3pt);
\fill[ochre] (2.074,1.073) ++(-1.15pt,-1.15pt) rectangle ++(2.3pt,2.3pt);
\fill[ochre] (2.074,1.144) ++(-1.15pt,-1.15pt) rectangle ++(2.3pt,2.3pt);
\fill[ochre] (2.196,0.071) ++(-1.15pt,-1.15pt) rectangle ++(2.3pt,2.3pt);
\fill[ochre] (2.196,0.357) ++(-1.15pt,-1.15pt) rectangle ++(2.3pt,2.3pt);
\fill[ochre] (2.196,1.215) ++(-1.15pt,-1.15pt) rectangle ++(2.3pt,2.3pt);
\fill[ochre] (2.318,0.071) ++(-1.15pt,-1.15pt) rectangle ++(2.3pt,2.3pt);
\fill[ochre] (2.318,0.143) ++(-1.15pt,-1.15pt) rectangle ++(2.3pt,2.3pt);
\fill[ochre] (2.318,0.214) ++(-1.15pt,-1.15pt) rectangle ++(2.3pt,2.3pt);
\fill[ochre] (2.318,0.429) ++(-1.15pt,-1.15pt) rectangle ++(2.3pt,2.3pt);
\fill[ochre] (2.318,0.572) ++(-1.15pt,-1.15pt) rectangle ++(2.3pt,2.3pt);
\fill[ochre] (2.44,0.5) ++(-1.15pt,-1.15pt) rectangle ++(2.3pt,2.3pt);
\fill[ochre] (2.562,0.357) ++(-1.15pt,-1.15pt) rectangle ++(2.3pt,2.3pt);
\fill[ochre] (2.562,0.572) ++(-1.15pt,-1.15pt) rectangle ++(2.3pt,2.3pt);
\fill[ochre] (2.562,0.715) ++(-1.15pt,-1.15pt) rectangle ++(2.3pt,2.3pt);
\fill[ochre] (2.562,0.786) ++(-1.15pt,-1.15pt) rectangle ++(2.3pt,2.3pt);
\fill[ochre] (2.562,0.929) ++(-1.15pt,-1.15pt) rectangle ++(2.3pt,2.3pt);
\fill[ochre] (2.684,0.071) ++(-1.15pt,-1.15pt) rectangle ++(2.3pt,2.3pt);
\fill[ochre] (2.684,0.214) ++(-1.15pt,-1.15pt) rectangle ++(2.3pt,2.3pt);
\fill[ochre] (2.684,0.5) ++(-1.15pt,-1.15pt) rectangle ++(2.3pt,2.3pt);
\fill[ochre] (2.684,0.643) ++(-1.15pt,-1.15pt) rectangle ++(2.3pt,2.3pt);
\fill[ochre] (2.684,1.073) ++(-1.15pt,-1.15pt) rectangle ++(2.3pt,2.3pt);
\fill[ochre] (2.806,0.214) ++(-1.15pt,-1.15pt) rectangle ++(2.3pt,2.3pt);
\fill[ochre] (2.806,1.215) ++(-1.15pt,-1.15pt) rectangle ++(2.3pt,2.3pt);
\fill[ochre] (2.806,1.573) ++(-1.15pt,-1.15pt) rectangle ++(2.3pt,2.3pt);
\fill[ochre] (2.928,0.071) ++(-1.15pt,-1.15pt) rectangle ++(2.3pt,2.3pt);
\fill[ochre] (3.05,0.071) ++(-1.15pt,-1.15pt) rectangle ++(2.3pt,2.3pt);
\fill[ochre] (3.05,0.214) ++(-1.15pt,-1.15pt) rectangle ++(2.3pt,2.3pt);
\fill[ochre] (3.05,0.286) ++(-1.15pt,-1.15pt) rectangle ++(2.3pt,2.3pt);
\fill[ochre] (3.05,0.5) ++(-1.15pt,-1.15pt) rectangle ++(2.3pt,2.3pt);
\fill[ochre] (3.05,1.501) ++(-1.15pt,-1.15pt) rectangle ++(2.3pt,2.3pt);
\fill[ochre] (3.172,0.214) ++(-1.15pt,-1.15pt) rectangle ++(2.3pt,2.3pt);
\fill[ochre] (3.294,0.071) ++(-1.15pt,-1.15pt) rectangle ++(2.3pt,2.3pt);
\fill[ochre] (3.294,0.929) ++(-1.15pt,-1.15pt) rectangle ++(2.3pt,2.3pt);
\fill[ochre] (3.294,1.358) ++(-1.15pt,-1.15pt) rectangle ++(2.3pt,2.3pt);
\fill[ochre] (3.294,1.787) ++(-1.15pt,-1.15pt) rectangle ++(2.3pt,2.3pt);
\fill[ochre] (3.416,0.071) ++(-1.15pt,-1.15pt) rectangle ++(2.3pt,2.3pt);
\fill[ochre] (3.538,0.357) ++(-1.15pt,-1.15pt) rectangle ++(2.3pt,2.3pt);
\fill[ochre] (3.782,0.071) ++(-1.15pt,-1.15pt) rectangle ++(2.3pt,2.3pt);
\fill[ochre] (3.782,2.073) ++(-1.15pt,-1.15pt) rectangle ++(2.3pt,2.3pt);
\fill[ochre] (3.904,0.357) ++(-1.15pt,-1.15pt) rectangle ++(2.3pt,2.3pt);
\fill[ochre] (3.904,1.215) ++(-1.15pt,-1.15pt) rectangle ++(2.3pt,2.3pt);
\fill[ochre] (3.904,1.644) ++(-1.15pt,-1.15pt) rectangle ++(2.3pt,2.3pt);
\fill[ochre] (4.148,0.5) ++(-1.15pt,-1.15pt) rectangle ++(2.3pt,2.3pt);
\fill[ochre] (4.392,0.929) ++(-1.15pt,-1.15pt) rectangle ++(2.3pt,2.3pt);
\fill[ochre] (4.514,0.071) ++(-1.15pt,-1.15pt) rectangle ++(2.3pt,2.3pt);
\fill[ochre] (4.514,0.429) ++(-1.15pt,-1.15pt) rectangle ++(2.3pt,2.3pt);
\fill[ochre] (4.514,0.5) ++(-1.15pt,-1.15pt) rectangle ++(2.3pt,2.3pt);
\fill[ochre] (4.636,0.071) ++(-1.15pt,-1.15pt) rectangle ++(2.3pt,2.3pt);
\fill[ochre] (4.636,1.93) ++(-1.15pt,-1.15pt) rectangle ++(2.3pt,2.3pt);
\fill[ochre] (5.246,0.5) ++(-1.15pt,-1.15pt) rectangle ++(2.3pt,2.3pt);
\fill[ochre] (5.246,1.859) ++(-1.15pt,-1.15pt) rectangle ++(2.3pt,2.3pt);
\fill[ochre] (5.734,0.572) ++(-1.15pt,-1.15pt) rectangle ++(2.3pt,2.3pt);
\fill[ochre] (5.734,2.002) ++(-1.15pt,-1.15pt) rectangle ++(2.3pt,2.3pt);
\fill[ochre] (6.71,0.071) ++(-1.15pt,-1.15pt) rectangle ++(2.3pt,2.3pt);
\fill[ochre] (7.076,2.502) ++(-1.15pt,-1.15pt) rectangle ++(2.3pt,2.3pt);
\fill[ochre] (7.442,0.715) ++(-1.15pt,-1.15pt) rectangle ++(2.3pt,2.3pt);
\fill[ochre] (7.564,2.645) ++(-1.15pt,-1.15pt) rectangle ++(2.3pt,2.3pt);
\fill[ochre] (7.93,0.786) ++(-1.15pt,-1.15pt) rectangle ++(2.3pt,2.3pt);
\fill[ochre] (8.906,3.146) ++(-1.15pt,-1.15pt) rectangle ++(2.3pt,2.3pt);
\fill[ochre] (9.394,3.289) ++(-1.15pt,-1.15pt) rectangle ++(2.3pt,2.3pt);
\fill[ochre] (9.638,0.929) ++(-1.15pt,-1.15pt) rectangle ++(2.3pt,2.3pt);
\fill[ochre] (10.126,1.001) ++(-1.15pt,-1.15pt) rectangle ++(2.3pt,2.3pt);
\fill[ochre] (10.736,3.789) ++(-1.15pt,-1.15pt) rectangle ++(2.3pt,2.3pt);
\fill[ochre] (11.224,3.932) ++(-1.15pt,-1.15pt) rectangle ++(2.3pt,2.3pt);
\fill[wine] (0.61,0.286) ++(0,1.9pt) -- ++(1.9pt,-1.9pt) -- ++(-1.9pt,-1.9pt) -- ++(-1.9pt,1.9pt) -- cycle;
\fill[wine] (0.854,0.071) ++(0,1.9pt) -- ++(1.9pt,-1.9pt) -- ++(-1.9pt,-1.9pt) -- ++(-1.9pt,1.9pt) -- cycle;
\fill[wine] (0.854,0.286) ++(0,1.9pt) -- ++(1.9pt,-1.9pt) -- ++(-1.9pt,-1.9pt) -- ++(-1.9pt,1.9pt) -- cycle;
\fill[wine] (0.854,0.357) ++(0,1.9pt) -- ++(1.9pt,-1.9pt) -- ++(-1.9pt,-1.9pt) -- ++(-1.9pt,1.9pt) -- cycle;
\fill[wine] (0.976,0.5) ++(0,1.9pt) -- ++(1.9pt,-1.9pt) -- ++(-1.9pt,-1.9pt) -- ++(-1.9pt,1.9pt) -- cycle;
\fill[wine] (1.098,0.071) ++(0,1.9pt) -- ++(1.9pt,-1.9pt) -- ++(-1.9pt,-1.9pt) -- ++(-1.9pt,1.9pt) -- cycle;
\fill[wine] (1.098,0.572) ++(0,1.9pt) -- ++(1.9pt,-1.9pt) -- ++(-1.9pt,-1.9pt) -- ++(-1.9pt,1.9pt) -- cycle;
\fill[wine] (1.22,0.071) ++(0,1.9pt) -- ++(1.9pt,-1.9pt) -- ++(-1.9pt,-1.9pt) -- ++(-1.9pt,1.9pt) -- cycle;
\fill[wine] (1.22,0.214) ++(0,1.9pt) -- ++(1.9pt,-1.9pt) -- ++(-1.9pt,-1.9pt) -- ++(-1.9pt,1.9pt) -- cycle;
\fill[wine] (1.342,0.143) ++(0,1.9pt) -- ++(1.9pt,-1.9pt) -- ++(-1.9pt,-1.9pt) -- ++(-1.9pt,1.9pt) -- cycle;
\fill[wine] (1.464,0.071) ++(0,1.9pt) -- ++(1.9pt,-1.9pt) -- ++(-1.9pt,-1.9pt) -- ++(-1.9pt,1.9pt) -- cycle;
\fill[wine] (1.464,0.5) ++(0,1.9pt) -- ++(1.9pt,-1.9pt) -- ++(-1.9pt,-1.9pt) -- ++(-1.9pt,1.9pt) -- cycle;
\fill[wine] (1.586,0.572) ++(0,1.9pt) -- ++(1.9pt,-1.9pt) -- ++(-1.9pt,-1.9pt) -- ++(-1.9pt,1.9pt) -- cycle;
\fill[wine] (2.074,0.143) ++(0,1.9pt) -- ++(1.9pt,-1.9pt) -- ++(-1.9pt,-1.9pt) -- ++(-1.9pt,1.9pt) -- cycle;
\fill[wine] (2.074,0.715) ++(0,1.9pt) -- ++(1.9pt,-1.9pt) -- ++(-1.9pt,-1.9pt) -- ++(-1.9pt,1.9pt) -- cycle;
\fill[wine] (3.416,1.215) ++(0,1.9pt) -- ++(1.9pt,-1.9pt) -- ++(-1.9pt,-1.9pt) -- ++(-1.9pt,1.9pt) -- cycle;
\fill[wine] (3.904,1.358) ++(0,1.9pt) -- ++(1.9pt,-1.9pt) -- ++(-1.9pt,-1.9pt) -- ++(-1.9pt,1.9pt) -- cycle;
\draw[ink, line width=0.5pt] (2.074,0.715) circle (3.2pt);
\draw[ink, line width=0.5pt] (3.904,1.358) circle (3.2pt);
\draw[ink, line width=0.5pt] (5.734,2.002) circle (3.2pt);
\draw[ink, line width=0.5pt] (7.564,2.645) circle (3.2pt);
\draw[ink, line width=0.5pt] (9.394,3.289) circle (3.2pt);
\draw[ink, line width=0.5pt] (11.224,3.932) circle (3.2pt);
\draw[ink, line width=0.5pt] (0.122,0) -- (12.322,0);
\draw[ink, line width=0.5pt] (0.122,0) -- (0.122,7.15);
\draw[ink, line width=0.5pt] (1.22,0) -- (1.22,-0.08);
\node[ink, below] at (1.22,-0.08) {$10$};
\draw[ink, line width=0.5pt] (2.44,0) -- (2.44,-0.08);
\node[ink, below] at (2.44,-0.08) {$20$};
\draw[ink, line width=0.5pt] (3.66,0) -- (3.66,-0.08);
\node[ink, below] at (3.66,-0.08) {$30$};
\draw[ink, line width=0.5pt] (4.88,0) -- (4.88,-0.08);
\node[ink, below] at (4.88,-0.08) {$40$};
\draw[ink, line width=0.5pt] (6.1,0) -- (6.1,-0.08);
\node[ink, below] at (6.1,-0.08) {$50$};
\draw[ink, line width=0.5pt] (7.32,0) -- (7.32,-0.08);
\node[ink, below] at (7.32,-0.08) {$60$};
\draw[ink, line width=0.5pt] (8.54,0) -- (8.54,-0.08);
\node[ink, below] at (8.54,-0.08) {$70$};
\draw[ink, line width=0.5pt] (9.76,0) -- (9.76,-0.08);
\node[ink, below] at (9.76,-0.08) {$80$};
\draw[ink, line width=0.5pt] (10.98,0) -- (10.98,-0.08);
\node[ink, below] at (10.98,-0.08) {$90$};
\draw[ink, line width=0.5pt] (12.2,0) -- (12.2,-0.08);
\node[ink, below] at (12.2,-0.08) {$100$};
\draw[ink, line width=0.5pt] (0.122,0.071) -- (0.042,0.071);
\node[ink, left] at (0.042,0.071) {$1$};
\draw[ink, line width=0.5pt] (0.122,0.929) -- (0.042,0.929);
\node[ink, left] at (0.042,0.929) {$13$};
\draw[ink, line width=0.5pt] (0.122,1.787) -- (0.042,1.787);
\node[ink, left] at (0.042,1.787) {$25$};
\draw[ink, line width=0.5pt] (0.122,3.575) -- (0.042,3.575);
\node[ink, left] at (0.042,3.575) {$50$};
\draw[ink, line width=0.5pt] (0.122,5.362) -- (0.042,5.362);
\node[ink, left] at (0.042,5.362) {$75$};
\draw[ink, line width=0.5pt] (0.122,7.078) -- (0.042,7.078);
\node[ink, left] at (0.042,7.078) {$99$};
\node[ink] at (6.161,-0.75) {denominator $Q$};
\node[ink, rotate=90] at (-0.728,3.575) {numerator $P$};
\node[ink, right] at (0.582,6.998) {margin};
\fill[wine] (0.732,6.578) ++(0,2.4699999999999998pt) -- ++(2.4699999999999998pt,-2.4699999999999998pt) -- ++(-2.4699999999999998pt,-2.4699999999999998pt) -- ++(-2.4699999999999998pt,2.4699999999999998pt) -- cycle;
\node[ink, right] at (0.832,6.578) {1 or 2};
\fill[ochre] (0.732,6.218) ++(-1.4949999999999999pt,-1.4949999999999999pt) rectangle ++(2.9899999999999998pt,2.9899999999999998pt);
\node[ink, right] at (0.832,6.218) {3 to 9};
\fill[teal] (0.732,5.858) circle (1.3650000000000002pt);
\node[ink, right] at (0.832,5.858) {10 to 99};
\fill[navy] (0.732,5.498) circle (0.9750000000000001pt);
\node[ink, right] at (0.832,5.498) {100 or more};
\draw[slate, line width=0.5pt] (0.732,5.138) circle (1.3pt);
\node[ink, right] at (0.832,5.138) {one representation};
\draw[ink, line width=0.5pt] (0.732,4.778) circle (3.2pt);
\node[ink, right] at (0.832,4.778) {$(9k+1)/(15k+2)$};
\end{tikzpicture}

\caption{Computed evidence for \cref{conj:euclidean-minimality}. Each mark is
an occurrence $P/Q$ with $Q\leqslant100$, placed at its denominator and
numerator; its color and shape give the margin, the cost of the cheapest
other primitive representation of the label minus the occurrence cost
$P+Q-2$. Hollow circles are labels with a single representation, which have
no competitor. No margin is zero or negative. The small margins thin out
as the denominator grows; the ringed occurrences $(9k+1)/(15k+2)$,
$k=1,\dots,6$, are the last ones with a margin at most $k$, and
\cref{prop:two-ladders} explains them. The shaded rows $P\leqslant13$ are those
of \cref{thm:thirteen-rows}.}
\label{fig:margins}
\end{figure}

\subsection{How small the margin can be}

The rarefaction of small margins in \cref{fig:margins} follows a precise
law in the computed range. Among the occurrences with $P+Q\leqslant160$, a
margin at most $k$, for $k=1,\dots,6$, occurs only when
$$
 P+Q\leqslant24k+3,\qquad P\leqslant9k+1,\qquad Q\leqslant15k+2 ,
$$
and all three bounds are reached by the single occurrence
$(9k+1)/(15k+2)$, whose margin is exactly $k$ (\cref{fig:thresholds}). These
six fractions are rungs of one Farey ladder, the fractions
$(1+3j)/(2+5j)$ that run from $1/2$ to $3/5$, taken at $j=3k$; their labels
are the neighbors of $m_{3/5}=433$ in the Markov tree. On every third rung
the label is divisible by $5=(1+2i)(1-2i)$, and at the six computed rungs
the cheapest competitor of the occurrence pair is the element of the fiber
in which the Gaussian factor above $5$ is replaced by its conjugate. A companion ladder,
$(2+3j)/(3+5j)$ from $2/3$ to $3/5$, behaves in the same way at $j=3k-1$. The
mechanism can be proved for every rung, and it shows that thresholds of
this kind, if they exist, cannot be smaller.

The proof follows the occurrence vectors $v_j=H_je$ of the rungs, $H_j$ being
the Gram factor of \cref{prop:occurrence-cost}. Two rungs apart, the central
words grow by one fixed word at each end, so that $v_{j+2}$ is obtained from
$v_j$ by a matrix that depends only on the parity of $j$; six rungs apart,
the competitor obtained by conjugating the factor above $5$ is transformed
by the same matrix conjugated by the rotation that multiplies by
$(1-2i)/(1+2i)$.
Cost is additive under nonnegative matrices of determinant one, and the
whole statement reduces to the lengths of four explicit matrices.

\begin{lemma}[Finite facts on two ladders]\label{lem:ladder-facts}
For the ladder $(1+3j)/(2+5j)$ put $u=babab$ and $j_0=3$, and for the ladder
$(2+3j)/(3+5j)$ put $u=abbab$ and $j_0=2$. For each ladder and each parity
of $j\geqslant1$, let $T=H_j\,\mu(u)\,H_j^{-1}$, computed at the rungs $j=1$ and
$j=2$. At each of the four rungs $j_0$ and $j_0+3$ of the two ladders, exactly
one of the two matrices $R$ of multiplication by $(1-2i)/(1+2i)$ and by
$(1+2i)/(1-2i)$ maps $v_j$ to an integer vector, and exactly one diagonal sign
matrix $S$ makes that vector positive. Then:
\begin{enumerate}[label=\textup{(\roman*)},ref=\textup{(\roman*)},leftmargin=*]
\item\label{it:ladder-transfer} On the first ladder $T$ is
 $\left(\begin{smallmatrix}1105&463\\463&194\end{smallmatrix}\right)$ at even $j$
 and $\left(\begin{smallmatrix}1130&437\\437&169\end{smallmatrix}\right)$ at odd
 $j$; on the second ladder the two matrices are exchanged.
\item\label{it:ladder-length} At each of the four rungs, with $T$ the matrix of its parity, the
 matrix $N=SRT^3R^{-1}S^{-1}$ is a nonnegative integer matrix of determinant
 one whose expression as a word in $L_a$ and $L_b$ has $50$ letters, while
 $T^3$ has $48$.
\item\label{it:ladder-base} At each of the four rungs the vector $SRv_j$ is primitive, and its
 Euclidean cost exceeds that of $v_j$ by $1$ at $j_0$ and by $2$ at $j_0+3$.
\end{enumerate}
\end{lemma}

\begin{proof}
These are finitely many assertions about explicit integer matrices and
vectors. They were verified by computer in exact arithmetic, and the
verification was carried out twice independently: once through the Gram
factors and the matrices of \ref{it:ladder-length}, and once by Gaussian division of the
occurrence pairs, with the labels taken as continuants of the continued
fraction $[2,\sigma(z),2]$.
\end{proof}

\begin{proposition}[Margins along two ladders]\label{prop:two-ladders}
For every integer $k\geqslant1$, the labels of the reduced fractions
$(9k+1)/(15k+2)$ and $(9k-1)/(15k-2)$ are divisible by $5$. In the occurrence
pair of each, written as a Gaussian integer, replacing the factor above $5$
by its conjugate gives a primitive representation of the same label whose
Euclidean cost exceeds the occurrence cost by exactly $k$. Consequently the
margins of these two occurrences are at most $k$.
\end{proposition}

\begin{proof}
\emph{The words.} Two fractions $a/b<c/d$ with $bc-ad=1$ are Farey
neighbors, and the lower Christoffel word of their mediant is the
product $\gamma_{a/b}\gamma_{c/d}$ \cite{Reutenauer2009}. The fraction
$(1+3j)/(2+5j)$ is a neighbor of $3/5$ lying below it, since
$3(2+5j)-5(1+3j)=1$, and its mediant with $3/5$ is the next rung; starting
from $\gamma_{1/2}=ab$ and $\gamma_{3/5}=ababb$, induction gives
$\gamma_{(1+3j)/(2+5j)}=ab\,(ababb)^j$, with central word
$z_j=b\,(ababb)^{j-1}abab$. Concatenation gives
$z_{j+2}=babab\cdot z_j\cdot babab=u\,z_j\,u$. In the same way
$5(2+3j)-3(3+5j)=1$ places $(2+3j)/(3+5j)$ above $3/5$, the words are
$(ababb)^jabb$ with central words $z_j=babb\,(ababb)^{j-1}ab$, and
$z_{j+2}=babba\cdot z_j\cdot abbab=\widetilde u\,z_j\,u$.

\emph{The transfer.} In both cases the right half of the central word gains
the word $u$ at its far end while its middle letter is unchanged, so that
the Gram factor of \cref{prop:occurrence-cost} satisfies
$H_{j+2}=H_j\,\mu(u)$. Put $T_j=H_j\,\mu(u)\,H_j^{-1}$. Then
$T_jv_j=H_j\,\mu(u)\,e=H_{j+2}e=v_{j+2}$, and
$T_{j+2}=H_j\,\mu(u)\,\mu(u)\,\mu(u)^{-1}H_j^{-1}=T_j$, so that $T_j$ depends only
on the parity of $j$: it is the matrix $T$ of \cref{lem:ladder-facts}\ref{it:ladder-transfer}, and
$v_{j+6}=T^3v_j$.

\emph{The competitor.} Write $v_j=(s,r)^{\mathsf T}$, so that the label is
$m=s^2+r^2$. Fix one of the four rungs $j_0$, $j_0+3$ of
\cref{lem:ladder-facts}, with its matrices $R$, $S$ and $T$, and put
$c_j=SR\,v_j$ for every rung $j$ congruent to it modulo six. Since
$v_{j+6}=T^3v_j$,
$$
 c_{j+6}=SR\,T^3v_j=N\,c_j,\qquad N=SRT^3R^{-1}S^{-1}.
$$
The matrix $N$ has nonnegative integer entries and determinant one
(\cref{lem:ladder-facts}\ref{it:ladder-length}); it maps positive vectors to positive
vectors, since no row of it vanishes, and primitive integer vectors to
primitive integer vectors, since its inverse is integral. By
\cref{lem:ladder-facts}\ref{it:ladder-base} and induction on $j$ in steps of six, $c_j$ is
therefore a positive primitive integer vector at every rung $j=3k$ of the
first ladder and $j=3k-1$ of the second. The matrices $R$ and $S$ preserve
lengths, so $c_j$ has norm $m$: up to the order of its entries it is an
element of $\mathcal F(m)$. As a Gaussian integer, $c_j$ is $s+ir$ multiplied
by $(1\mp2i)/(1\pm2i)$, up to a sign and a conjugation; its integrality means
that $1\pm2i$ divides $s+ir$, so that $5$ divides $m$, and its primitivity that
$25$ does not, by \cref{prop:gaussian-fibre}. The element $c_j$ is thus the
occurrence pair with the factor above $5$ replaced by its conjugate.

\emph{Cost.} By \cref{lem:euclidean-bijection}, every positive primitive
column is $x=L_w{\bf1}$ for a unique word $w$, so that for every word $w'$,
$\ell(L_{w'}x)=\ell(L_{w'w}{\bf1})=\lvert w'\rvert+\ell(x)$. Every nonnegative integer
matrix $M=\left(\begin{smallmatrix}\alpha&\beta\\\gamma&\delta\end{smallmatrix}\right)\neq I$
of determinant one is such an $L_{w'}$. Indeed one of its rows dominates the
other entrywise: if $\alpha>\gamma$ and $\beta<\delta$, then
$\alpha\delta-\beta\gamma\geqslant(\gamma+1)(\beta+1)-\beta\gamma=\beta+\gamma+1$, which equals one
only for $M=I$; and if $\alpha<\gamma$ and $\beta>\delta$, then
$\alpha\delta-\beta\gamma\leqslant(\gamma-1)(\beta-1)-\beta\gamma<0$. Subtracting the dominated
row from the dominating one removes a first letter $L_a$ or $L_b$, keeps the
entries nonnegative and the determinant one, and lowers the sum of the
entries, so that the process ends at $I$. By \cref{lem:ladder-facts}\ref{it:ladder-length},
$$
 \ell(c_{j+6})=\ell(c_j)+50,\qquad \ell(v_{j+6})=\ell(v_j)+48,
$$
so that the excess $\ell(c_j)-\ell(v_j)$ grows by $2$ every six rungs. It is
$1$ and $2$ at the first two rungs of each family
(\cref{lem:ladder-facts}\ref{it:ladder-base}), hence $k$ at the $k$th, which is the rung
$(9k+1)/(15k+2)$ of the first ladder and $(9k-1)/(15k-2)$ of the second.
The margin is at most the excess of this particular competitor.
\end{proof}

The proposition explains the six ringed occurrences of \cref{fig:margins}:
together with the computations, which find no cheaper competitor at those
rungs, it gives their margins exactly. It also shows that the thresholds
observed above cannot be lowered: for every $k$ there is an occurrence with
$P+Q=24k+3$ and margin at most $k$.

\begin{figure}[tb]
\centering
% Authors: Dr. Denys Dutykh (Mathematics Department, Khalifa University of Science and
% Technology, Abu Dhabi, UAE) and Prof. Laurent Vuillon (Univ. Savoie Mont Blanc, CNRS, LAMA)
% Generated by supplement/code/make_figures.py from exact data; do not edit by hand.
\begin{tikzpicture}[font=\footnotesize, >={Stealth[length=4pt]}]
\fill[sand] (0.225,0) -- (12,3.271) -- (12,0) -- cycle;
\draw[slate, line width=0.7pt, dash pattern=on 5pt off 2pt] (0.225,0) -- (12,3.271);
\filldraw[navy, fill=white, line width=0.6pt] (0.6,0.5) circle (1.5pt);
\filldraw[navy, fill=white, line width=0.6pt] (0.675,0.5) circle (1.5pt);
\filldraw[navy, fill=white, line width=0.6pt] (0.675,1.5) circle (1.5pt);
\filldraw[navy, fill=white, line width=0.6pt] (0.75,1) circle (1.5pt);
\filldraw[navy, fill=white, line width=0.6pt] (0.825,1) circle (1.5pt);
\filldraw[navy, fill=white, line width=0.6pt] (0.825,3) circle (1.5pt);
\filldraw[navy, fill=white, line width=0.6pt] (0.9,1) circle (1.5pt);
\filldraw[navy, fill=white, line width=0.6pt] (0.975,0.5) circle (1.5pt);
\filldraw[navy, fill=white, line width=0.6pt] (0.975,1) circle (1.5pt);
\filldraw[navy, fill=white, line width=0.6pt] (1.05,1.5) circle (1.5pt);
\filldraw[navy, fill=white, line width=0.6pt] (1.125,0.5) circle (1.5pt);
\filldraw[navy, fill=white, line width=0.6pt] (1.2,3) circle (1.5pt);
\filldraw[navy, fill=white, line width=0.6pt] (1.2,3.5) circle (1.5pt);
\filldraw[navy, fill=white, line width=0.6pt] (1.2,4) circle (1.5pt);
\filldraw[navy, fill=white, line width=0.6pt] (1.275,1) circle (1.5pt);
\filldraw[navy, fill=white, line width=0.6pt] (1.275,1.5) circle (1.5pt);
\filldraw[navy, fill=white, line width=0.6pt] (1.275,2) circle (1.5pt);
\filldraw[navy, fill=white, line width=0.6pt] (1.35,1.5) circle (1.5pt);
\filldraw[navy, fill=white, line width=0.6pt] (1.35,2.5) circle (1.5pt);
\filldraw[navy, fill=white, line width=0.6pt] (1.425,0.5) circle (1.5pt);
\filldraw[navy, fill=white, line width=0.6pt] (1.425,1) circle (1.5pt);
\filldraw[navy, fill=white, line width=0.6pt] (1.425,1.5) circle (1.5pt);
\filldraw[navy, fill=white, line width=0.6pt] (1.425,2.5) circle (1.5pt);
\filldraw[navy, fill=white, line width=0.6pt] (1.425,3) circle (1.5pt);
\filldraw[navy, fill=white, line width=0.6pt] (1.425,4.5) circle (1.5pt);
\filldraw[navy, fill=white, line width=0.6pt] (1.5,2.5) circle (1.5pt);
\filldraw[navy, fill=white, line width=0.6pt] (1.5,3) circle (1.5pt);
\filldraw[navy, fill=white, line width=0.6pt] (1.575,0.5) circle (1.5pt);
\filldraw[navy, fill=white, line width=0.6pt] (1.575,2) circle (1.5pt);
\filldraw[navy, fill=white, line width=0.6pt] (1.575,2.5) circle (1.5pt);
\filldraw[navy, fill=white, line width=0.6pt] (1.575,3) circle (1.5pt);
\filldraw[navy, fill=white, line width=0.6pt] (1.65,4.5) circle (1.5pt);
\filldraw[navy, fill=white, line width=0.6pt] (1.725,1.5) circle (1.5pt);
\filldraw[navy, fill=white, line width=0.6pt] (1.725,2.5) circle (1.5pt);
\filldraw[navy, fill=white, line width=0.6pt] (1.725,3) circle (1.5pt);
\filldraw[navy, fill=white, line width=0.6pt] (1.725,3.5) circle (1.5pt);
\filldraw[navy, fill=white, line width=0.6pt] (1.8,2) circle (1.5pt);
\filldraw[navy, fill=white, line width=0.6pt] (1.875,2) circle (1.5pt);
\filldraw[navy, fill=white, line width=0.6pt] (1.875,2.5) circle (1.5pt);
\filldraw[navy, fill=white, line width=0.6pt] (1.875,3) circle (1.5pt);
\filldraw[navy, fill=white, line width=0.6pt] (1.875,5) circle (1.5pt);
\filldraw[navy, fill=white, line width=0.6pt] (1.95,3) circle (1.5pt);
\filldraw[navy, fill=white, line width=0.6pt] (1.95,3.5) circle (1.5pt);
\filldraw[navy, fill=white, line width=0.6pt] (1.95,4.5) circle (1.5pt);
\filldraw[navy, fill=white, line width=0.6pt] (2.025,0.5) circle (1.5pt);
\filldraw[navy, fill=white, line width=0.6pt] (2.025,2.5) circle (1.5pt);
\filldraw[navy, fill=white, line width=0.6pt] (2.025,4) circle (1.5pt);
\filldraw[navy, fill=white, line width=0.6pt] (2.1,2.5) circle (1.5pt);
\filldraw[navy, fill=white, line width=0.6pt] (2.1,3.5) circle (1.5pt);
\filldraw[navy, fill=white, line width=0.6pt] (2.175,1.5) circle (1.5pt);
\filldraw[navy, fill=white, line width=0.6pt] (2.175,4) circle (1.5pt);
\filldraw[navy, fill=white, line width=0.6pt] (2.175,4.5) circle (1.5pt);
\filldraw[navy, fill=white, line width=0.6pt] (2.175,5) circle (1.5pt);
\filldraw[navy, fill=white, line width=0.6pt] (2.325,3) circle (1.5pt);
\filldraw[navy, fill=white, line width=0.6pt] (2.325,4) circle (1.5pt);
\filldraw[navy, fill=white, line width=0.6pt] (2.325,5) circle (1.5pt);
\filldraw[navy, fill=white, line width=0.6pt] (2.4,1.5) circle (1.5pt);
\filldraw[navy, fill=white, line width=0.6pt] (2.4,2.5) circle (1.5pt);
\filldraw[navy, fill=white, line width=0.6pt] (2.4,3.5) circle (1.5pt);
\filldraw[navy, fill=white, line width=0.6pt] (2.4,4.5) circle (1.5pt);
\filldraw[navy, fill=white, line width=0.6pt] (2.475,4.5) circle (1.5pt);
\filldraw[navy, fill=white, line width=0.6pt] (2.475,5) circle (1.5pt);
\filldraw[navy, fill=white, line width=0.6pt] (2.55,1.5) circle (1.5pt);
\filldraw[navy, fill=white, line width=0.6pt] (2.625,2) circle (1.5pt);
\filldraw[navy, fill=white, line width=0.6pt] (2.625,5) circle (1.5pt);
\filldraw[navy, fill=white, line width=0.6pt] (2.775,4.5) circle (1.5pt);
\filldraw[navy, fill=white, line width=0.6pt] (2.85,3) circle (1.5pt);
\filldraw[navy, fill=white, line width=0.6pt] (2.925,4) circle (1.5pt);
\filldraw[navy, fill=white, line width=0.6pt] (2.925,5) circle (1.5pt);
\filldraw[navy, fill=white, line width=0.6pt] (3,3.5) circle (1.5pt);
\filldraw[navy, fill=white, line width=0.6pt] (3,4.5) circle (1.5pt);
\filldraw[navy, fill=white, line width=0.6pt] (3.075,4) circle (1.5pt);
\filldraw[navy, fill=white, line width=0.6pt] (3.15,5) circle (1.5pt);
\filldraw[navy, fill=white, line width=0.6pt] (3.225,3.5) circle (1.5pt);
\filldraw[navy, fill=white, line width=0.6pt] (3.225,5) circle (1.5pt);
\filldraw[navy, fill=white, line width=0.6pt] (3.3,3) circle (1.5pt);
\filldraw[navy, fill=white, line width=0.6pt] (3.3,5) circle (1.5pt);
\filldraw[navy, fill=white, line width=0.6pt] (3.375,1) circle (1.5pt);
\filldraw[navy, fill=white, line width=0.6pt] (3.375,3) circle (1.5pt);
\filldraw[navy, fill=white, line width=0.6pt] (3.375,5) circle (1.5pt);
\filldraw[navy, fill=white, line width=0.6pt] (3.45,3) circle (1.5pt);
\filldraw[navy, fill=white, line width=0.6pt] (3.45,4) circle (1.5pt);
\filldraw[navy, fill=white, line width=0.6pt] (3.675,4) circle (1.5pt);
\filldraw[navy, fill=white, line width=0.6pt] (3.675,4.5) circle (1.5pt);
\filldraw[navy, fill=white, line width=0.6pt] (3.75,2.5) circle (1.5pt);
\filldraw[navy, fill=white, line width=0.6pt] (3.825,1) circle (1.5pt);
\filldraw[navy, fill=white, line width=0.6pt] (3.9,4) circle (1.5pt);
\filldraw[navy, fill=white, line width=0.6pt] (3.975,5) circle (1.5pt);
\filldraw[navy, fill=white, line width=0.6pt] (4.125,2.5) circle (1.5pt);
\filldraw[navy, fill=white, line width=0.6pt] (4.125,3.5) circle (1.5pt);
\filldraw[navy, fill=white, line width=0.6pt] (4.125,5) circle (1.5pt);
\filldraw[navy, fill=white, line width=0.6pt] (4.2,4.5) circle (1.5pt);
\filldraw[navy, fill=white, line width=0.6pt] (4.5,4) circle (1.5pt);
\filldraw[navy, fill=white, line width=0.6pt] (4.875,3.5) circle (1.5pt);
\filldraw[navy, fill=white, line width=0.6pt] (5.175,1.5) circle (1.5pt);
\filldraw[navy, fill=white, line width=0.6pt] (5.325,3.5) circle (1.5pt);
\filldraw[navy, fill=white, line width=0.6pt] (5.625,1.5) circle (1.5pt);
\filldraw[navy, fill=white, line width=0.6pt] (5.7,3.5) circle (1.5pt);
\filldraw[navy, fill=white, line width=0.6pt] (6.9,4.5) circle (1.5pt);
\filldraw[navy, fill=white, line width=0.6pt] (6.975,2) circle (1.5pt);
\filldraw[navy, fill=white, line width=0.6pt] (7.275,4.5) circle (1.5pt);
\filldraw[navy, fill=white, line width=0.6pt] (7.425,2) circle (1.5pt);
\filldraw[navy, fill=white, line width=0.6pt] (8.775,2.5) circle (1.5pt);
\filldraw[navy, fill=white, line width=0.6pt] (9.225,2.5) circle (1.5pt);
\filldraw[navy, fill=white, line width=0.6pt] (10.575,3) circle (1.5pt);
\filldraw[navy, fill=white, line width=0.6pt] (11.025,3) circle (1.5pt);
\fill[wine] (2.025,0.5) circle (1.6pt);
\draw[wine, line width=0.5pt] (2.025,0.5) circle (3.2pt);
\fill[wine] (3.825,1) circle (1.6pt);
\draw[wine, line width=0.5pt] (3.825,1) circle (3.2pt);
\fill[wine] (5.625,1.5) circle (1.6pt);
\draw[wine, line width=0.5pt] (5.625,1.5) circle (3.2pt);
\fill[wine] (7.425,2) circle (1.6pt);
\draw[wine, line width=0.5pt] (7.425,2) circle (3.2pt);
\fill[wine] (9.225,2.5) circle (1.6pt);
\draw[wine, line width=0.5pt] (9.225,2.5) circle (3.2pt);
\fill[wine] (11.025,3) circle (1.6pt);
\draw[wine, line width=0.5pt] (11.025,3) circle (3.2pt);
\draw[ink, line width=0.5pt] (0,0) -- (12,0);
\draw[ink, line width=0.5pt] (0,0) -- (0,5);
\draw[ink, line width=0.5pt] (1.5,0) -- (1.5,-0.08);
\node[ink, below] at (1.5,-0.08) {$20$};
\draw[ink, line width=0.5pt] (3,0) -- (3,-0.08);
\node[ink, below] at (3,-0.08) {$40$};
\draw[ink, line width=0.5pt] (4.5,0) -- (4.5,-0.08);
\node[ink, below] at (4.5,-0.08) {$60$};
\draw[ink, line width=0.5pt] (6,0) -- (6,-0.08);
\node[ink, below] at (6,-0.08) {$80$};
\draw[ink, line width=0.5pt] (7.5,0) -- (7.5,-0.08);
\node[ink, below] at (7.5,-0.08) {$100$};
\draw[ink, line width=0.5pt] (9,0) -- (9,-0.08);
\node[ink, below] at (9,-0.08) {$120$};
\draw[ink, line width=0.5pt] (10.5,0) -- (10.5,-0.08);
\node[ink, below] at (10.5,-0.08) {$140$};
\draw[ink, line width=0.5pt] (12,0) -- (12,-0.08);
\node[ink, below] at (12,-0.08) {$160$};
\draw[ink, line width=0.5pt] (0,0) -- (-0.08,0);
\node[ink, left] at (-0.08,0) {$0$};
\draw[ink, line width=0.5pt] (0,1) -- (-0.08,1);
\node[ink, left] at (-0.08,1) {$2$};
\draw[ink, line width=0.5pt] (0,2) -- (-0.08,2);
\node[ink, left] at (-0.08,2) {$4$};
\draw[ink, line width=0.5pt] (0,3) -- (-0.08,3);
\node[ink, left] at (-0.08,3) {$6$};
\draw[ink, line width=0.5pt] (0,4) -- (-0.08,4);
\node[ink, left] at (-0.08,4) {$8$};
\draw[ink, line width=0.5pt] (0,5) -- (-0.08,5);
\node[ink, left] at (-0.08,5) {$10$};
\node[ink] at (6,-0.75) {$P+Q$};
\node[ink, rotate=90] at (-0.75,2.5) {margin};
\node[slate, font=\scriptsize, rotate=15.524, anchor=north] at (10.275,2.492) {$(P+Q-3)/24$};
\end{tikzpicture}

\caption{The margin against $P+Q$ for the $120$ occurrences with
$P+Q\leqslant160$ and margin at most $10$ (hollow circles; several
occurrences may share a mark). The dashed line is $(P+Q-3)/24$, and the
shaded region below it contains no occurrence. The filled wine circles are
the occurrences $(9k+1)/(15k+2)$, $k=1,\dots,6$, which lie on the line and
whose margin is exactly $k$; \cref{prop:two-ladders} shows that the margin
of the next ones on the line is at most $k$ for every $k$.}
\label{fig:thresholds}
\end{figure}

\begin{conjecture}[Thresholds of the margin]\label{conj:thresholds}
For every integer $k\geqslant0$, every occurrence $P/Q$ with $P+Q>24k+3$ has
margin larger than $k$, the margin of a label with a single primitive
representation being taken as infinite.
\end{conjecture}

The case $k=0$ asserts that every occurrence is strictly minimal, which is
stronger than \cref{conj:euclidean-minimality} and therefore implies the
Frobenius--Markov uniqueness conjecture. The conjecture holds for every
occurrence with $P+Q\leqslant160$, and \cref{prop:two-ladders} shows that it
would be false with any smaller threshold.

The minimality is not a consequence of the digit pattern alone. The
representation of a Markov label by its occurrence pair has a half continued
fraction with digits $1$ and $2$ only \cite[Corollary~4.4]{CanakciSchiffler2020},
but such digits do not force minimality for other integers. The integer
$135328989953=37\cdot113\cdot761\cdot42533$, which is not a Markov number,
has eight primitive representations; two of them,
$196223^2+311168^2$ and $131468^2+343577^2$, have continued fractions with
digits $1$ and $2$ only, and their costs are $28$ and $27$. The first is
therefore not minimal. Whatever makes the conjecture true must use the
Markov structure of the label, as the path comparison of this paper does on
thirteen rows.

\subsection{Directives of the Christoffel words of a given length}\label{sec:directives}

The comparison of this paper fixes a label and varies its representation as
a sum of two squares. One may instead fix an integer $m$ and vary the
Christoffel word of length $m$. For each $n$ prime to $m$ with $1\leqslant n<m/2$,
the lower Christoffel word $\gamma_{n/m}$ has a central word, and the central
word has a directive, in the sense of the iterated palindromic closure of de
Luca. This subsection shows that these directives are the Euclidean
directives of \cref{lem:euclidean-bijection}, characterizes those that are
images of the Thue--Morse morphism, and states a conjecture on them that
implies the uniqueness conjecture. It also recasts
\cref{conj:euclidean-minimality} in terms of the square roots of $-1$ modulo
the label.

The palindromic closure $w^{(+)}$ of a word $w$ is the shortest palindrome
having $w$ as a prefix, and the iterated palindromic closure is defined by
$\operatorname{Pal}(\varepsilon)=\varepsilon$ and
$\operatorname{Pal}(vx)=(\operatorname{Pal}(v)x)^{(+)}$ for a word $v$ and a letter $x$
\cite{DeLuca1997}. The map $\operatorname{Pal}$ is injective, and $d$ is called the
\emph{directive} of the central word $\operatorname{Pal}(d)$; the map
$d\mapsto a\operatorname{Pal}(d)b$ is a bijection from the binary words onto the
proper lower Christoffel words
\cite[Section~2 and the proof of Theorem~1]{ReutenauerVuillon2017}, a consequence of
\cite[Corollary~3.1]{BertheDeLucaReutenauer2008}. The Thue--Morse morphism
$\theta$ sends $a$ to $ab$ and $b$ to $ba$; its image is the set $\{ab,ba\}^*$
of the products of the words $ab$ and $ba$.

\begin{lemma}[Directives of Christoffel words]\label{lem:directives}
Let $1\leqslant n<m/2$ be coprime integers, write $\gamma_{n/m}=azb$, and let $d$
be the directive of $z$. Then $L_d{\bf1}=(m-n,n)^{\mathsf T}$: the word $d$ is the
directive of the pair $(m-n,n)$ in the sense of \cref{lem:euclidean-bijection},
and $\abs d=\ell(n,m)-1$. If $n=1$, then $d=a^{m-2}$. If $n\geqslant2$ and
$m/n=[c_0;c_1,\ldots,c_k]$ is the canonical continued fraction, the runs of $d$
have the lengths $c_0-1,c_1,\ldots,c_{k-1},c_k-1$, alternately of the letters
$a$ and $b$, beginning with $a$. In particular, $d$ lies in $\{ab,ba\}^*$ if and
only if
\begin{equation}\label{eq:block-form}
 \frac mn=[2;B_1,\ldots,B_j,2],\qquad j\geqslant0,\quad B_i\in\{(2),\,(1,1)\},
\end{equation}
that is, if and only if the first and the last partial quotients equal $2$
and those in between split into consecutive blocks $(2)$ and $(1,1)$.
\end{lemma}

\begin{proof}
\emph{The pair.} Berth\'e, de Luca and Reutenauer attach to a word $d$ the
incidence matrix $M_d$ of a composition of morphisms, with $M_{dd'}=M_dM_{d'}$,
$M_a=\left(\begin{smallmatrix}1&1\\0&1\end{smallmatrix}\right)=L_a$ and
$M_b=\left(\begin{smallmatrix}1&0\\1&1\end{smallmatrix}\right)=L_b$
\cite[proofs of Proposition~3.1 and Corollary~3.2]{BertheDeLucaReutenauer2008},
so that $M_d=L_d$. By \cite[Corollary~3.1]{BertheDeLucaReutenauer2008}, the
Christoffel word $a\operatorname{Pal}(d)b$ has as many letters $a$ as the sum
of the entries of the first row of $L_d$, and as many letters $b$ as the sum
of the entries of its second row; that is, its two letter counts are the
entries of $L_d{\bf1}$. For $z=\operatorname{Pal}(d)$, \eqref{eq:farey-counts}
therefore gives $L_d{\bf1}=(m-n,n)^{\mathsf T}$, and \cref{lem:euclidean-bijection}
shows that $d$ is the directive of $(m-n,n)$, of length $\ell(m-n,n)$. Since
$m-n>n$, the first subtraction carries the pair $(n,m)$ to $(n,m-n)$, so that
$\ell(n,m)=1+\ell(n,m-n)$.

\emph{The runs.} If $n=1$, the pair $(m-1,1)$ reaches $(1,1)$ after $m-2$
subtractions of its second coordinate, and its directive is $a^{m-2}$. If
$n\geqslant2$, then $m/n$ is not an integer, so $k\geqslant1$; moreover $c_0\geqslant2$
because $m>2n$, and $(m-n)/n=m/n-1=[c_0-1;c_1,\ldots,c_k]$ is a canonical
expansion. The description of directives after \cref{lem:euclidean-bijection}
gives the runs $c_0-1,c_1,\ldots,c_{k-1},c_k-1$, alternately of $a$ and $b$,
beginning with $a$.

\emph{The image of $\theta$.} A word lies in $\{ab,ba\}^*$ exactly when its
length is even and its letters in the positions $2i-1$ and $2i$ differ for
every $i$, that is, when a run ends after every odd position. If $n=1$, the
word $a^{m-2}$, with $m\geqslant3$, has no such run end, and \eqref{eq:block-form}
fails since $m/1$ has a single partial quotient. Let $n\geqslant2$, and write
$r_0=c_0-1$, $r_i=c_i$ for $0<i<k$ and $r_k=c_k-1$, all positive; the runs end
after the positions $r_0$, $r_0+r_1$, \dots, $r_0+\cdots+r_{k-1}$. Suppose that
$d$ lies in $\{ab,ba\}^*$, of length $2h$. A run ending after the position $1$
forces $r_0=1$, that is $c_0=2$. A run ending after the last odd position
$2h-1$ leaves a last run of length $1$, so $r_k=1$, that is $c_k=2$. Between
two consecutive odd positions $2i-1$ and $2i+1$, there is either no other run
end, which makes a run of length $2$, or one run end after the position $2i$,
which makes two runs of length $1$; so the runs $r_1,\ldots,r_{k-1}$, which are
the quotients $c_1,\ldots,c_{k-1}$, split into consecutive blocks $(2)$ and
$(1,1)$. Conversely, when \eqref{eq:block-form} holds, the run lengths are
$1$, then blocks $(2)$ and $(1,1)$, then $1$, so the word has even length and
a run ends after every odd position.
\end{proof}

Every occurrence supplies an integer of the form \eqref{eq:block-form}. Let
$P/Q$ be an occurrence of the label $m$, with central word $z$, and write
$\mu(z)e=(x,y)^{\mathsf T}$. Since $\mu(a)=P_1^2$, $\mu(b)=P_2^2$ and
$P_2(1,0)^{\mathsf T}=e$, the column
$P_2\,\mu(z)\,P_2(1,0)^{\mathsf T}=P_2\,\mu(z)\,e=(2x+y,x)^{\mathsf T}$ is the image
of $(1,0)^{\mathsf T}$ under the product of the matrices $P_c$, $c$ running
through the sequence $(2,\sigma(z),2)$ of \cref{sec:occurrences}, and its
first entry $2x+y=e^{\mathsf T}\mu(z)e$ is $m$ by \eqref{eq:square-count-identity}.
By \cref{lem:cost-descriptions}, the integer $u_{P/Q}=x$, which we call the
\emph{root} of the occurrence, satisfies
\begin{equation}\label{eq:root-cf}
 \frac m{u_{P/Q}}=[2;\sigma(z),2].
\end{equation}
The sequence $\sigma(z)$ consists of the blocks $(1,1)$ and $(2,2)$, so the
root is of the form \eqref{eq:block-form}, and $m/3<u_{P/Q}<m/2$. The matrix
$P_2\,\mu(z)\,P_2$ is symmetric of determinant one, with first column
$(m,u_{P/Q})^{\mathsf T}$, so that $u_{P/Q}^2\equiv-1\pmod m$. Since $\mu(a)e_1=e$,
$\mu(b)e_1=2e+e_1$ and $\mu(z)$ is symmetric, the upper left entry of
$\mu(\gamma_{P/Q})=\mu(a)\mu(z)\mu(b)$ is
$e^{\mathsf T}\mu(z)(2e+e_1)=2m+u_{P/Q}$. By the first step of the proof of
\cref{prop:occurrence-in-cube}, $u_{P/Q}\,m_R\equiv-m_L\pmod m$, and since
$0<u_{P/Q}<m/2$, the root is the characteristic number of the Markov triple
$(m,m_L,m_R)$ in the normalization of \cite[Section~2.2]{Aigner2013}, the
smaller of the two solutions in $(0,m)$ of $m_Rx\equiv\pm m_L\pmod m$.

\begin{conjecture}[Directives in the image of the Thue--Morse morphism]\label{conj:directives}
For every Markov number $m\geqslant5$, exactly one integer $n$ with
$1\leqslant n<m/2$, prime to $m$, has the property that the directive of the
central word of $\gamma_{n/m}$ lies in $\{ab,ba\}^*$; equivalently, by
\cref{lem:directives}, exactly one such $n$ satisfies \eqref{eq:block-form}.
\end{conjecture}

\begin{proposition}[The conjecture implies uniqueness]\label{prop:directives-uniqueness}
\Cref{conj:directives} implies the Frobenius--Markov uniqueness conjecture.
In terms of the Thue--Morse morphism, it states that for every Markov number
$m\geqslant5$ exactly one word $w$ beginning with the letter $a$ satisfies
$\abs{\operatorname{Pal}(\theta(w))}=m-2$.
\end{proposition}

\begin{proof}
\emph{Uniqueness.} Every occurrence $P/Q$ of a Markov number $m$ supplies the
integer $n=u_{P/Q}$, which satisfies \eqref{eq:block-form} by
\eqref{eq:root-cf}. Two distinct occurrences supply distinct integers. Indeed,
the canonical continued fraction \eqref{eq:root-cf} of $m/u_{P/Q}$ determines
the sequence $\sigma(z)$; this sequence determines $z$, because $\sigma$ replaces
each letter by two equal entries that identify it; and $z$ determines the word
$\gamma_{P/Q}=azb$ and the fraction $P/Q$, by \eqref{eq:farey-counts}. If the
integer of the conjecture is unique, so is the occurrence.

\emph{The second form.} Let $n$ be an integer as in the conjecture, and $d$ the
directive of the central word of $\gamma_{n/m}$. By \cref{lem:directives}, $d$
begins with $a$, and it lies in $\{ab,ba\}^*$ exactly when $d=\theta(w)$ for a
word $w$, which then begins with $a$ and is unique, the morphism $\theta$ being
injective. Conversely, let $w$ begin with $a$. Then $d=\theta(w)$ begins with
$a$, and $a\operatorname{Pal}(d)b$ is a proper lower Christoffel word whose numbers
of letters $a$ and $b$ are the entries of $L_d{\bf1}$, by the first step of the
proof of \cref{lem:directives}; the first entry is the larger, since
$L_d{\bf1}=L_aL_{d'}{\bf1}$ with $d=ad'$. This word is therefore
$\gamma_{n/m}$ for coprime integers with $n<m/2$ and
$m=\abs{\operatorname{Pal}(\theta(w))}+2$, and the directive of its central word is
$\theta(w)$. The integers $n$ of the conjecture for the label $m$ thus
correspond one to one to the words $w$ beginning with $a$ such that
$\abs{\operatorname{Pal}(\theta(w))}=m-2$.
\end{proof}

Reutenauer and Vuillon proved that, for every word $v$, the number
$\abs{\operatorname{Pal}(\theta(\operatorname{Pal}(av)))}+2$ is a Markov number, that every
Markov number $m\geqslant5$ arises in this way, and that the injectivity of the
map $v\mapsto\abs{\operatorname{Pal}(\theta(\operatorname{Pal}(av)))}$ is equivalent to the
uniqueness conjecture \cite[Theorem~1]{ReutenauerVuillon2017}. The words
$\operatorname{Pal}(av)$ are the central words beginning with $a$. In the second
form of \cref{prop:directives-uniqueness}, the conjecture contains this
injectivity, and asserts in addition that no word $w$ beginning with $a$ that
is not a central word gives a Markov number. The next lemma supports it for
every Markov number up to $10^{12}$.

\begin{lemma}[Finite facts on directives]\label{lem:directive-facts}
\begin{enumerate}[label=\textup{(\roman*)},leftmargin=*]
\item There are $150$ Markov numbers $m$ with $5\leqslant m\leqslant10^{12}$. Each
 is the label of exactly one occurrence $P/Q$, and its root $u_{P/Q}$ is the
 only integer $n$ prime to $m$ with $1\leqslant n<m/2$ satisfying
 \eqref{eq:block-form}.
\item The integers $31$ and $1130$, which are not Markov numbers, have exactly
 two such integers prime to them each: $31/12=[2;1,1,2,2]$ and $31/13=[2;2,1,1,2]$, and
 $1130/437=[2;\sigma(abba),2]$ and $1130/467=[2;\sigma(baab),2]$.
\item For $m=m_{9/10}=6\,625\,109$, the root is $u_{9/10}=2\,744\,210$, with
 $m/u_{9/10}=[2;2,\ldots,2]$, eighteen partial quotients equal to $2$, so that
 $\ell(u_{9/10},m)=35$. The integer $n=2\,561\,343$ is prime to $m$ and smaller
 than $m/2$, and
 $$
 \frac mn=[2;1,1,2,2,1,1,2,1,1,1,2,1,1,1,2,1,2,2,1,1,1,1,2,2],
 $$
 so that $\ell(n,m)=34$.
\end{enumerate}
\end{lemma}

\begin{proof}
These are finitely many assertions about integers. They were verified by
computer in exact arithmetic, and the verification was carried out twice
independently. The first verification generates the labels from the Markov
triples of \cref{lem:farey-triangle}, takes the roots from the residues of
$-m_L/m_R$ modulo $m$ given by the first step of the proof of
\cref{prop:occurrence-in-cube}, and searches, for each label, the sequences of
blocks from the left. If a prefix $[2;B_1,\ldots,B_i]$ has $p$ and $p'$ as the
numerators of its last two convergents, every completion has a tail column
$(X,Y)^{\mathsf T}$ with coprime entries $0<Y<X$ and $m=pX+p'Y$; hence
$Y<m/(p+p')$ and $Y$ lies in one residue class modulo $p$, so that once
$p(p+p')>m$ at most one candidate remains, and it is tested directly. The
second verification enumerates all the continued fractions
\eqref{eq:block-form} with numerator at most $10^{12}$ and compares their
numerators and denominators with the labels and roots computed from the
Cohn matrices \eqref{eq:cohn-source}. The continued fractions of the last two
items can also be checked by hand.
\end{proof}

The integers $31$ and $1130$ of \cref{lem:directive-facts} show that
\cref{conj:directives} depends on the Markov structure of $m$. Whenever the block sequence $B_1,\ldots,B_j$ of \eqref{eq:block-form}
is not a palindrome, the reversed expansion has the same numerator, since
the product of the symmetric matrices $P_c$ in reverse order is the
transpose, and a different denominator; this is the case of $31$. For
$1130$, the two expansions come from the palindromes $abba$ and $baab$, which
are not central words and have the same label $e^{\mathsf T}\mu(z)e$. The
conjecture therefore asserts that a Markov number is the numerator of no
expansion \eqref{eq:block-form} with a non-palindromic block sequence, and of
no palindromic one other than that of its root.

The roots of the occurrences are square roots of $-1$, and the square roots
of $-1$ carry the whole two-square fiber, as the next lemma shows. This gives
\cref{conj:euclidean-minimality} a form in which only the label and its square
roots of $-1$ appear. In substance the lemma is classical. Smith derived from
two identities of Euler that the numerator of a palindromic continued fraction
with an even number of partial quotients is a sum of two coprime squares, given
by its first half. He proved in this way that every prime $p\equiv1\pmod4$ is a
sum of two squares, since a parity argument shows that one of the expansions of
$p/u$ with $1<u<p/2$ is a palindrome, and he observed that the same argument
applies to every divisor of an integer $1+x^2$ \cite{Smith1855}. Brillhart read
the two squares off the remainders of the Euclidean algorithm on $p$ and the
square root $u$ of $-1$, refining a method of Hermite \cite{Brillhart1972}.
\c{C}anak\c{c}{\i} and Schiffler relate the primitive representations of an
integer to its even palindromic continued fractions \cite[Theorem~3.10 and
Corollary~3.14]{CanakciSchiffler2020}. What the lemma adds is a complete proof,
for every label, of the correspondence between the square roots of $-1$ in
$(1,m/2)$ and the fiber $\mathcal F(m)$, and the identity \eqref{eq:root-cost}
between the costs, on which the reformulation of the conjecture below rests.

\begin{lemma}[Square roots of $-1$]\label{lem:square-roots}
Let $m\geqslant5$ be an integer with a primitive representation as a sum of two
squares, and let $u$ be an integer with $1<u<m/2$ and $u^2\equiv-1\pmod m$.
Then the canonical continued fraction of $m/u$ is a palindrome
$[c_0;c_1,\ldots,c_{2h-1}]$ of even length. If $s/r=[c_{h-1};c_{h-2},\ldots,c_0]$
in lowest terms, then $m=r^2+s^2$, $0<r<s$, and
\begin{equation}\label{eq:root-cost}
 \ell(u,m)=2\,\ell(r,s)+1 .
\end{equation}
The map $u\mapsto(r,s)$ is a bijection from these integers $u$ onto
$\mathcal F(m)$, and it sends the root of an occurrence to its occurrence pair.
\end{lemma}

\begin{proof}
\emph{The palindrome.} Write $m/u=[c_0;c_1,\ldots,c_k]$ canonically. Here
$k\geqslant1$, because $u>1$ is prime to $m$; $c_0\geqslant2$, because $u<m/2$; and
$c_k\geqslant2$. By \cref{lem:cost-descriptions}, the product
$\Pi=P_{c_0}\cdots P_{c_k}$ has first column $(m,u)^{\mathsf T}$; write
$\Pi=\left(\begin{smallmatrix}m&m'\\u&u'\end{smallmatrix}\right)$. Each $P_c$ is
symmetric, so $\Pi^{\mathsf T}=P_{c_k}\cdots P_{c_0}$, whose first column is
$(m,m')^{\mathsf T}$, and \cref{lem:cost-descriptions} gives
$m/m'=[c_k;\ldots,c_0]$, a canonical expansion since $c_0\geqslant2$. As $c_k\geqslant2$
and $k\geqslant1$, $m/m'>2$, so that $0<m'<m/2$. The determinant of $\Pi$ is
$(-1)^{k+1}$, so $mu'-m'u=(-1)^{k+1}$, and $m'u\equiv(-1)^k\pmod m$. Since
$u^2\equiv-1$, the inverse of $u$ modulo $m$ is $-u$, whence
$m'\equiv(-1)^{k+1}u\pmod m$. If $k$ were even, then $m'\equiv-u$ and
$0<m'<m$ would give $m'=m-u>m/2$, a contradiction. Hence $k$ is odd,
$m'\equiv u$, and $m'=u$, both lying in $(0,m)$. The two canonical expansions
$[c_0;\ldots,c_k]$ and $[c_k;\ldots,c_0]$ of $m/u$ then coincide, so that
$c_i=c_{k-i}$ for every $i$, and $k+1=2h$ is even.

\emph{The two squares.} By the symmetry of the expansion, the second half of
the product is $P_{c_{h-1}}\cdots P_{c_0}=\Pi_0^{\mathsf T}$, where
$\Pi_0=P_{c_0}\cdots P_{c_{h-1}}$, so that $\Pi=\Pi_0\Pi_0^{\mathsf T}$. If
$(s',r')$ is the first row of $\Pi_0$, the upper left entry of $\Pi$ is
$m=s'^2+r'^2$. The first column of $\Pi_0^{\mathsf T}=P_{c_{h-1}}\cdots P_{c_0}$ is
$(s',r')^{\mathsf T}$, so \cref{lem:cost-descriptions} shows that $s'$ and $r'$ are
coprime and positive, with $s'/r'=[c_{h-1};\ldots,c_0]$, a canonical
expansion since $c_0\geqslant2$; hence $(s',r')=(s,r)$, and $r<s$ because this
expansion is not $[1]$. Finally \eqref{eq:quotient-cost} gives
$\ell(u,m)=2(c_0+\cdots+c_{h-1})-1$ and $\ell(r,s)=c_0+\cdots+c_{h-1}-1$, which
is \eqref{eq:root-cost}.

\emph{The bijection.} The pair $(r,s)$ determines $u$, because the canonical
expansion of $s/r$ is the reversed first half of the palindrome, which gives
the whole palindrome and hence $m/u$. The map is therefore injective. By \cref{prop:gaussian-fibre},
$\abs{\mathcal F(m)}=2^{t-1}$, where $t$ is the number of distinct odd primes
dividing $m$, as in \cref{prop:gaussian-fibre}. By the Chinese remainder theorem and Hensel's lemma, $-1$ has
exactly $2^{t}$ square roots modulo $m$, two modulo each odd prime power
dividing $m$ and one modulo $2$ when $m$ is even. They pair off as $u$ and
$m-u$, and none of them is $1$ or $m/2$, because $m\geqslant5$; so exactly
$2^{t-1}$ of them lie in $(1,m/2)$. An injective map between two finite sets
of the same size is a bijection. Finally, for the root of an occurrence, the
expansion is \eqref{eq:root-cf}, a palindrome because $z$ is one, and the
reversal of its first half is the right half $[\sigma(w),2]$ or
$[j,\sigma(w),2]$ of the paragraph after \cref{prop:occurrence-cost}, which is
the continued fraction of the occurrence pair.
\end{proof}

By \cref{lem:square-roots} and \cref{prop:occurrence-cost},
\cref{conj:euclidean-minimality} is equivalent to the following statement:
for every occurrence $P/Q$ of a label $m$, no square root of $-1$ modulo $m$
in $(1,m/2)$ has a shorter Euclidean algorithm with $m$ than the root, whose
cost is $\ell(u_{P/Q},m)=2(P+Q-2)+1$. \Cref{fig:directives} shows the label
$m_{4/5}=985$, whose two representations differ in cost by one. The
restriction to square roots of $-1$ cannot be dropped, even on the rows of
\cref{thm:thirteen-rows}: by the last item of \cref{lem:directive-facts}, the
integer $n=2\,561\,343$, which is not a square root of $-1$ modulo $m_{9/10}$
since its cost is even, has a shorter Euclidean algorithm with $m_{9/10}$
than the root. Among all the integers $n$ prime to $m$, the condition
\eqref{eq:block-form} isolates the root where the minimum of the cost does
not.

\begin{figure}[tb]
\centering
% Authors: Dr. Denys Dutykh (Mathematics Department, Khalifa University of Science and
% Technology, Abu Dhabi, UAE) and Prof. Laurent Vuillon (Univ. Savoie Mont Blanc, CNRS, LAMA)
% Generated by supplement/code/make_figures.py from exact data; do not edit by hand.
\begin{tikzpicture}[font=\footnotesize]
\draw[slate, line width=0.6pt, dash pattern=on 4pt off 2pt] (-0.2,0.458) -- (12,0.458);
\filldraw[navy, fill=white, line width=0.45pt] (0.024,5.182) circle (1.15pt);
\filldraw[navy, fill=white, line width=0.45pt] (0.048,4.401) circle (1.15pt);
\filldraw[navy, fill=white, line width=0.45pt] (0.072,3.948) circle (1.15pt);
\filldraw[navy, fill=white, line width=0.45pt] (0.096,3.63) circle (1.15pt);
\filldraw[navy, fill=white, line width=0.45pt] (0.144,3.193) circle (1.15pt);
\filldraw[navy, fill=white, line width=0.45pt] (0.168,3.012) circle (1.15pt);
\filldraw[navy, fill=white, line width=0.45pt] (0.192,2.896) circle (1.15pt);
\filldraw[navy, fill=white, line width=0.45pt] (0.216,2.748) circle (1.15pt);
\filldraw[navy, fill=white, line width=0.45pt] (0.264,2.542) circle (1.15pt);
\filldraw[navy, fill=white, line width=0.45pt] (0.288,2.518) circle (1.15pt);
\filldraw[navy, fill=white, line width=0.45pt] (0.312,2.362) circle (1.15pt);
\filldraw[navy, fill=white, line width=0.45pt] (0.336,2.29) circle (1.15pt);
\filldraw[navy, fill=white, line width=0.45pt] (0.384,2.148) circle (1.15pt);
\filldraw[navy, fill=white, line width=0.45pt] (0.408,2.245) circle (1.15pt);
\filldraw[navy, fill=white, line width=0.45pt] (0.432,2.023) circle (1.15pt);
\filldraw[navy, fill=white, line width=0.45pt] (0.456,2.004) circle (1.15pt);
\filldraw[navy, fill=white, line width=0.45pt] (0.504,1.965) circle (1.15pt);
\filldraw[navy, fill=white, line width=0.45pt] (0.528,1.84) circle (1.15pt);
\filldraw[navy, fill=white, line width=0.45pt] (0.552,1.817) circle (1.15pt);
\filldraw[navy, fill=white, line width=0.45pt] (0.576,2.096) circle (1.15pt);
\filldraw[navy, fill=white, line width=0.45pt] (0.624,1.747) circle (1.15pt);
\filldraw[navy, fill=white, line width=0.45pt] (0.648,1.817) circle (1.15pt);
\filldraw[navy, fill=white, line width=0.45pt] (0.672,1.647) circle (1.15pt);
\filldraw[navy, fill=white, line width=0.45pt] (0.696,2.042) circle (1.15pt);
\filldraw[navy, fill=white, line width=0.45pt] (0.744,1.537) circle (1.15pt);
\filldraw[navy, fill=white, line width=0.45pt] (0.768,1.507) circle (1.15pt);
\filldraw[navy, fill=white, line width=0.45pt] (0.792,1.507) circle (1.15pt);
\filldraw[navy, fill=white, line width=0.45pt] (0.816,2.042) circle (1.15pt);
\filldraw[navy, fill=white, line width=0.45pt] (0.864,1.415) circle (1.15pt);
\filldraw[navy, fill=white, line width=0.45pt] (0.888,1.382) circle (1.15pt);
\filldraw[navy, fill=white, line width=0.45pt] (0.912,1.537) circle (1.15pt);
\filldraw[navy, fill=white, line width=0.45pt] (0.936,1.477) circle (1.15pt);
\filldraw[navy, fill=white, line width=0.45pt] (0.984,2.096) circle (1.15pt);
\filldraw[navy, fill=white, line width=0.45pt] (1.008,1.313) circle (1.15pt);
\filldraw[navy, fill=white, line width=0.45pt] (1.032,1.415) circle (1.15pt);
\filldraw[navy, fill=white, line width=0.45pt] (1.056,1.241) circle (1.15pt);
\filldraw[navy, fill=white, line width=0.45pt] (1.104,1.202) circle (1.15pt);
\filldraw[navy, fill=white, line width=0.45pt] (1.128,1.673) circle (1.15pt);
\filldraw[navy, fill=white, line width=0.45pt] (1.152,1.382) circle (1.15pt);
\filldraw[navy, fill=white, line width=0.45pt] (1.176,1.348) circle (1.15pt);
\filldraw[navy, fill=white, line width=0.45pt] (1.224,1.202) circle (1.15pt);
\filldraw[navy, fill=white, line width=0.45pt] (1.248,1.477) circle (1.15pt);
\filldraw[navy, fill=white, line width=0.45pt] (1.272,1.122) circle (1.15pt);
\filldraw[navy, fill=white, line width=0.45pt] (1.296,1.202) circle (1.15pt);
\filldraw[navy, fill=white, line width=0.45pt] (1.344,1.079) circle (1.15pt);
\filldraw[navy, fill=white, line width=0.45pt] (1.368,1.079) circle (1.15pt);
\filldraw[navy, fill=white, line width=0.45pt] (1.392,2.245) circle (1.15pt);
\filldraw[navy, fill=white, line width=0.45pt] (1.416,1.035) circle (1.15pt);
\filldraw[navy, fill=white, line width=0.45pt] (1.464,1.122) circle (1.15pt);
\filldraw[navy, fill=white, line width=0.45pt] (1.488,1.202) circle (1.15pt);
\filldraw[navy, fill=white, line width=0.45pt] (1.512,1.035) circle (1.15pt);
\filldraw[navy, fill=white, line width=0.45pt] (1.536,0.989) circle (1.15pt);
\filldraw[navy, fill=white, line width=0.45pt] (1.584,1.278) circle (1.15pt);
\filldraw[navy, fill=white, line width=0.45pt] (1.608,1.035) circle (1.15pt);
\filldraw[navy, fill=white, line width=0.45pt] (1.632,1.348) circle (1.15pt);
\filldraw[navy, fill=white, line width=0.45pt] (1.656,0.94) circle (1.15pt);
\filldraw[navy, fill=white, line width=0.45pt] (1.704,1.163) circle (1.15pt);
\filldraw[navy, fill=white, line width=0.45pt] (1.728,1.035) circle (1.15pt);
\filldraw[navy, fill=white, line width=0.45pt] (1.752,1.817) circle (1.15pt);
\filldraw[navy, fill=white, line width=0.45pt] (1.776,0.94) circle (1.15pt);
\filldraw[navy, fill=white, line width=0.45pt] (1.824,1.537) circle (1.15pt);
\filldraw[navy, fill=white, line width=0.45pt] (1.848,0.94) circle (1.15pt);
\filldraw[navy, fill=white, line width=0.45pt] (1.872,0.89) circle (1.15pt);
\filldraw[navy, fill=white, line width=0.45pt] (1.896,0.989) circle (1.15pt);
\filldraw[navy, fill=white, line width=0.45pt] (1.944,0.989) circle (1.15pt);
\filldraw[navy, fill=white, line width=0.45pt] (1.968,2.518) circle (1.15pt);
\filldraw[navy, fill=white, line width=0.45pt] (1.992,0.989) circle (1.15pt);
\filldraw[navy, fill=white, line width=0.45pt] (2.016,0.94) circle (1.15pt);
\filldraw[navy, fill=white, line width=0.45pt] (2.064,0.989) circle (1.15pt);
\filldraw[navy, fill=white, line width=0.45pt] (2.088,1.035) circle (1.15pt);
\filldraw[navy, fill=white, line width=0.45pt] (2.112,0.94) circle (1.15pt);
\filldraw[navy, fill=white, line width=0.45pt] (2.136,1.241) circle (1.15pt);
\filldraw[navy, fill=white, line width=0.45pt] (2.184,0.89) circle (1.15pt);
\filldraw[navy, fill=white, line width=0.45pt] (2.208,0.838) circle (1.15pt);
\filldraw[navy, fill=white, line width=0.45pt] (2.232,0.838) circle (1.15pt);
\filldraw[navy, fill=white, line width=0.45pt] (2.256,1.079) circle (1.15pt);
\filldraw[navy, fill=white, line width=0.45pt] (2.304,0.89) circle (1.15pt);
\filldraw[navy, fill=white, line width=0.45pt] (2.328,0.989) circle (1.15pt);
\filldraw[navy, fill=white, line width=0.45pt] (2.352,1.313) circle (1.15pt);
\filldraw[navy, fill=white, line width=0.45pt] (2.376,1.313) circle (1.15pt);
\filldraw[navy, fill=white, line width=0.45pt] (2.424,1.477) circle (1.15pt);
\filldraw[navy, fill=white, line width=0.45pt] (2.448,0.989) circle (1.15pt);
\filldraw[navy, fill=white, line width=0.45pt] (2.472,0.838) circle (1.15pt);
\filldraw[navy, fill=white, line width=0.45pt] (2.496,0.989) circle (1.15pt);
\filldraw[navy, fill=white, line width=0.45pt] (2.544,0.725) circle (1.15pt);
\filldraw[navy, fill=white, line width=0.45pt] (2.568,0.89) circle (1.15pt);
\filldraw[navy, fill=white, line width=0.45pt] (2.592,0.94) circle (1.15pt);
\filldraw[navy, fill=white, line width=0.45pt] (2.616,1.537) circle (1.15pt);
\filldraw[navy, fill=white, line width=0.45pt] (2.664,1.163) circle (1.15pt);
\filldraw[navy, fill=white, line width=0.45pt] (2.688,0.838) circle (1.15pt);
\filldraw[navy, fill=white, line width=0.45pt] (2.712,0.838) circle (1.15pt);
\filldraw[navy, fill=white, line width=0.45pt] (2.736,0.725) circle (1.15pt);
\filldraw[navy, fill=white, line width=0.45pt] (2.784,1.537) circle (1.15pt);
\filldraw[navy, fill=white, line width=0.45pt] (2.808,0.664) circle (1.15pt);
\filldraw[navy, fill=white, line width=0.45pt] (2.832,0.89) circle (1.15pt);
\filldraw[navy, fill=white, line width=0.45pt] (2.856,0.783) circle (1.15pt);
\filldraw[navy, fill=white, line width=0.45pt] (2.904,0.989) circle (1.15pt);
\filldraw[navy, fill=white, line width=0.45pt] (2.928,1.079) circle (1.15pt);
\filldraw[navy, fill=white, line width=0.45pt] (2.952,2.896) circle (1.15pt);
\filldraw[navy, fill=white, line width=0.45pt] (2.976,1.163) circle (1.15pt);
\filldraw[navy, fill=white, line width=0.45pt] (3.024,0.989) circle (1.15pt);
\filldraw[navy, fill=white, line width=0.45pt] (3.048,0.94) circle (1.15pt);
\filldraw[navy, fill=white, line width=0.45pt] (3.072,0.725) circle (1.15pt);
\filldraw[navy, fill=white, line width=0.45pt] (3.096,0.94) circle (1.15pt);
\filldraw[navy, fill=white, line width=0.45pt] (3.144,0.989) circle (1.15pt);
\filldraw[navy, fill=white, line width=0.45pt] (3.168,0.989) circle (1.15pt);
\filldraw[navy, fill=white, line width=0.45pt] (3.192,0.783) circle (1.15pt);
\filldraw[navy, fill=white, line width=0.45pt] (3.216,0.725) circle (1.15pt);
\filldraw[navy, fill=white, line width=0.45pt] (3.264,0.89) circle (1.15pt);
\filldraw[navy, fill=white, line width=0.45pt] (3.288,0.725) circle (1.15pt);
\filldraw[navy, fill=white, line width=0.45pt] (3.312,0.838) circle (1.15pt);
\filldraw[navy, fill=white, line width=0.45pt] (3.336,0.94) circle (1.15pt);
\filldraw[navy, fill=white, line width=0.45pt] (3.384,2.305) circle (1.15pt);
\filldraw[navy, fill=white, line width=0.45pt] (3.408,1.079) circle (1.15pt);
\filldraw[navy, fill=white, line width=0.45pt] (3.432,1.202) circle (1.15pt);
\filldraw[navy, fill=white, line width=0.45pt] (3.456,0.783) circle (1.15pt);
\filldraw[navy, fill=white, line width=0.45pt] (3.504,1.163) circle (1.15pt);
\filldraw[navy, fill=white, line width=0.45pt] (3.528,1.035) circle (1.15pt);
\filldraw[navy, fill=white, line width=0.45pt] (3.552,0.89) circle (1.15pt);
\filldraw[navy, fill=white, line width=0.45pt] (3.576,0.783) circle (1.15pt);
\filldraw[navy, fill=white, line width=0.45pt] (3.624,0.89) circle (1.15pt);
\filldraw[navy, fill=white, line width=0.45pt] (3.648,1.035) circle (1.15pt);
\filldraw[navy, fill=white, line width=0.45pt] (3.672,0.838) circle (1.15pt);
\filldraw[navy, fill=white, line width=0.45pt] (3.696,0.838) circle (1.15pt);
\filldraw[navy, fill=white, line width=0.45pt] (3.744,0.725) circle (1.15pt);
\filldraw[navy, fill=white, line width=0.45pt] (3.768,0.725) circle (1.15pt);
\filldraw[navy, fill=white, line width=0.45pt] (3.792,0.664) circle (1.15pt);
\filldraw[navy, fill=white, line width=0.45pt] (3.816,0.838) circle (1.15pt);
\filldraw[navy, fill=white, line width=0.45pt] (3.864,0.989) circle (1.15pt);
\filldraw[navy, fill=white, line width=0.45pt] (3.888,1.035) circle (1.15pt);
\filldraw[navy, fill=white, line width=0.45pt] (3.912,1.348) circle (1.15pt);
\filldraw[navy, fill=white, line width=0.45pt] (3.936,3.193) circle (1.15pt);
\filldraw[navy, fill=white, line width=0.45pt] (3.984,1.241) circle (1.15pt);
\filldraw[navy, fill=white, line width=0.45pt] (4.008,0.838) circle (1.15pt);
\filldraw[navy, fill=white, line width=0.45pt] (4.032,0.725) circle (1.15pt);
\filldraw[navy, fill=white, line width=0.45pt] (4.056,0.725) circle (1.15pt);
\filldraw[navy, fill=white, line width=0.45pt] (4.104,0.783) circle (1.15pt);
\filldraw[navy, fill=white, line width=0.45pt] (4.128,1.035) circle (1.15pt);
\filldraw[navy, fill=white, line width=0.45pt] (4.152,0.599) circle (1.15pt);
\filldraw[navy, fill=white, line width=0.45pt] (4.176,1.163) circle (1.15pt);
\filldraw[navy, fill=white, line width=0.45pt] (4.224,0.89) circle (1.15pt);
\filldraw[navy, fill=white, line width=0.45pt] (4.248,0.725) circle (1.15pt);
\filldraw[navy, fill=white, line width=0.45pt] (4.272,0.989) circle (1.15pt);
\filldraw[navy, fill=white, line width=0.45pt] (4.296,2.542) circle (1.15pt);
\filldraw[navy, fill=white, line width=0.45pt] (4.344,0.664) circle (1.15pt);
\filldraw[navy, fill=white, line width=0.45pt] (4.368,0.838) circle (1.15pt);
\filldraw[navy, fill=white, line width=0.45pt] (4.416,0.89) circle (1.15pt);
\filldraw[navy, fill=white, line width=0.45pt] (4.464,0.531) circle (1.15pt);
\filldraw[navy, fill=white, line width=0.45pt] (4.488,0.664) circle (1.15pt);
\filldraw[navy, fill=white, line width=0.45pt] (4.512,0.664) circle (1.15pt);
\filldraw[navy, fill=white, line width=0.45pt] (4.536,0.599) circle (1.15pt);
\filldraw[navy, fill=white, line width=0.45pt] (4.584,0.664) circle (1.15pt);
\filldraw[navy, fill=white, line width=0.45pt] (4.608,0.89) circle (1.15pt);
\filldraw[navy, fill=white, line width=0.45pt] (4.632,0.89) circle (1.15pt);
\filldraw[navy, fill=white, line width=0.45pt] (4.656,1.278) circle (1.15pt);
\filldraw[navy, fill=white, line width=0.45pt] (4.704,1.771) circle (1.15pt);
\filldraw[navy, fill=white, line width=0.45pt] (4.752,1.723) circle (1.15pt);
\filldraw[navy, fill=white, line width=0.45pt] (4.776,1.313) circle (1.15pt);
\filldraw[navy, fill=white, line width=0.45pt] (4.824,1.348) circle (1.15pt);
\filldraw[navy, fill=white, line width=0.45pt] (4.848,1.035) circle (1.15pt);
\filldraw[navy, fill=white, line width=0.45pt] (4.872,0.664) circle (1.15pt);
\filldraw[navy, fill=white, line width=0.45pt] (4.896,0.725) circle (1.15pt);
\filldraw[navy, fill=white, line width=0.45pt] (4.944,0.531) circle (1.15pt);
\filldraw[navy, fill=white, line width=0.45pt] (4.968,0.838) circle (1.15pt);
\filldraw[navy, fill=white, line width=0.45pt] (4.992,0.531) circle (1.15pt);
\filldraw[navy, fill=white, line width=0.45pt] (5.016,0.989) circle (1.15pt);
\filldraw[navy, fill=white, line width=0.45pt] (5.064,2.29) circle (1.15pt);
\filldraw[navy, fill=white, line width=0.45pt] (5.088,0.599) circle (1.15pt);
\filldraw[navy, fill=white, line width=0.45pt] (5.112,1.382) circle (1.15pt);
\filldraw[navy, fill=white, line width=0.45pt] (5.136,0.94) circle (1.15pt);
\filldraw[navy, fill=white, line width=0.45pt] (5.184,0.725) circle (1.15pt);
\filldraw[navy, fill=white, line width=0.45pt] (5.208,0.783) circle (1.15pt);
\filldraw[navy, fill=white, line width=0.45pt] (5.232,1.079) circle (1.15pt);
\filldraw[navy, fill=white, line width=0.45pt] (5.256,2.748) circle (1.15pt);
\filldraw[navy, fill=white, line width=0.45pt] (5.304,0.725) circle (1.15pt);
\filldraw[navy, fill=white, line width=0.45pt] (5.328,0.664) circle (1.15pt);
\filldraw[navy, fill=white, line width=0.45pt] (5.352,1.122) circle (1.15pt);
\filldraw[navy, fill=white, line width=0.45pt] (5.376,0.989) circle (1.15pt);
\filldraw[navy, fill=white, line width=0.45pt] (5.424,0.599) circle (1.15pt);
\filldraw[navy, fill=white, line width=0.45pt] (5.448,1.163) circle (1.15pt);
\filldraw[navy, fill=white, line width=0.45pt] (5.472,0.94) circle (1.15pt);
\filldraw[navy, fill=white, line width=0.45pt] (5.496,0.725) circle (1.15pt);
\filldraw[navy, fill=white, line width=0.45pt] (5.544,0.599) circle (1.15pt);
\filldraw[navy, fill=white, line width=0.45pt] (5.568,1.035) circle (1.15pt);
\filldraw[navy, fill=white, line width=0.45pt] (5.592,0.838) circle (1.15pt);
\filldraw[navy, fill=white, line width=0.45pt] (5.616,0.664) circle (1.15pt);
\filldraw[navy, fill=white, line width=0.45pt] (5.664,0.89) circle (1.15pt);
\filldraw[navy, fill=white, line width=0.45pt] (5.688,0.783) circle (1.15pt);
\filldraw[navy, fill=white, line width=0.45pt] (5.712,0.725) circle (1.15pt);
\filldraw[navy, fill=white, line width=0.45pt] (5.736,0.89) circle (1.15pt);
\filldraw[navy, fill=white, line width=0.45pt] (5.784,1.079) circle (1.15pt);
\filldraw[navy, fill=white, line width=0.45pt] (5.808,1.035) circle (1.15pt);
\filldraw[navy, fill=white, line width=0.45pt] (5.832,1.163) circle (1.15pt);
\filldraw[navy, fill=white, line width=0.45pt] (5.856,1.537) circle (1.15pt);
\filldraw[navy, fill=white, line width=0.45pt] (5.904,3.63) circle (1.15pt);
\filldraw[navy, fill=white, line width=0.45pt] (5.928,2.443) circle (1.15pt);
\filldraw[navy, fill=white, line width=0.45pt] (5.952,1.62) circle (1.15pt);
\filldraw[navy, fill=white, line width=0.45pt] (5.976,1.241) circle (1.15pt);
\filldraw[navy, fill=white, line width=0.45pt] (6.024,0.94) circle (1.15pt);
\filldraw[navy, fill=white, line width=0.45pt] (6.048,1.415) circle (1.15pt);
\filldraw[navy, fill=white, line width=0.45pt] (6.072,0.725) circle (1.15pt);
\filldraw[navy, fill=white, line width=0.45pt] (6.096,0.838) circle (1.15pt);
\filldraw[navy, fill=white, line width=0.45pt] (6.144,0.599) circle (1.15pt);
\filldraw[navy, fill=white, line width=0.45pt] (6.168,1.817) circle (1.15pt);
\filldraw[navy, fill=white, line width=0.45pt] (6.192,1.313) circle (1.15pt);
\filldraw[navy, fill=white, line width=0.45pt] (6.216,0.94) circle (1.15pt);
\filldraw[navy, fill=white, line width=0.45pt] (6.264,0.531) circle (1.15pt);
\filldraw[navy, fill=white, line width=0.45pt] (6.288,0.664) circle (1.15pt);
\filldraw[navy, fill=white, line width=0.45pt] (6.312,0.89) circle (1.15pt);
\filldraw[navy, fill=white, line width=0.45pt] (6.336,0.531) circle (1.15pt);
\filldraw[navy, fill=white, line width=0.45pt] (6.384,0.458) circle (1.15pt);
\filldraw[navy, fill=white, line width=0.45pt] (6.408,0.531) circle (1.15pt);
\filldraw[navy, fill=white, line width=0.45pt] (6.432,0.89) circle (1.15pt);
\filldraw[navy, fill=white, line width=0.45pt] (6.456,1.241) circle (1.15pt);
\filldraw[navy, fill=white, line width=0.45pt] (6.504,0.599) circle (1.15pt);
\filldraw[navy, fill=white, line width=0.45pt] (6.528,0.458) circle (1.15pt);
\filldraw[navy, fill=white, line width=0.45pt] (6.552,0.531) circle (1.15pt);
\filldraw[navy, fill=white, line width=0.45pt] (6.576,0.664) circle (1.15pt);
\filldraw[navy, fill=white, line width=0.45pt] (6.624,0.664) circle (1.15pt);
\filldraw[navy, fill=white, line width=0.45pt] (6.648,1.507) circle (1.15pt);
\filldraw[navy, fill=white, line width=0.45pt] (6.672,0.599) circle (1.15pt);
\filldraw[navy, fill=white, line width=0.45pt] (6.696,0.664) circle (1.15pt);
\filldraw[navy, fill=white, line width=0.45pt] (6.744,1.883) circle (1.15pt);
\filldraw[navy, fill=white, line width=0.45pt] (6.768,1.62) circle (1.15pt);
\filldraw[navy, fill=white, line width=0.45pt] (6.792,0.89) circle (1.15pt);
\filldraw[navy, fill=white, line width=0.45pt] (6.816,0.664) circle (1.15pt);
\filldraw[navy, fill=white, line width=0.45pt] (6.864,1.537) circle (1.15pt);
\filldraw[navy, fill=white, line width=0.45pt] (6.888,0.599) circle (1.15pt);
\filldraw[navy, fill=white, line width=0.45pt] (6.912,0.458) circle (1.15pt);
\filldraw[navy, fill=white, line width=0.45pt] (6.936,0.531) circle (1.15pt);
\filldraw[navy, fill=white, line width=0.45pt] (6.984,1.241) circle (1.15pt);
\filldraw[navy, fill=white, line width=0.45pt] (7.008,0.725) circle (1.15pt);
\filldraw[navy, fill=white, line width=0.45pt] (7.032,0.664) circle (1.15pt);
\filldraw[navy, fill=white, line width=0.45pt] (7.056,0.725) circle (1.15pt);
\filldraw[navy, fill=white, line width=0.45pt] (7.104,1.163) circle (1.15pt);
\filldraw[navy, fill=white, line width=0.45pt] (7.128,0.664) circle (1.15pt);
\filldraw[navy, fill=white, line width=0.45pt] (7.152,0.531) circle (1.15pt);
\filldraw[navy, fill=white, line width=0.45pt] (7.176,1.079) circle (1.15pt);
\filldraw[navy, fill=white, line width=0.45pt] (7.224,1.415) circle (1.15pt);
\filldraw[navy, fill=white, line width=0.45pt] (7.248,0.599) circle (1.15pt);
\filldraw[navy, fill=white, line width=0.45pt] (7.272,2.362) circle (1.15pt);
\filldraw[navy, fill=white, line width=0.45pt] (7.296,0.989) circle (1.15pt);
\filldraw[navy, fill=white, line width=0.45pt] (7.344,0.599) circle (1.15pt);
\filldraw[navy, fill=white, line width=0.45pt] (7.368,0.94) circle (1.15pt);
\filldraw[navy, fill=white, line width=0.45pt] (7.392,1.241) circle (1.15pt);
\filldraw[navy, fill=white, line width=0.45pt] (7.416,1.202) circle (1.15pt);
\filldraw[navy, fill=white, line width=0.45pt] (7.464,2.004) circle (1.15pt);
\filldraw[navy, fill=white, line width=0.45pt] (7.488,0.599) circle (1.15pt);
\filldraw[navy, fill=white, line width=0.45pt] (7.512,0.89) circle (1.15pt);
\filldraw[navy, fill=white, line width=0.45pt] (7.536,0.725) circle (1.15pt);
\filldraw[navy, fill=white, line width=0.45pt] (7.584,0.783) circle (1.15pt);
\filldraw[navy, fill=white, line width=0.45pt] (7.608,1.035) circle (1.15pt);
\filldraw[navy, fill=white, line width=0.45pt] (7.632,0.94) circle (1.15pt);
\filldraw[navy, fill=white, line width=0.45pt] (7.656,0.89) circle (1.15pt);
\filldraw[navy, fill=white, line width=0.45pt] (7.704,0.989) circle (1.15pt);
\filldraw[navy, fill=white, line width=0.45pt] (7.728,1.477) circle (1.15pt);
\filldraw[navy, fill=white, line width=0.45pt] (7.752,1.241) circle (1.15pt);
\filldraw[navy, fill=white, line width=0.45pt] (7.776,1.537) circle (1.15pt);
\filldraw[navy, fill=white, line width=0.45pt] (7.824,1.883) circle (1.15pt);
\filldraw[navy, fill=white, line width=0.45pt] (7.848,2.443) circle (1.15pt);
\filldraw[navy, fill=white, line width=0.45pt] (7.872,3.948) circle (1.15pt);
\filldraw[navy, fill=white, line width=0.45pt] (7.896,3.179) circle (1.15pt);
\filldraw[navy, fill=white, line width=0.45pt] (7.944,1.747) circle (1.15pt);
\filldraw[navy, fill=white, line width=0.45pt] (7.968,1.507) circle (1.15pt);
\filldraw[navy, fill=white, line width=0.45pt] (7.992,1.278) circle (1.15pt);
\filldraw[navy, fill=white, line width=0.45pt] (8.016,1.163) circle (1.15pt);
\filldraw[navy, fill=white, line width=0.45pt] (8.064,0.94) circle (1.15pt);
\filldraw[navy, fill=white, line width=0.45pt] (8.088,1.537) circle (1.15pt);
\filldraw[navy, fill=white, line width=0.45pt] (8.112,0.989) circle (1.15pt);
\filldraw[navy, fill=white, line width=0.45pt] (8.136,0.838) circle (1.15pt);
\filldraw[navy, fill=white, line width=0.45pt] (8.184,1.747) circle (1.15pt);
\filldraw[navy, fill=white, line width=0.45pt] (8.208,1.035) circle (1.15pt);
\filldraw[navy, fill=white, line width=0.45pt] (8.232,0.838) circle (1.15pt);
\filldraw[navy, fill=white, line width=0.45pt] (8.256,0.783) circle (1.15pt);
\filldraw[navy, fill=white, line width=0.45pt] (8.304,0.783) circle (1.15pt);
\filldraw[navy, fill=white, line width=0.45pt] (8.328,0.94) circle (1.15pt);
\filldraw[navy, fill=white, line width=0.45pt] (8.352,0.838) circle (1.15pt);
\filldraw[navy, fill=white, line width=0.45pt] (8.376,0.531) circle (1.15pt);
\filldraw[navy, fill=white, line width=0.45pt] (8.424,0.725) circle (1.15pt);
\filldraw[navy, fill=white, line width=0.45pt] (8.448,1.278) circle (1.15pt);
\filldraw[navy, fill=white, line width=0.45pt] (8.472,0.599) circle (1.15pt);
\filldraw[navy, fill=white, line width=0.45pt] (8.496,0.989) circle (1.15pt);
\filldraw[navy, fill=white, line width=0.45pt] (8.544,0.599) circle (1.15pt);
\filldraw[navy, fill=white, line width=0.45pt] (8.568,0.599) circle (1.15pt);
\filldraw[navy, fill=white, line width=0.45pt] (8.592,1.84) circle (1.15pt);
\filldraw[navy, fill=white, line width=0.45pt] (8.616,0.783) circle (1.15pt);
\filldraw[navy, fill=white, line width=0.45pt] (8.664,0.664) circle (1.15pt);
\filldraw[navy, fill=white, line width=0.45pt] (8.688,0.664) circle (1.15pt);
\filldraw[navy, fill=white, line width=0.45pt] (8.712,1.348) circle (1.15pt);
\filldraw[navy, fill=white, line width=0.45pt] (8.736,1.202) circle (1.15pt);
\filldraw[navy, fill=white, line width=0.45pt] (8.784,0.599) circle (1.15pt);
\filldraw[navy, fill=white, line width=0.45pt] (8.808,0.838) circle (1.15pt);
\filldraw[navy, fill=white, line width=0.45pt] (8.832,0.838) circle (1.15pt);
\filldraw[navy, fill=white, line width=0.45pt] (8.856,1.747) circle (1.15pt);
\filldraw[navy, fill=white, line width=0.45pt] (8.904,0.783) circle (1.15pt);
\filldraw[navy, fill=white, line width=0.45pt] (8.928,0.664) circle (1.15pt);
\filldraw[navy, fill=white, line width=0.45pt] (8.952,0.458) circle (1.15pt);
\filldraw[navy, fill=white, line width=0.45pt] (8.976,0.458) circle (1.15pt);
\filldraw[navy, fill=white, line width=0.45pt] (9.024,0.531) circle (1.15pt);
\filldraw[navy, fill=white, line width=0.45pt] (9.048,0.531) circle (1.15pt);
\filldraw[navy, fill=white, line width=0.45pt] (9.072,0.599) circle (1.15pt);
\filldraw[navy, fill=white, line width=0.45pt] (9.096,1.673) circle (1.15pt);
\filldraw[navy, fill=white, line width=0.45pt] (9.144,0.725) circle (1.15pt);
\filldraw[navy, fill=white, line width=0.45pt] (9.168,0.599) circle (1.15pt);
\filldraw[navy, fill=white, line width=0.45pt] (9.192,2.023) circle (1.15pt);
\filldraw[navy, fill=white, line width=0.45pt] (9.216,1.035) circle (1.15pt);
\filldraw[navy, fill=white, line width=0.45pt] (9.264,0.94) circle (1.15pt);
\filldraw[navy, fill=white, line width=0.45pt] (9.288,1.647) circle (1.15pt);
\filldraw[navy, fill=white, line width=0.45pt] (9.312,1.507) circle (1.15pt);
\filldraw[navy, fill=white, line width=0.45pt] (9.336,0.664) circle (1.15pt);
\filldraw[navy, fill=white, line width=0.45pt] (9.384,0.989) circle (1.15pt);
\filldraw[navy, fill=white, line width=0.45pt] (9.408,1.313) circle (1.15pt);
\filldraw[navy, fill=white, line width=0.45pt] (9.432,1.723) circle (1.15pt);
\filldraw[navy, fill=white, line width=0.45pt] (9.504,1.163) circle (1.15pt);
\filldraw[navy, fill=white, line width=0.45pt] (9.528,0.89) circle (1.15pt);
\filldraw[navy, fill=white, line width=0.45pt] (9.552,0.783) circle (1.15pt);
\filldraw[navy, fill=white, line width=0.45pt] (9.576,0.989) circle (1.15pt);
\filldraw[navy, fill=white, line width=0.45pt] (9.624,0.599) circle (1.15pt);
\filldraw[navy, fill=white, line width=0.45pt] (9.648,0.783) circle (1.15pt);
\filldraw[navy, fill=white, line width=0.45pt] (9.672,1.84) circle (1.15pt);
\filldraw[navy, fill=white, line width=0.45pt] (9.696,0.599) circle (1.15pt);
\filldraw[navy, fill=white, line width=0.45pt] (9.744,0.664) circle (1.15pt);
\filldraw[navy, fill=white, line width=0.45pt] (9.768,0.458) circle (1.15pt);
\filldraw[navy, fill=white, line width=0.45pt] (9.816,0.531) circle (1.15pt);
\filldraw[navy, fill=white, line width=0.45pt] (9.864,0.94) circle (1.15pt);
\filldraw[navy, fill=white, line width=0.45pt] (9.888,0.531) circle (1.15pt);
\filldraw[navy, fill=white, line width=0.45pt] (9.912,0.989) circle (1.15pt);
\filldraw[navy, fill=white, line width=0.45pt] (9.936,0.94) circle (1.15pt);
\filldraw[navy, fill=white, line width=0.45pt] (9.984,0.531) circle (1.15pt);
\filldraw[navy, fill=white, line width=0.45pt] (10.008,0.725) circle (1.15pt);
\filldraw[navy, fill=white, line width=0.45pt] (10.032,0.725) circle (1.15pt);
\filldraw[navy, fill=white, line width=0.45pt] (10.056,0.725) circle (1.15pt);
\filldraw[navy, fill=white, line width=0.45pt] (10.104,1.163) circle (1.15pt);
\filldraw[navy, fill=white, line width=0.45pt] (10.128,3.012) circle (1.15pt);
\filldraw[navy, fill=white, line width=0.45pt] (10.152,1.348) circle (1.15pt);
\filldraw[navy, fill=white, line width=0.45pt] (10.176,0.89) circle (1.15pt);
\filldraw[navy, fill=white, line width=0.45pt] (10.224,0.94) circle (1.15pt);
\filldraw[navy, fill=white, line width=0.45pt] (10.248,0.664) circle (1.15pt);
\filldraw[navy, fill=white, line width=0.45pt] (10.272,0.725) circle (1.15pt);
\filldraw[navy, fill=white, line width=0.45pt] (10.296,0.531) circle (1.15pt);
\filldraw[navy, fill=white, line width=0.45pt] (10.344,2.148) circle (1.15pt);
\filldraw[navy, fill=white, line width=0.45pt] (10.368,1.079) circle (1.15pt);
\filldraw[navy, fill=white, line width=0.45pt] (10.392,0.531) circle (1.15pt);
\filldraw[navy, fill=white, line width=0.45pt] (10.416,0.531) circle (1.15pt);
\filldraw[navy, fill=white, line width=0.45pt] (10.464,1.122) circle (1.15pt);
\filldraw[navy, fill=white, line width=0.45pt] (10.488,1.079) circle (1.15pt);
\filldraw[navy, fill=white, line width=0.45pt] (10.512,2.042) circle (1.15pt);
\filldraw[navy, fill=white, line width=0.45pt] (10.536,0.783) circle (1.15pt);
\filldraw[navy, fill=white, line width=0.45pt] (10.584,0.599) circle (1.15pt);
\filldraw[navy, fill=white, line width=0.45pt] (10.608,0.89) circle (1.15pt);
\filldraw[navy, fill=white, line width=0.45pt] (10.632,0.783) circle (1.15pt);
\filldraw[navy, fill=white, line width=0.45pt] (10.656,0.531) circle (1.15pt);
\filldraw[navy, fill=white, line width=0.45pt] (10.704,0.783) circle (1.15pt);
\filldraw[navy, fill=white, line width=0.45pt] (10.728,0.838) circle (1.15pt);
\filldraw[navy, fill=white, line width=0.45pt] (10.752,1.507) circle (1.15pt);
\filldraw[navy, fill=white, line width=0.45pt] (10.776,0.725) circle (1.15pt);
\filldraw[navy, fill=white, line width=0.45pt] (10.824,0.599) circle (1.15pt);
\filldraw[navy, fill=white, line width=0.45pt] (10.848,0.531) circle (1.15pt);
\filldraw[navy, fill=white, line width=0.45pt] (10.872,0.664) circle (1.15pt);
\filldraw[navy, fill=white, line width=0.45pt] (10.896,0.783) circle (1.15pt);
\filldraw[navy, fill=white, line width=0.45pt] (10.944,1.202) circle (1.15pt);
\filldraw[navy, fill=white, line width=0.45pt] (10.968,0.664) circle (1.15pt);
\filldraw[navy, fill=white, line width=0.45pt] (10.992,0.725) circle (1.15pt);
\filldraw[navy, fill=white, line width=0.45pt] (11.016,0.725) circle (1.15pt);
\filldraw[navy, fill=white, line width=0.45pt] (11.064,0.783) circle (1.15pt);
\filldraw[navy, fill=white, line width=0.45pt] (11.088,0.838) circle (1.15pt);
\filldraw[navy, fill=white, line width=0.45pt] (11.112,0.783) circle (1.15pt);
\filldraw[navy, fill=white, line width=0.45pt] (11.136,0.989) circle (1.15pt);
\filldraw[navy, fill=white, line width=0.45pt] (11.184,0.783) circle (1.15pt);
\filldraw[navy, fill=white, line width=0.45pt] (11.208,0.838) circle (1.15pt);
\filldraw[navy, fill=white, line width=0.45pt] (11.232,0.783) circle (1.15pt);
\filldraw[navy, fill=white, line width=0.45pt] (11.256,1.965) circle (1.15pt);
\filldraw[navy, fill=white, line width=0.45pt] (11.304,1.382) circle (1.15pt);
\filldraw[navy, fill=white, line width=0.45pt] (11.328,1.382) circle (1.15pt);
\filldraw[navy, fill=white, line width=0.45pt] (11.352,1.035) circle (1.15pt);
\filldraw[navy, fill=white, line width=0.45pt] (11.376,0.94) circle (1.15pt);
\filldraw[navy, fill=white, line width=0.45pt] (11.424,0.989) circle (1.15pt);
\filldraw[navy, fill=white, line width=0.45pt] (11.448,0.989) circle (1.15pt);
\filldraw[navy, fill=white, line width=0.45pt] (11.472,1.348) circle (1.15pt);
\filldraw[navy, fill=white, line width=0.45pt] (11.496,1.163) circle (1.15pt);
\filldraw[navy, fill=white, line width=0.45pt] (11.544,1.382) circle (1.15pt);
\filldraw[navy, fill=white, line width=0.45pt] (11.568,1.673) circle (1.15pt);
\filldraw[navy, fill=white, line width=0.45pt] (11.592,1.348) circle (1.15pt);
\filldraw[navy, fill=white, line width=0.45pt] (11.616,1.537) circle (1.15pt);
\filldraw[navy, fill=white, line width=0.45pt] (11.664,1.673) circle (1.15pt);
\filldraw[navy, fill=white, line width=0.45pt] (11.688,1.84) circle (1.15pt);
\filldraw[navy, fill=white, line width=0.45pt] (11.712,2.042) circle (1.15pt);
\filldraw[navy, fill=white, line width=0.45pt] (11.736,2.305) circle (1.15pt);
\filldraw[navy, fill=white, line width=0.45pt] (11.784,3.179) circle (1.15pt);
\filldraw[navy, fill=white, line width=0.45pt] (11.808,4.401) circle (1.15pt);
\fill[ochre] (4.392,0.599) circle (1.9pt);
\node[ochre, below] at (4.392,0.368) {$183$};
\fill[wine] (9.792,0.458) circle (1.9pt);
\node[wine, below] at (9.792,0.368) {$408$};
\draw[wine, line width=0.6pt, fill=none] (9.792,0.458) circle (3.2pt);
\draw[ink, line width=0.5pt] (-0.2,0) -- (12,0);
\draw[ink, line width=0.5pt] (-0.2,0) -- (-0.2,5.2);
\draw[ink, line width=0.5pt] (0,0) -- (0,-0.08);
\node[ink, below] at (0,-0.08) {$0$};
\draw[ink, line width=0.5pt] (2.4,0) -- (2.4,-0.08);
\node[ink, below] at (2.4,-0.08) {$100$};
\draw[ink, line width=0.5pt] (4.8,0) -- (4.8,-0.08);
\node[ink, below] at (4.8,-0.08) {$200$};
\draw[ink, line width=0.5pt] (7.2,0) -- (7.2,-0.08);
\node[ink, below] at (7.2,-0.08) {$300$};
\draw[ink, line width=0.5pt] (9.6,0) -- (9.6,-0.08);
\node[ink, below] at (9.6,-0.08) {$400$};
\draw[ink, line width=0.5pt] (12,0) -- (12,-0.08);
\node[ink, below] at (12,-0.08) {$500$};
\draw[ink, line width=0.5pt] (-0.2,0) -- (-0.28,0);
\node[ink, left] at (-0.28,0) {$10$};
\draw[ink, line width=0.5pt] (-0.2,0.458) -- (-0.28,0.458);
\node[ink, left] at (-0.28,0.458) {$15$};
\draw[ink, line width=0.5pt] (-0.2,0.783) -- (-0.28,0.783);
\node[ink, left] at (-0.28,0.783) {$20$};
\draw[ink, line width=0.5pt] (-0.2,1.817) -- (-0.28,1.817);
\node[ink, left] at (-0.28,1.817) {$50$};
\draw[ink, line width=0.5pt] (-0.2,2.6) -- (-0.28,2.6);
\node[ink, left] at (-0.28,2.6) {$100$};
\draw[ink, line width=0.5pt] (-0.2,3.383) -- (-0.28,3.383);
\node[ink, left] at (-0.28,3.383) {$200$};
\draw[ink, line width=0.5pt] (-0.2,4.417) -- (-0.28,4.417);
\node[ink, left] at (-0.28,4.417) {$500$};
\draw[ink, line width=0.5pt] (-0.2,5.2) -- (-0.28,5.2);
\node[ink, left] at (-0.28,5.2) {$1000$};
\node[ink] at (6,-0.75) {$n$};
\node[ink, rotate=90] at (-1.2,2.6) {$\ell(n,985)$};
\end{tikzpicture}

\caption{The Euclidean cost $\ell(n,m)$ for $m=m_{4/5}=985$ and the $392$
integers $n$ prime to $m$ with $1\leqslant n<m/2$ (hollow circles), on a
logarithmic scale; the directive of the central word of $\gamma_{n/m}$ has
$\ell(n,m)-1$ letters. The two square roots of $-1$ modulo $m$ in $(1,m/2)$ are
filled: the root $408$ of the occurrence $4/5$, of cost $15=2\cdot7+1$, and
$183$, of cost $17=2\cdot8+1$, which correspond by \cref{lem:square-roots} to the
representations $12^2+29^2$ and $16^2+27^2$. Six other integers, $266$,
$272$, $288$, $373$, $374$ and $407$, share the least cost $15$ with the root
(dashed line), but only the root, ringed, satisfies \eqref{eq:block-form}.}
\label{fig:directives}
\end{figure}

\Cref{conj:directives} is a rule of the kind discussed in
\cref{rem:pair-determines}: it designates the root from the label alone, and
with it the occurrence pair, which \cref{lem:square-roots} reads off the
first half of the continued fraction of $m/u_{P/Q}$. Like
\cref{conj:euclidean-minimality}, it can hold only because of the Markov
structure of the label, as the integers $31$ and $1130$ of
\cref{lem:directive-facts} show, and it
reaches the uniqueness conjecture by a different road: through the
palindromic closure and the Thue--Morse morphism rather than through the
comparison of costs.
